\documentclass[11pt,a4paper]{article}
\usepackage[T1]{fontenc}
\usepackage[utf8]{inputenc}
\usepackage{lmodern,amsmath,amssymb,amsthm,mathtools,microtype,mathrsfs}
\usepackage[margin=25mm,top=23mm,bottom=25mm]{geometry}
\usepackage{booktabs,array,longtable,enumitem,xcolor,fancyhdr}
\usepackage[colorlinks=true,linkcolor=blue!45!black,urlcolor=blue!45!black,citecolor=blue!45!black]{hyperref}
\hypersetup{pdftitle={Reading the Alpoge--Claude RH and Dirichlet-GRH proof: working version 0.49},pdfauthor={ChatGPT},pdfsubject={Partial human-readable reconstruction with explicit arithmetic inputs},pdfkeywords={Riemann hypothesis, Lean, theta lifts, prime sums, Arakelov slopes}}
\setlength{\parindent}{0pt}
\setlength{\parskip}{5pt plus 1pt}
\setlist{topsep=3pt,itemsep=3pt,leftmargin=*}
\setlength{\headheight}{14pt}
\pagestyle{fancy}\fancyhf{}
\fancyhead[L]{\small RH and Dirichlet GRH: a reconstruction}
\fancyhead[R]{\small Working version 0.49}
\fancyfoot[C]{\thepage}
\renewcommand{\headrulewidth}{0.3pt}
\numberwithin{equation}{section}
\newtheorem{theorem}{Theorem}[section]
\newtheorem{proposition}[theorem]{Proposition}
\newtheorem{lemma}[theorem]{Lemma}
\newtheorem{corollary}[theorem]{Corollary}
\theoremstyle{remark}\newtheorem{remark}[theorem]{Remark}
\newcommand{\Rea}{\operatorname{Re}}
\newcommand{\HS}{\mathrm{HS}}
\newcommand{\tr}{\operatorname{tr}}
\newcommand{\codim}{\operatorname{codim}}
\newcommand{\covol}{\operatorname{covol}}
\newcommand{\OK}{\mathcal O_K}
\newcommand{\src}[2]{\href{https://github.com/CocoSgt/riemann-hypothesis-lean/blob/7cbf7a0c32fca1badf05931879a37c3f8a2b03b4/src/#1}{\texttt{#2}}}
\newcommand{\status}[1]{\textbf{Status:} #1}
\begin{document}
\begin{center}
{\LARGE\bfseries RH and Dirichlet GRH: new reconstruction chapters}\\[7pt]
{\large Version 0.49 --- 23 November 2026}\\[5pt]
{\large Uniform local gain and holomorphic families}
\end{center}
\textbf{Scope.} This edition contains the current audit and the new
chapters. It remains a partial human reconstruction. Earlier named
claims are in the cumulative manuscript and preserved component files.
The underlying formal proof is accepted as correct; the exact remaining
human reconstruction boundaries are stated below.
\setcounter{tocdepth}{1}
\tableofcontents
\clearpage
\begingroup\emergencystretch=3em
\section{Version 0.49: uniform local gain and holomorphic families}
\label{sec:dependency-audit-v49}

\subsection{Current result and the completed selected estimate}

The pinned formal source is accepted as fully checked and correct for
reconstruction. The task is to write its mathematical argument,
allowing simpler proofs of the same original statements. Version 0.49
adds 74 numbered class-B claims to the 2940 preserved predecessor
claims, for 3014 in all. The claim graph records these human proofs;
it is not a mechanically extracted formal dependency graph.

The selected \texttt{GLOBAL\_\allowbreak LOCAL\_\allowbreak GAIN} contract is now complete.
The existing built-class membership already contains one joint
geometric witness with the given comparison and kernel data. This
supplies the original \texttt{Group\allowbreak Compatibility\allowbreak Family} after one
family threshold. The unchanged D3 chosen-data condition supplies
the required tops and admission. The completed quantitative argument
therefore gives all five original \texttt{Link\allowbreak G\_\allowbreak at} clauses with
\[
 C_G=4C_{\mathrm{count}}+1,
 \qquad\delta(x)=x^{-\epsilon/2}\longrightarrow0.
\]
These are chosen before every subsequent actual index. The stronger
historical \texttt{WPt}/\texttt{TSel} construction is bypassed for
this estimate. It is not declared solved. In particular, compatibility
of a supplied member does not assert existence of the entire common
design family.

The selected sufficient route retains 121 nodes and 161 edges.
Its whole-obligation register decreases from 25 to 24:
17 class C and 7 class B. There are 17 shared RH/Dirichlet-GRH
obligations, 7 additional Dirichlet-GRH obligations, and no
RH-only obligations in this route. One whole obligation can contain
many substantial producers; these counts are an inventory.

\subsection{The four new chapters}

\paragraph{UCE1--UCE19: ordinary points and coefficient extension.}
An inertia separator and specialization of ordered flags give the
literal ordinary local condition over every allowed complete local
domain, including characteristic~\(\ell\). Integral spectral
projectors keep the prescribed character order. The five
\texttt{Ord\allowbreak Pt} fields use the actual local embedding and the full
away-inertia condition.

Actual entry deviations give the finite weighted presentation.
Conjugation by the formal torus \(T(1+T)\) proves homogeneous
projection stability integrally. Finite monomial generators and a
characteristic-prime DVR curve argument prove the whole original
\texttt{Fin\allowbreak Crit}, with no reducedness or coefficient-flatness
restriction. Finite-index characteristic polynomials force constant
diagonals and one triangular orientation. The argument includes
the permitted abstract characteristic-three case. For the final
arithmetic datum, the original \texttt{GD.\allowbreak gen} guard
\(\ell-1\nmid 2e\) itself excludes \(\ell=3\); no new
prime exclusion is introduced.

Finite weight zero controls arbitrary entry-algebra maps through
finite homogeneous modules. Their continuity is proved at the
required stage. Torus rescaling gives an actual finite-coefficient
ordinary representative, and the irreducibility argument works on
every original open subgroup. Both automorphic conclusions transport
back at the same \(\iota\), including every supplied full
outer-corner extension. Hence
\[
 \mathrm{GdEnd}\quad\Longrightarrow\quad
 \mathrm{AutomorphyUnramifiedCoefficients}
\]
is reconstructed in its complete original scope. The same marked
\(g_1\) supplies \(\mathrm{FinA}\), \(\mathrm{AutB}\) and
\(\mathrm{LemW}\). Its upstream \(\mathrm{FinEnd}\),
\(\mathrm{AutEnd3}\) and \(\mathrm{AutEnd2}\) have not been
constructed here. The common CM frame, both seeds and full
\texttt{Pfrom\allowbreak Aut}/\texttt{Mod\allowbreak UAlt} remain explicit.
Thus \texttt{AUT\_\allowbreak UNRAM\_\allowbreak COEF} stays open.

\paragraph{BWP1--BWP24: the original normalized wave packet.}
An explicit polynomial chart of the original unipotent group has
Haar Jacobian one and real dimension \(2r^2+2rk+r\).
Its inverse-coordinate bound gives integrability at the exact source
exponent \(N=4(2r+k+1)^2\sqrt{2(r+1)}\).
Harish--Chandra admissibility gives a finite-dimensional full compact
isotypic space. A conjugation-averaged approximate identity and
Cayley--Hamilton then produce an exact compact reproducing function
in any prescribed identity neighborhood.

The printed cubic gap dominates every coefficient-decay bound.
This proves the whole \texttt{Coef\allowbreak Weighted\allowbreak Loc} for each supplied
discrete series and lowest unit vector. The wave packet is absolutely
convergent, continuous and equivariant, with the exact weighted
quotient integrability and square norm. Both compact reproduction
and reproduction by the original formal-degree kernel are proved.
The retained \texttt{Bessel\allowbreak NVGap} supplies its nonzero norm;
therefore \texttt{Wloc\allowbreak Exists} is complete on its original antecedents.

Transport uses separate conjugacies on the ambient and Bessel groups,
and an equivariant map on their homogeneous spaces. This preserves
the factor \(4\pi/g_D\) in the character and all supplied Haar
and quotient measures, even when the actual frame is not unitary.
The normalized tensor product is chosen before the later measured
block, and has the required decay for every such block. This gives
the original \texttt{WFrom\allowbreak Loc}, and \texttt{WExists} conditional
on \texttt{Bessel\allowbreak NVGap}. The full analytic setting also uses the
inherited original \texttt{Forms0} input. Bessel nonvanishing,
multiplicity one, the orbital/fiber ties and Siegel--Weil comparisons
remain separate.

\paragraph{HUF1--HUF21: the whole finite holomorphic family.}
Starting with the supplied moduli scheme, the argument extracts the
actual exponential kernel, perfect integral CM datum and all-level
Tate/genus comparison from each point. Relative exponentials and
continued lattice bases give local lifts with the prescribed initial
period and every positive Weil level on one neighborhood.

Period-kernel coordinates have the exact dimension \(fpq\).
Marked uniqueness and the local inverse theorem give holomorphy even
for arbitrary pointwise choices of points and exponentials. A finite
affine cover and finite semialgebraic components give one finite
positive family. The moduli space need not be compact. The family
and its total maps are selected before all period and base choices.

The primitive semicharacter gives fixed half characteristics and the
exact theta multiplier. Proper analytic base change applied to the
ratio line, followed by its zero rigidification, gives a global
normalized theta-line isomorphism. Projective GAGA supplies the
algebraic section separately on every fiber. All five
\texttt{Is\allowbreak Theta\allowbreak Frame} clauses, every \texttt{Family\allowbreak Unif} field,
both markings, the negative Weil sign and zero-dimensional cases
are checked. Thus the whole original \texttt{N1c} is complete
on its supplied \texttt{Moduli} input. Principal framed charts,
\texttt{Is\allowbreak Unf}, reduced-test \texttt{From\allowbreak Test} and the later theta
spanning statements remain adjoining producers.

\paragraph{WRS1--WRS10: full uniform local gain.}
Every local/Taylor kernel field survives division by arbitrary
original submodules. A map on the dependent comparison-kernel tuple
keeps the actual object and kernel together. Extracting the joint
witness already contained in built-class membership gives common
compatibility for the original family. Existing selected tops give
the other original input to the local-gain implication.

One coefficient bound is chosen before the family, with
\(\mathrm{pre}=\mathrm{t}=0\), \(\mathrm{mid}=4129\), and
\(\mathrm{last}=\max(5600002,C_{\mathrm{orb}}(\epsilon))\).
The original closedness cutoff is \(1/104\). The resulting estimate
retains the actual lattice divisibility and lengths, zero conditions,
orbit bounds, signed gains, served-orbit coverage and mean bound.
The universal installer from the original \texttt{TSel}/\texttt{Duty}
premises is also proved. The direct selected-class proof adds no
new selector premise. This is the one whole selected closure of
the present edition.

\subsection{Foundations, preservation and independent review}

All 74 numbered claims are class-B deductions. Named classical
foundations remain separate inputs, including the inherited
Harish--Chandra admissibility, Appell--Humbert and projective GAGA,
the injective holomorphic map theorem in equal dimension, finite
semialgebraic components, and Grauert's proper analytic base-change
theorem with its corrigendum. Their hypotheses and exact uses are
recorded. For the ratio line, properness, a reduced base, flatness
and constant \(h^0=1\) are checked. The rigidification supplies
global triviality; base change alone does not.

Every new claim has independent mathematical and source review by
an agent other than its author. Reports bind the exact proof file
and reviewed blocks. All cited source bytes are compared with the
pinned git blobs. Every predecessor claim and historical proof body
is preserved, and mutable predecessor metadata is frozen in
\texttt{history/\allowbreak version\_\allowbreak 0\_\allowbreak 48}.

Draft TeX compilation checks syntax, references and missing characters
in both readers. It is not PDF visual review or formal replay.
This remains a partial human reconstruction of RH and unrestricted
Dirichlet GRH. The following version-0.49 register is the current
scope authority. Historical conditional statements retain their
printed hypotheses even where a later chapter supplies them.

\subsection{The complete current selected register}

\paragraph{1. GLOBAL\_\allowbreak DENOMINATOR — C — Actual LinkN local table and binding producer}

Closed(Des $\mathcal C$ 3) LinkN: the actual etaSum is bounded with exact dZ0 and same-host vvTerm, subject only to the literal ShaFinOdd and summability guards.

\emph{Shared objects and choice order.} one FamP; one class $\mathcal A$ and its target-compatible subclass; one coefficient field per family datum; one x and i at each application; original selected rows, words and levels; thresholds chosen before the common family; actual D.nd; actual etaJ and Aj; sha\_\allowbreak eq

\emph{Source.} \href{https://github.com/CocoSgt/riemann-hypothesis-lean/blob/7cbf7a0c32fca1badf05931879a37c3f8a2b03b4/src/Formalism/SharpRestriction/GlobalTheoremInputs.lean#L45-L48}{LinkN}.

\emph{Source.} \href{https://github.com/CocoSgt/riemann-hypothesis-lean/blob/7cbf7a0c32fca1badf05931879a37c3f8a2b03b4/src/Formalism/SharpRestriction/GlobalTheoremInputs.lean#L40-L44}{LinkN\_\allowbreak at}.

\emph{Source.} \href{https://github.com/CocoSgt/riemann-hypothesis-lean/blob/7cbf7a0c32fca1badf05931879a37c3f8a2b03b4/src/Development/AutomorphicArithmetic/CharacterBoxes/DivinNOfHypV.lean#L117-L134}{divinN\_\allowbreak of\_\allowbreak hypV}.



\paragraph{2. GLOBAL\_\allowbreak RETAINED\_\allowbreak DEFECT — C — Original retained-row defect bound LinkIRR}

Rows.DelpNonneg and delp$\le$exp(-x\textasciicircum{}(ez/7)) uniformly for every late actual row datum. CURRENT VERSION0.47 QUALIFICATION: Finish the same-D/common-family adelic setting after ASG: normalized nonzero Bessel W and decay, measured-field block, original Siegel-Weil rescaling and measure/character ties, Bessel uniqueness/local ties, then the original AutTie/RowDictTie/LabTie. RowsFinDim is now proved universally by RFD. Supply the actual CountIdempotentsClass theorem or its three sufficient arithmetic cylinder counts, with one simultaneous family threshold; reuse the completed ARO structure. Completed new original producers: ARO1-ARO30, RFD1-RFD9, ASG1-ASG18. Any earlier narrative naming these completed producers as outstanding is historical; this qualification is the current scope authority. CURRENT VERSION0.48 QUALIFICATION: Finish the same-D/common-family adelic setting after ASG and AMD: normalized nonzero Bessel W and decay, original orbital/fiber and Siegel-Weil measure/character ties, Bessel uniqueness/local ties, and the full original transport theorem, then AutTie/RowDictTie/LabTie. The complete MExists measured-data block and unipotent covolume normalization are now supplied, and universal RowsFinDim remains available. Supply the actual CountIdempotentsClass theorem or its three sufficient arithmetic cylinder counts with one simultaneous family threshold. Completed current producers or conditional consumers: ARO1-ARO30, RFD1-RFD9, ASG1-ASG18, AMD1-AMD17. Earlier frontier descriptions are historical; this qualification is the current scope authority. Original whole-contract premises and quantifiers are retained. CURRENT VERSION0.49 QUALIFICATION: Prove the exact original BesselNVGap nonzero integral on its supplied discrete-series, compact-representation and character data. CoefWeightedLoc, WlocExists with its original nonvanishing antecedent, WFromLoc, the normalized global product and all-later-measured-block decay are complete. On the inherited Forms0 scope this leaves BesselNVGap for ConcreteSettingExists. Continue original orbital/fiber and Siegel-Weil measure/character ties, the required Bessel uniqueness/local ties and the full original transport theorem, then AutTie/RowDictTie/LabTie. Supply the actual CountIdempotentsClass theorem or its three sufficient arithmetic cylinder counts with one simultaneous family threshold. Completed current producers or conditional consumers: ARO1-ARO30, RFD1-RFD9, ASG1-ASG18, AMD1-AMD17, BWP1-BWP24 at their exact BesselNVGap/Forms0 scope. Earlier frontier descriptions are historical; this qualification is the current scope authority. Original whole-contract premises and quantifiers are retained.

\emph{Shared objects and choice order.} one FamP; one class $\mathcal A$ and its target-compatible subclass; one coefficient field per family datum; one x and i at each application; original selected rows, words and levels; thresholds chosen before the common family; AllRows A; LabelCountGoodInjective; AllJoinA; CountAllRows

\emph{New completed inputs.} ARO1-ARO30, RFD1-RFD9, ASG1-ASG18, AMD1-AMD17, BWP1-BWP24 at their exact BesselNVGap/Forms0 scope

\emph{Current remaining subfrontier.} Prove the exact original BesselNVGap nonzero integral on its supplied discrete-series, compact-representation and character data. CoefWeightedLoc, WlocExists with its original nonvanishing antecedent, WFromLoc, the normalized global product and all-later-measured-block decay are complete. On the inherited Forms0 scope this leaves BesselNVGap for ConcreteSettingExists. Continue original orbital/fiber and Siegel-Weil measure/character ties, the required Bessel uniqueness/local ties and the full original transport theorem, then AutTie/RowDictTie/LabTie. Supply the actual CountIdempotentsClass theorem or its three sufficient arithmetic cylinder counts with one simultaneous family threshold.

\emph{Source.} \href{https://github.com/CocoSgt/riemann-hypothesis-lean/blob/7cbf7a0c32fca1badf05931879a37c3f8a2b03b4/src/Formalism/SharpRestriction/GlobalTheoremInputs.lean#L55-L58}{LinkIRR}.

\emph{Source.} \href{https://github.com/CocoSgt/riemann-hypothesis-lean/blob/7cbf7a0c32fca1badf05931879a37c3f8a2b03b4/src/Development/AbelianVarietyTheory/Proofs/GlobalInterpolation/LinkIRROfJoinEx.lean#L17-L38}{linkIRR\_\allowbreak of\_\allowbreak join\_\allowbreak ex}.

\emph{Source.} \href{https://github.com/CocoSgt/riemann-hypothesis-lean/blob/7cbf7a0c32fca1badf05931879a37c3f8a2b03b4/src/Development/AbelianVarietyTheory/Proofs/GlobalInterpolation/LinkIRROfJoinEx.lean#L143-L166}{link\_\allowbreak irr}.



\paragraph{3. GLOBAL\_\allowbreak PERIOD\_\allowbreak HEIGHT — C — Full actual period and height comparison LinkDIVt}

Closed conditional TuSize -> DGamSize -> (LinkEta -> DIVinN -> LinkDIV), with every original LinkDIV\_\allowbreak at clause.

\emph{Shared objects and choice order.} one FamP; one class $\mathcal A$ and its target-compatible subclass; one coefficient field per family datum; one x and i at each application; original selected rows, words and levels; thresholds chosen before the common family; same RR target; same Bu/besS; same Nu and Bk

\emph{Source.} \href{https://github.com/CocoSgt/riemann-hypothesis-lean/blob/7cbf7a0c32fca1badf05931879a37c3f8a2b03b4/src/Formalism/SharpRestriction/GlobalTheoremInputs.lean#L77-L90}{LinkDIV\_\allowbreak at}.

\emph{Source.} \href{https://github.com/CocoSgt/riemann-hypothesis-lean/blob/7cbf7a0c32fca1badf05931879a37c3f8a2b03b4/src/Formalism/SharpRestriction/GlobalTheoremInputs.lean#L94-L96}{LinkDIVt}.

\emph{Source.} \href{https://github.com/CocoSgt/riemann-hypothesis-lean/blob/7cbf7a0c32fca1badf05931879a37c3f8a2b03b4/src/Development/AbelianVarietyTheory/Proofs/GlobalInterpolation/LinkBrest.lean#L7-L25}{GlobalTheory.Top.link\_\allowbreak brest}.

\emph{Source.} \href{https://github.com/CocoSgt/riemann-hypothesis-lean/blob/7cbf7a0c32fca1badf05931879a37c3f8a2b03b4/src/Development/AbelianVarietyTheory/Proofs/GlobalInterpolation/WireSel.lean#L11-L41}{GlobalTheory.Top.wire\_\allowbreak sel}.



\paragraph{4. GLOBAL\_\allowbreak SPECTRAL — C — Original normalized kernel/Gram estimates LinkS}

Closed TuSize -> LinkS with exact Gamma normalization, off-diagonal Hilbert-Schmidt bound, diagonal error and Dall/Dstar estimate, errors exp(-x\textasciicircum{}(epsilon/8)). CURRENT VERSION0.48 QUALIFICATION: Complete the original same-family spectral/Bessel and Siegel-Weil inputs with all measure, character, transport and common-object ties. The complete MExists measured block and prescribed-measure covolume normalization are supplied; these are not the full normalized spectral theorem. Completed current producers or conditional consumers: ASG1-ASG18, AMD1-AMD17, GCF1-GCF18 at their stated scopes. Earlier frontier descriptions are historical; this qualification is the current scope authority. Original whole-contract premises and quantifiers are retained. CURRENT VERSION0.49 QUALIFICATION: Complete the original same-family spectral/Bessel and Siegel-Weil inputs with every measure, character, transport and common-object tie. The normalized packet and decay are now constructed conditional on the exact BesselNVGap, and the analytic setting uses the original supplied Forms0 scope. These do not imply the full normalized spectral theorem. Completed current producers or conditional consumers: ASG1-ASG18, AMD1-AMD17, GCF1-GCF18 at their stated scopes, BWP1-BWP24 at their exact BesselNVGap/Forms0 scope, WRS1-WRS10 original local gain. Earlier frontier descriptions are historical; this qualification is the current scope authority. Original whole-contract premises and quantifiers are retained.

\emph{Shared objects and choice order.} one FamP; one class $\mathcal A$ and its target-compatible subclass; one coefficient field per family datum; one x and i at each application; original selected rows, words and levels; thresholds chosen before the common family; actual KD; same S.Cu spaces; GamOf; Dstar; Dall

\emph{New completed inputs.} ASG1-ASG18, AMD1-AMD17, GCF1-GCF18 at their stated scopes, BWP1-BWP24 at their exact BesselNVGap/Forms0 scope, WRS1-WRS10 original local gain

\emph{Current remaining subfrontier.} Complete the original same-family spectral/Bessel and Siegel-Weil inputs with every measure, character, transport and common-object tie. The normalized packet and decay are now constructed conditional on the exact BesselNVGap, and the analytic setting uses the original supplied Forms0 scope. These do not imply the full normalized spectral theorem.

\emph{Source.} \href{https://github.com/CocoSgt/riemann-hypothesis-lean/blob/7cbf7a0c32fca1badf05931879a37c3f8a2b03b4/src/Formalism/SharpRestriction/GlobalTheoremInputs.lean#L145-L149}{LinkS\_\allowbreak at}.

\emph{Source.} \href{https://github.com/CocoSgt/riemann-hypothesis-lean/blob/7cbf7a0c32fca1badf05931879a37c3f8a2b03b4/src/Formalism/SharpRestriction/GlobalTheoremInputs.lean#L152-L157}{LinkS}.

\emph{Source.} \href{https://github.com/CocoSgt/riemann-hypothesis-lean/blob/7cbf7a0c32fca1badf05931879a37c3f8a2b03b4/src/Development/AbelianVarietyTheory/Proofs/GlobalInterpolation/LinkAlt.lean#L16-L50}{GlobalTheory.Top.link\_\allowbreak alt}.



\paragraph{5. AUT\_\allowbreak UNRAM\_\allowbreak COEF — C — Actual unramified-coefficient automorphy}

For each actual Dat with coefficient ramification one and fixed iota, choose one marked g1 supplying FinA, AutB and LemW together. CURRENT VERSION0.48 QUALIFICATION: Construct the actual AutomorphyUnramifiedCoefficients package: for every original Dat with coefficient ramification one and fixed iota, choose one marked g1 that supplies FinA, AutB and LemW together. The original InpZ record, full RigII, CNR3, finite PropGK consumer and now the entire OneChar implication are supplied. Continue RH Cases I/V and the independent additional Dirichlet-GRH CaseIII/promodularity, SeedIV, tensor-factor and rank-four-chain tasks. Completed current producers or conditional consumers: OCH1-OCH30 (the full consumer implication), KIZ original InpZ, KTR1-KTR13, RNC16, PGK1-PGK6. Earlier frontier descriptions are historical; this qualification is the current scope authority. Original whole-contract premises and quantifiers are retained. CURRENT VERSION0.49 QUALIFICATION: Construct one common GdEnd for each original datum and fixed iota: the same marked g1 must supply FinEnd, AutEnd3 and AutEnd2. UCE now proves all ordinary local/curve/coefficient consumers and the exact implication to FinA/AutB/LemW, including arbitrary entry maps and every original outer-corner extension. The producer still needs the actual common abelian-CM GdFrameEx and LinkedEndE with both rank-two and rank-three runs, full ModUAlt/PfromAut and the two actual seeds, followed by the original local/automorphic descent. At this pin the alternative upstream route gives GdEndAlt; its additional local/descent fields must be matched to the actual original points, so no name-only GdEndAlt-to-GdEnd identification is allowed. Keep the supplied InpZ, RigII, CNR3, finite PropGK and whole OneChar implication. Continue RH Cases I/V and the independent additional Dirichlet-GRH CaseIII/promodularity, SeedIV, tensor-factor and rank-four-chain tasks. Completed current producers or conditional consumers: UCE1-UCE19 (complete GdEnd consumer), OCH1-OCH30, KIZ original InpZ, KTR1-KTR13, RNC16, PGK1-PGK6. Earlier frontier descriptions are historical; this qualification is the current scope authority. Original whole-contract premises and quantifiers are retained.

\emph{New completed inputs.} UCE1-UCE19 (complete GdEnd consumer), OCH1-OCH30, KIZ original InpZ, KTR1-KTR13, RNC16, PGK1-PGK6

\emph{Current remaining subfrontier.} Construct one common GdEnd for each original datum and fixed iota: the same marked g1 must supply FinEnd, AutEnd3 and AutEnd2. UCE now proves all ordinary local/curve/coefficient consumers and the exact implication to FinA/AutB/LemW, including arbitrary entry maps and every original outer-corner extension. The producer still needs the actual common abelian-CM GdFrameEx and LinkedEndE with both rank-two and rank-three runs, full ModUAlt/PfromAut and the two actual seeds, followed by the original local/automorphic descent. At this pin the alternative upstream route gives GdEndAlt; its additional local/descent fields must be matched to the actual original points, so no name-only GdEndAlt-to-GdEnd identification is allowed. Keep the supplied InpZ, RigII, CNR3, finite PropGK and whole OneChar implication. Continue RH Cases I/V and the independent additional Dirichlet-GRH CaseIII/promodularity, SeedIV, tensor-factor and rank-four-chain tasks.

\emph{Source.} \href{https://github.com/CocoSgt/riemann-hypothesis-lean/blob/7cbf7a0c32fca1badf05931879a37c3f8a2b03b4/src/Development/AutomorphicArithmetic/AutomorphicVarieties/AutomorphicCases/AutomorphyUnramifiedCoefficients.lean#L24-L27}{AutomorphyUnramifiedCoefficients}.

\emph{Source.} \href{https://github.com/CocoSgt/riemann-hypothesis-lean/blob/7cbf7a0c32fca1badf05931879a37c3f8a2b03b4/src/Development/AutomorphicArithmetic/AutomorphicVarieties/AutomorphicCases/AutomorphyUnramifiedCoefficients.lean#L12-L13}{FinA}.

\emph{Source.} \href{https://github.com/CocoSgt/riemann-hypothesis-lean/blob/7cbf7a0c32fca1badf05931879a37c3f8a2b03b4/src/Development/AutomorphicArithmetic/AutomorphicVarieties/AutomorphicCases/AutomorphyUnramifiedCoefficients.lean#L15-L17}{AutB}.

\emph{Source.} \href{https://github.com/CocoSgt/riemann-hypothesis-lean/blob/7cbf7a0c32fca1badf05931879a37c3f8a2b03b4/src/Development/AutomorphicArithmetic/AutomorphicVarieties/AutomorphicCases/AutomorphyUnramifiedCoefficients.lean#L19-L22}{LemW}.



\paragraph{6. TASK\_\allowbreak CASE3\_\allowbreak AUT\_\allowbreak JOIN — B — Complete the original Case-III automorphy/image join}

Install the actual Case-III ProModAlt, SeedIVAlt and shared PMExAlt with the inherited CohomologyHyp/Ferrero-Washington/HFit/SeedII facts into the original whole Case-III automorphy/image conclusion.

\emph{Source.} \href{https://github.com/CocoSgt/riemann-hypothesis-lean/blob/7cbf7a0c32fca1badf05931879a37c3f8a2b03b4/src/Development/AutomorphicArithmetic/AutomorphicVarieties/ModularityArgument/ThresholdArguments/DenseThresholdCases/PromodularityCaseThree.lean#L43-L44}{AutomorphyRankFourDense}.

\emph{Source.} \href{https://github.com/CocoSgt/riemann-hypothesis-lean/blob/7cbf7a0c32fca1badf05931879a37c3f8a2b03b4/src/Development/AutomorphicArithmetic/AutomorphicVarieties/ModularityArgument/ThresholdArguments/DenseThresholdCases/CuFinal.lean#L44-L44}{main4\_\allowbreak alt}.



\paragraph{7. TASK\_\allowbreak CASE3\_\allowbreak PROMOD\_\allowbreak JOIN — C — Reconstruct the actual Case-III chain-to-promodularity proof}

Use the original tensor-factor tower, ChainRankFour, ChainDesc4, seeded rank-four chain reduction, actual PMExAlt and SeedIVAlt to construct ProModAlt on every original Ctx3 star point at the same fixed iota, preserving all finite-coefficient, inertia and support maps.

\emph{Source.} \href{https://github.com/CocoSgt/riemann-hypothesis-lean/blob/7cbf7a0c32fca1badf05931879a37c3f8a2b03b4/src/Development/AutomorphicArithmetic/AutomorphicVarieties/ModularityArgument/ThresholdArguments/DenseThresholdCases/PromodularityCaseThreeOfLemmas.lean#L131-L153}{UpperTheory.V.Cu.promodularity\_\allowbreak case\_\allowbreak three\_\allowbreak of\_\allowbreak chain\_\allowbreak reduction}.

\emph{Source.} \href{https://github.com/CocoSgt/riemann-hypothesis-lean/blob/7cbf7a0c32fca1badf05931879a37c3f8a2b03b4/src/Development/AutomorphicArithmetic/AutomorphicVarieties/ModularityArgument/ThresholdArguments/DenseThresholdCases/PromodularityCaseThree.lean#L40-L42}{PromodularityCaseThreeDense}.



\paragraph{8. SEED\_\allowbreak IV\_\allowbreak ALT — C — Actual Case-III seed with prechosen residual datum}

Base Ctx3/Four tower, iota and Amb -> w,s before every later Dies tower -> compatible finite-coefficient seed with exact residual support/root entries.

\emph{Source.} \href{https://github.com/CocoSgt/riemann-hypothesis-lean/blob/7cbf7a0c32fca1badf05931879a37c3f8a2b03b4/src/Development/AutomorphicArithmetic/AutomorphicVarieties/ModularityArgument/ThresholdArguments/DenseThresholdCases/PromodularityCaseThree.lean#L19-L25}{SeedIVAlt}.



\paragraph{9. TENSOR\_\allowbreak FACTOR\_\allowbreak TOWER — B — Actual tensor-factor tower and characters}

Original finite-field, tame, factor and character constraints; source-level construction is sizeable and has not been reconstructed here.

\emph{Source.} \href{https://github.com/CocoSgt/riemann-hypothesis-lean/blob/7cbf7a0c32fca1badf05931879a37c3f8a2b03b4/src/Development/AutomorphicArithmetic/AutomorphicVarieties/ModularityArgument/ThresholdArguments/DenseThresholdCases/PromodularityCaseThreeOfLemmas.lean#L131-L153}{UpperTheory.V.Cu.promodularity\_\allowbreak case\_\allowbreak three\_\allowbreak of\_\allowbreak chain\_\allowbreak reduction}.



\paragraph{10. CHAIN\_\allowbreak RANK\_\allowbreak FOUR — C — Actual SomeIn-to-AllIn propagation}

Same source Space and PM ideal, residual point Good4, original bound 30(C+1)\textasciicircum{}2 <= d; geometry must preserve all identities.

\emph{Source.} \href{https://github.com/CocoSgt/riemann-hypothesis-lean/blob/7cbf7a0c32fca1badf05931879a37c3f8a2b03b4/src/Development/AutomorphicArithmetic/AutomorphicVarieties/ModularityArgument/ThresholdArguments/DenseThresholdCases/PromodularityCaseThree.lean#L26-L32}{ChainRankFour}.



\paragraph{11. CHAIN\_\allowbreak DESC4 — B — Original residual chain descent}

Finite residual descent retaining source Chain4 compatibility.

\emph{Source.} \href{https://github.com/CocoSgt/riemann-hypothesis-lean/blob/7cbf7a0c32fca1badf05931879a37c3f8a2b03b4/src/Development/AutomorphicArithmetic/AutomorphicVarieties/ModularityArgument/ThresholdArguments/DenseThresholdCases/PromodularityCaseThreeOfLemmas.lean#L131-L153}{UpperTheory.V.Cu.promodularity\_\allowbreak case\_\allowbreak three\_\allowbreak of\_\allowbreak chain\_\allowbreak reduction}.



\paragraph{12. RANK4\_\allowbreak CHAIN\_\allowbreak REDUCTION — C — Seeded reduction to the original residual chain}

Select T1 before every later Deep tower; retain one compatible seed and Chain4 for every original residual lift over Q.

\emph{Source.} \href{https://github.com/CocoSgt/riemann-hypothesis-lean/blob/7cbf7a0c32fca1badf05931879a37c3f8a2b03b4/src/Development/AutomorphicArithmetic/AutomorphicVarieties/ModularityArgument/ThresholdArguments/DenseThresholdCases/PromodularityCaseThree.lean#L33-L39}{RankFourChainReductionDense}.

\emph{Source.} \href{https://github.com/CocoSgt/riemann-hypothesis-lean/blob/7cbf7a0c32fca1badf05931879a37c3f8a2b03b4/src/Development/AutomorphicArithmetic/AutomorphicVarieties/ModularityArgument/ThresholdArguments/DenseThresholdCases/CuPChainIV.lean#L16-L82}{UpperTheory.V.Cu.chainIV\_\allowbreak alt}.



\paragraph{13. TASK\_\allowbreak COUNT\_\allowbreak V\_\allowbreak INSTALL — B — Complete the original CountV finite assembly}

From exactly the listed source host/cohomology inputs, complete countV\_\allowbreak of\_\allowbreak finiteness\_\allowbreak duality\_\allowbreak euler\_\allowbreak host, including original host/frame formation, normalized length, the CaseV guard and R1<=2\textasciicircum{}c2.

\emph{Source.} \href{https://github.com/CocoSgt/riemann-hypothesis-lean/blob/7cbf7a0c32fca1badf05931879a37c3f8a2b03b4/src/Development/AutomorphicArithmetic/AutomorphicVarieties/HeightAndVanishing/Counts.lean#L69-L69}{CountV}.

\emph{Source.} \href{https://github.com/CocoSgt/riemann-hypothesis-lean/blob/7cbf7a0c32fca1badf05931879a37c3f8a2b03b4/src/Development/AutomorphicArithmetic/AutomorphicVarieties/ModularityArgument/ThresholdArguments/DenseThresholdCases/CountVOf.lean#L33-L60}{UpperTheory.V.UnitaryThreeInertia.countV\_\allowbreak of\_\allowbreak finiteness\_\allowbreak duality\_\allowbreak euler\_\allowbreak host}.



\paragraph{14. THM\_\allowbreak S6\_\allowbreak S — C — Original unitary-three ordinary support and automorphy}

Original host/frame and fixed iota -> one marked g1 -> every universal R,u -> weight-zero finiteness and every Type-I BPrime\_\allowbreak s point.

\emph{Source.} \href{https://github.com/CocoSgt/riemann-hypothesis-lean/blob/7cbf7a0c32fca1badf05931879a37c3f8a2b03b4/src/Development/AutomorphicArithmetic/AutomorphicVarieties/ModularityArgument/ThresholdArguments/DenseThresholdCases/ThmS6CoreS.lean#L17-L26}{ThmS6\_\allowbreak s}.

\emph{Source.} \href{https://github.com/CocoSgt/riemann-hypothesis-lean/blob/7cbf7a0c32fca1badf05931879a37c3f8a2b03b4/src/Development/AutomorphicArithmetic/AutomorphicVarieties/ModularityArgument/ThresholdArguments/DenseThresholdCases/ThmS6CoreS.lean#L12-L16}{BPrime\_\allowbreak s}.



\paragraph{15. HOST\_\allowbreak FRAME\_\allowbreak LENGTH — B — Original Case-V host/frame and normalized-length installation}

Actual HostData, frame, local-condition equality, abstract count and normalized original Selmer length must all refer to the same objects.

\emph{Source.} \href{https://github.com/CocoSgt/riemann-hypothesis-lean/blob/7cbf7a0c32fca1badf05931879a37c3f8a2b03b4/src/Development/AutomorphicArithmetic/AutomorphicVarieties/ModularityArgument/ThresholdArguments/DenseThresholdCases/CountVOf.lean#L33-L60}{UpperTheory.V.UnitaryThreeInertia.countV\_\allowbreak of\_\allowbreak finiteness\_\allowbreak duality\_\allowbreak euler\_\allowbreak host}.



\paragraph{16. HOST\_\allowbreak NUMBER\_\allowbreak S — C — Uniform original-host rank-three catalogue}

One exponent before all closures, coefficient rings, levels, characters, frames, embeddings and marked points; one separating trace function controls all later finite pairwise nonisomorphic collections.

\emph{Source.} \href{https://github.com/CocoSgt/riemann-hypothesis-lean/blob/7cbf7a0c32fca1badf05931879a37c3f8a2b03b4/src/Development/AutomorphicArithmetic/AutomorphicVarieties/ModularityArgument/ThresholdArguments/DenseThresholdCases/HostNumberS.lean#L18-L28}{HostNumber\_\allowbreak s}.



\paragraph{17. TASK\_\allowbreak HOST\_\allowbreak VAN\_\allowbreak E\_\allowbreak INSTALL — B — Transport K/duality to literal HostVanE}

Use the already supplied LemmaK/local-global duality and finiteness on the actual CountV host/frame and exact original coefficient modules.

\emph{Source.} \href{https://github.com/CocoSgt/riemann-hypothesis-lean/blob/7cbf7a0c32fca1badf05931879a37c3f8a2b03b4/src/Development/AutomorphicArithmetic/AutomorphicVarieties/UnitaryThreeDeformations/HostNumber.lean#L9-L12}{HostVanE}.

\emph{Source.} \href{https://github.com/CocoSgt/riemann-hypothesis-lean/blob/7cbf7a0c32fca1badf05931879a37c3f8a2b03b4/src/Development/AutomorphicArithmetic/AutomorphicVarieties/LevelCohomology/SoHostVanE.lean#L11-L18}{UpperTheory.V.X.So.hostVanE}.



\paragraph{18. TASK\_\allowbreak HOST\_\allowbreak XPP\_\allowbreak INSTALL — B — Transport CorX to literal HostXpp}

Preserve Psi\textasciicircum{}-2 eta, the prime-to-ell quotient, actual local conditions and the original CountV host while installing BR7.

\emph{Source.} \href{https://github.com/CocoSgt/riemann-hypothesis-lean/blob/7cbf7a0c32fca1badf05931879a37c3f8a2b03b4/src/Development/AutomorphicArithmetic/AutomorphicVarieties/UnitaryThreeDeformations/HostNumber.lean#L14-L17}{HostXpp}.

\emph{Source.} \href{https://github.com/CocoSgt/riemann-hypothesis-lean/blob/7cbf7a0c32fca1badf05931879a37c3f8a2b03b4/src/Development/AutomorphicArithmetic/AutomorphicVarieties/LevelCohomology/Xpp.lean#L11-L84}{UpperTheory.V.X.xpp}.



\paragraph{19. DESIGN\_\allowbreak CHOICES\_\allowbreak COMMON — C — Construct the original common design family and bundle}

For the one target HostExW, construct the original compatible DesignChoices on Cdes=WithTw(WithSm Cprime), where Cprime=WithChN(A intersect target class), A is the one Built construction class, and cT(eps,y)=cTmin(y). Retain every original field, tuple, word, metric and common bundle duty. DSZ10 and RMP11 give size/rank properties of class members, not member existence. The original RecExP/RecAtG and common object/row route remains coupled. Construct ONE original common design family R and its recursive/geometric half with original Inputs/Pins/Holds/ExAt/InGeo and required CoreP. PKT now reconstructs full exists\_\allowbreak pkT, rowPk\_\allowbreak nonempty and exists\_\allowbreak moved with Copy/Coh on those original given inputs and a supplied fitting later stage: one catalogue, both ChT, complete LfT, actual integral lattices and all metric/measure/period constants are retained; one numerical record precedes every eta and V/S/N/G. RDR supplies one Side0 before all readers, ReadPk/pins and full ReadL with exact o.eta=q.eta. The initial actual geometry, Fits/division/central conditions, primitive rational tests and eta agreement must still be produced. Inherited CSR supplies finite columns, one ReadC before later twists, period/evaluation and scalar/error transports, RowSideDetermined and same-R Bundle/BTw/BWit/BunU consumers. RPI supplies exact RowPairInt for actual pairs and finite translates with original mu\_\allowbreak b/fd0; first literal DivDef3 at the same updated R,F,Y,Z,C,sEE follows from its retained original readings, Pins/Holds, stored nonzero IsItW chain, five real-group predicates and group tie. Remaining common fields include host/frame, MultEq nonzero-part support, that original nonzero chain/group tie, eigensystem/carrier/parts and all-embedding logLam, HD/d3b, NEb, actual rational tests and base rows. Retain every simultaneous family hypothesis and eventual threshold. Stronger second DivRows3 remains separate; universal quotient-height CuspRapidDecay and actual common member existence are not inferred from the comparison constructors. Version 0.41 supplies the finite phase-correct rational theta route: exact coefficient extension, split representatives, square-law/PhiIs, Pins from CoversAt, class-sum/point-identity transfer to OneMapOverL, actual-good-projector RatG and both-level LiftTestB. Its original higher covers, all-tuple class-sum uniqueness including phase rescaling, one PointIdentCore before O, base ClassSumsSpan and geometric LatFull remain explicit producers. MLat is available from the original LatFull input. This sharpens the rational-test work without claiming actual common geometry, Fits or eta agreement. Version0.46 proves finite generation of the original whole integral cusp lattice by actual integral-frame evaluation. It also proves the exact identity-cover, analytic unfolding, original-field coefficient descent, global clearing-scalar and both-level lattice-transport implications. Full spanning retains the supplied original IsNab connection and the five explicit principal-chart, level-system, raw-Gaussian, short-rule/dictionary and ramified-jet producers. The higher tuple covers, actual common class sums, PointIdentCore and ClassSumsSpan remain separate. CURRENT VERSION0.47 QUALIFICATION: Construct one original common geometric family and the higher CoversAt/all-tuple class sums and uniqueness with original phase rescalings, one PointIdentCore with common lambda,z, and base ClassSumsSpan. CNC now supplies the full original IsNab; PCL supplies the conditional trace-dual principal-lattice and quantitative normalization. Actual uniformization/HasPrinc and the five principal-chart, finite-level-system, Gaussian/reproduction, short-rule/simultaneous-dictionary and ramified-jet producer groups remain. Keep all original both-level, coefficient-field, Fits and eta agreements. Completed new original producers: ICL1-ICL21, CNC1-CNC16, PCL1-PCL12. Any earlier narrative naming these completed producers as outstanding is historical; this qualification is the current scope authority. CURRENT VERSION0.48 QUALIFICATION: Construct the original finite holomorphic FamilyUnif/N1c, normalized theta families and local lifting, then the common principal charts and finite-level theta systems. PUN now supplies the exact fixed-space models, all-family center lattices, A=1 exponent matching and individual primitive points with the full point-equality criterion; CNC supplies the full canonical connection and PCL the quantitative trace-dual normalization. Complete higher CoversAt and all-tuple class sums/uniqueness with the original phase rescalings, one PointIdentCore with common lambda,z, and base ClassSumsSpan. The Gaussian/reproduction, short-rule/simultaneous-dictionary and ramified-jet producers, and all original both-level, coefficient-field, Fits and eta agreements remain explicit. Completed current producers or conditional consumers: ICL1-ICL21, CNC1-CNC16, PCL1-PCL12, PUN1-PUN24, GCF1-GCF18 at their stated scopes. Earlier frontier descriptions are historical; this qualification is the current scope authority. Original whole-contract premises and quantifiers are retained. CURRENT VERSION0.49 QUALIFICATION: Construct the adjoining simultaneous principal framed points and test-ring dictionary. Full finite holomorphic FamilyUnif/N1c, normalized theta families and local lifts are now supplied by HUF. The next bounded target is allTest\_\allowbreak of\_\allowbreak secBaseChange: represent FrPt by the Lie-eigenframe bundles and theta-section frame line, use the full n=1 SecBaseChangeSt on every nonflat coefficient base change, and prove finite-generated subalgebra descent with the exact TRing/FromTest transport. Then supply the full FramesOK/CoverTri/FibTri and IsUnf/hasPrinc package, including adapted mod-eight symplectic basis, normalized theta multiplier/weight and arithmetic invariance. Continue higher CoversAt, all-tuple class sums and uniqueness with original phase rescalings, one PointIdentCore with common lambda,z, base ClassSumsSpan, Gaussian/reproduction, short-rule dictionaries and ramified jets, retaining both-level, coefficient-field, Fits and eta agreements. Common-family nonemptiness remains distinct from compatibility of an already supplied member. Completed current producers or conditional consumers: ICL1-ICL21, CNC1-CNC16, PCL1-PCL12, PUN1-PUN24, HUF1-HUF21 (full original FamilyUnif/N1c), GCF1-GCF18 and WRS1-WRS10 at their stated scopes. Earlier frontier descriptions are historical; this qualification is the current scope authority. Original whole-contract premises and quantifiers are retained.

\emph{New completed inputs.} ICL1-ICL21, CNC1-CNC16, PCL1-PCL12, PUN1-PUN24, HUF1-HUF21 (full original FamilyUnif/N1c), GCF1-GCF18 and WRS1-WRS10 at their stated scopes

\emph{Current remaining subfrontier.} Construct the adjoining simultaneous principal framed points and test-ring dictionary. Full finite holomorphic FamilyUnif/N1c, normalized theta families and local lifts are now supplied by HUF. The next bounded target is allTest\_\allowbreak of\_\allowbreak secBaseChange: represent FrPt by the Lie-eigenframe bundles and theta-section frame line, use the full n=1 SecBaseChangeSt on every nonflat coefficient base change, and prove finite-generated subalgebra descent with the exact TRing/FromTest transport. Then supply the full FramesOK/CoverTri/FibTri and IsUnf/hasPrinc package, including adapted mod-eight symplectic basis, normalized theta multiplier/weight and arithmetic invariance. Continue higher CoversAt, all-tuple class sums and uniqueness with original phase rescalings, one PointIdentCore with common lambda,z, base ClassSumsSpan, Gaussian/reproduction, short-rule dictionaries and ramified jets, retaining both-level, coefficient-field, Fits and eta agreements. Common-family nonemptiness remains distinct from compatibility of an already supplied member.



\paragraph{20. LOCAL\_\allowbreak NBUILT\_\allowbreak COMMON — C — Actual local specialization on the common selected family}

Construct, with the required closed-class property on the same Des(C,3), one eventual family of actual NBuilt certificates for the original Fm.D(x,i), retaining the certificate before forgetting to TSel. HLP now reconstructs full HeckeLift0, all-index commutation, HeckeStd0 and HeckeIntTw from the original MtInput and complex Hecke facts, in the full ambient GenAt range k<=nV. Every maximal ideal is an actual supported residual character class; literal MtSub is the O-span of projections of the original MtLat generators, is finite free, and localizes onto the exact NCov.part. Literal Psi has its exact range and off-space projection error, both original span assertions hold, and ActEx/ActHK are constructed on the same lattice under original semisimplicity. Original MjLat relation quotients, universal coefficient-map relation fields, selected tame generators, coefficient embeddings and all identity transports give the complete nonempty\_\allowbreak spec constructor on its full printed input list. A single compositum combines the four original coefficient requirements with big fixed beforehand. Whole ndSpec\_\allowbreak exists still needs actual common geometric/choice MtInput, complex HeckeStdPlain and the no-Jordan producer, original F3 bridges with valuation/GoodCoef/iq, original F4 common-overfield reading of the entire coefficient field, and exact arithmetic character/reciprocity/admissibility suppliers. Original AutSpec, Galois specialization and the remaining NBuilt fields on that same installation also remain. Do not replace the actual projected lattice or maps by abstract objects with the same dimensions.



\paragraph{21. SECOND\_\allowbreak LAYER\_\allowbreak COMPATIBLE\_\allowbreak SELECTION — C — Original second-layer selector and class duties}

For the same A, supply Duty on A and a Closed common original TSel family whose selected records are the original true-column rec7 records (or have the entire dependent tuple transported), with their real witness and original coefficient-field root data. Existence/compatibility remains arithmetic; merely abstractly isomorphic fields or independently selected columns do not suffice. CSR21–22 now construct the same-R Bundle/BTw/BWit/BunU outputs from the exact selected realization/row-side packet and NEb, with the original EvP/ParEv quantifier order. Actual common selection, the five stated realization packets and NEb remain upstream; none is replaced by an independently chosen compatible record. Version 0.38 closes the selected analytic extension and reconstructs the first updated DivDef3 consumer; it leaves the original common selection, geometric readings, rational tests, base rows, host/parts/NEb and exact retained DI premises upstream, with unchanged EvP/ParEv thresholds. Version 0.39 supplies the original given-data reader constructors and exact RowPairInt/first DivDef3 on original premises. It leaves one actual common geometric/comparison/test member, eta agreement, compatible selection, rational tests, base rows, host/parts/NEb and original EvP/ParEv thresholds upstream. Version 0.40 also supplies original comparison/RowPk/fitting-stage Moved algorithms and exact common data transport. The actual initial geometric member, fitting tests, eta agreement and compatible second-layer selector with original eventual thresholds remain upstream. Version 0.41 supplies the finite phase-correct rational theta route: exact coefficient extension, split representatives, square-law/PhiIs, Pins from CoversAt, class-sum/point-identity transfer to OneMapOverL, actual-good-projector RatG and both-level LiftTestB. Its original higher covers, all-tuple class-sum uniqueness including phase rescaling, one PointIdentCore before O, base ClassSumsSpan and geometric LatFull remain explicit producers. MLat is available from the original LatFull input. This sharpens the rational-test work without claiming actual common geometry, Fits or eta agreement. CURRENT VERSION0.49 QUALIFICATION: Continue the distinct global denominator, period/height, normalized spectral and second-layer compatible-selection producers, together with existence of the original common arithmetic/geometric family. GLOBAL\_\allowbreak LOCAL\_\allowbreak GAIN and its actual GroupCompatibilityFamily supplier are now complete on the unchanged selected built class; no additional WPt/TSel realization is needed for that estimate. Do not reopen local first-growth, mean, auxiliary admission, selected D3 tops/Adm2F, the single GIs2 stage or the exact comparison/kernel transport. A full RecT or TSel may remain useful for other tasks, but is not claimed by this local-gain closure. Completed current producers or conditional consumers: GCF1-GCF18 at their stated scopes, WRS1-WRS10 original selected-class compatibility/local gain. Earlier frontier descriptions are historical; this qualification is the current scope authority. Original whole-contract premises and quantifiers are retained.

\emph{New completed inputs.} GCF1-GCF18 at their stated scopes, WRS1-WRS10 original selected-class compatibility/local gain

\emph{Current remaining subfrontier.} Continue the distinct global denominator, period/height, normalized spectral and second-layer compatible-selection producers, together with existence of the original common arithmetic/geometric family. GLOBAL\_\allowbreak LOCAL\_\allowbreak GAIN and its actual GroupCompatibilityFamily supplier are now complete on the unchanged selected built class; no additional WPt/TSel realization is needed for that estimate. Do not reopen local first-growth, mean, auxiliary admission, selected D3 tops/Adm2F, the single GIs2 stage or the exact comparison/kernel transport. A full RecT or TSel may remain useful for other tasks, but is not claimed by this local-gain closure.



\paragraph{22. BASE\_\allowbreak THETA\_\allowbreak READINGS — C — Actual base reading on the common bundle}

For every family in the original class, eventually on the SAME selected Bundle and for each i<m, supply the actual R,F,Read,ReadL, weak models/Weil data and phase-one pins. Identify their LiftTie.theta, mtRead, mjRead and ThDef3 basis with that Bundle.P.A3, P.Zs, P.S0 and its stored maps. Exact field, embeddings, metrics, measures, row family and weak choices are retained. This remains one part of the coupled BTH12 producer. PKT now supplies the complete original comparison packet, RowPk and fitting-stage Moved with Copy/Coh on original given geometry, Inputs/Pins/Holds/ExAt/InGeo and Fits. Its ChT/LfT/Reps, actual lattice carriers, original Haar and norm constants, period maps, representatives and all-place RrT are constructed together; its one numerical record precedes every eta/V/S/N/G. RDR gives the original level-zero side, moved-Read and ReadPk/pin constructors on those objects, with one Side0 before every reader and all twelve Read ties. Full ReadL retains the exact o.eta(j,v)=q.eta(j,v) guard for every 1<=j<=m and place v; arbitrary Over does not imply it. Inherited CSR fixes one ReadC before all TwArgs, transports both metrics and all thirteen reading maps, and gives actual period/evaluation with kappa=(c*nP)\textasciicircum{}(-1/2). RPI supplies exact row-pair integrability, so first updated DivDef3 follows from the retained original Pins/Holds, nonzero chain, five real-group predicates and group tie. Initial theta extension remains complete. Remaining work is actual initial common geometry and its original predicates, Fits/primitive rational tests, eta agreement and outer family/eventual/same-Bundle duties, reFace compatibility, MultEq support, original chain/group data, actual base rows and any stronger DivRows3 theta square. No new compatible reading, family threshold or Haar normalization is substituted. Version 0.41 supplies the finite phase-correct rational theta route: exact coefficient extension, split representatives, square-law/PhiIs, Pins from CoversAt, class-sum/point-identity transfer to OneMapOverL, actual-good-projector RatG and both-level LiftTestB. Its original higher covers, all-tuple class-sum uniqueness including phase rescaling, one PointIdentCore before O, base ClassSumsSpan and geometric LatFull remain explicit producers. MLat is available from the original LatFull input. This sharpens the rational-test work without claiming actual common geometry, Fits or eta agreement.



\paragraph{23. BASE\_\allowbreak THETA\_\allowbreak RATIONAL\_\allowbreak TESTS — C — Simultaneous original base rational tests}

On the SAME original Bundle/reading/theta map, for every i<m supply one finite separating complex-linear test family whose values lie in the fixed embedded coefficient field on theta of EVERY original rational source vector and on EVERY stored target-coordinate basis vector. BTH8/BTH11 use this same LiftTestB for rational adjoint image and matrix coordinates. RTG now reconstructs prescribed-field coefficient extension, original split class representatives/RepsOK, the self-dual split-lattice square law and finite residue PhiIs; every Pins field follows from the actual CoversAt. With the base coefficient record literally fixed, exact class-sum existence/uniqueness including phase-rescaled O, PointIdentCore with common lambda,z before every O, and base ClassSumsSpan imply all-tuple finite lift span and OneMapOverL. The phase is exactly cj(twAt embLL A3.A.v0)(R.phs i). RatG uses the actual good projector eOp(y), and full LiftTestB follows at BOTH original levels. These finite implication proofs are complete. Remaining original producers are the prescribed higher covers, those actual class sums/uniqueness, PointIdentCore, base ClassSumsSpan and the base/common geometric datum; exact MLat is supplied only from original LatFull, whose geometric production remains separate. Finite dimension or an unrelated rational theta model cannot replace these same-object identities. Original outer eventual-family/Bundle quantifiers remain. Version0.46 proves finite generation of the original whole integral cusp lattice by actual integral-frame evaluation. It also proves the exact identity-cover, analytic unfolding, original-field coefficient descent, global clearing-scalar and both-level lattice-transport implications. Full spanning retains the supplied original IsNab connection and the five explicit principal-chart, level-system, raw-Gaussian, short-rule/dictionary and ramified-jet producers. The higher tuple covers, actual common class sums, PointIdentCore and ClassSumsSpan remain separate. CURRENT VERSION0.47 QUALIFICATION: Construct one original common geometric family and the higher CoversAt/all-tuple class sums and uniqueness with original phase rescalings, one PointIdentCore with common lambda,z, and base ClassSumsSpan. CNC now supplies the full original IsNab; PCL supplies the conditional trace-dual principal-lattice and quantitative normalization. Actual uniformization/HasPrinc and the five principal-chart, finite-level-system, Gaussian/reproduction, short-rule/simultaneous-dictionary and ramified-jet producer groups remain. Keep all original both-level, coefficient-field, Fits and eta agreements. Completed new original producers: ICL1-ICL21, CNC1-CNC16, PCL1-PCL12. Any earlier narrative naming these completed producers as outstanding is historical; this qualification is the current scope authority. CURRENT VERSION0.48 QUALIFICATION: Construct the original finite holomorphic FamilyUnif/N1c, normalized theta families and local lifting, then the common principal charts and finite-level theta systems. PUN now supplies the exact fixed-space models, all-family center lattices, A=1 exponent matching and individual primitive points with the full point-equality criterion; CNC supplies the full canonical connection and PCL the quantitative trace-dual normalization. Complete higher CoversAt and all-tuple class sums/uniqueness with the original phase rescalings, one PointIdentCore with common lambda,z, and base ClassSumsSpan. The Gaussian/reproduction, short-rule/simultaneous-dictionary and ramified-jet producers, and all original both-level, coefficient-field, Fits and eta agreements remain explicit. Completed current producers or conditional consumers: ICL1-ICL21, CNC1-CNC16, PCL1-PCL12, PUN1-PUN24. Earlier frontier descriptions are historical; this qualification is the current scope authority. Original whole-contract premises and quantifiers are retained. CURRENT VERSION0.49 QUALIFICATION: Construct the adjoining simultaneous principal framed points and test-ring dictionary. Full finite holomorphic FamilyUnif/N1c, normalized theta families and local lifts are now supplied by HUF. The next bounded target is allTest\_\allowbreak of\_\allowbreak secBaseChange: represent FrPt by the Lie-eigenframe bundles and theta-section frame line, use the full n=1 SecBaseChangeSt on every nonflat coefficient base change, and prove finite-generated subalgebra descent with the exact TRing/FromTest transport. Then supply the full FramesOK/CoverTri/FibTri and IsUnf/hasPrinc package, including adapted mod-eight symplectic basis, normalized theta multiplier/weight and arithmetic invariance. Continue higher CoversAt, all-tuple class sums and uniqueness with original phase rescalings, one PointIdentCore with common lambda,z, base ClassSumsSpan, Gaussian/reproduction, short-rule dictionaries and ramified jets, retaining both-level, coefficient-field, Fits and eta agreements. Completed current producers or conditional consumers: ICL1-ICL21, CNC1-CNC16, PCL1-PCL12, PUN1-PUN24, HUF1-HUF21 (full original FamilyUnif/N1c and theta family). Earlier frontier descriptions are historical; this qualification is the current scope authority. Original whole-contract premises and quantifiers are retained.

\emph{New completed inputs.} ICL1-ICL21, CNC1-CNC16, PCL1-PCL12, PUN1-PUN24, HUF1-HUF21 (full original FamilyUnif/N1c and theta family)

\emph{Current remaining subfrontier.} Construct the adjoining simultaneous principal framed points and test-ring dictionary. Full finite holomorphic FamilyUnif/N1c, normalized theta families and local lifts are now supplied by HUF. The next bounded target is allTest\_\allowbreak of\_\allowbreak secBaseChange: represent FrPt by the Lie-eigenframe bundles and theta-section frame line, use the full n=1 SecBaseChangeSt on every nonflat coefficient base change, and prove finite-generated subalgebra descent with the exact TRing/FromTest transport. Then supply the full FramesOK/CoverTri/FibTri and IsUnf/hasPrinc package, including adapted mod-eight symplectic basis, normalized theta multiplier/weight and arithmetic invariance. Continue higher CoversAt, all-tuple class sums and uniqueness with original phase rescalings, one PointIdentCore with common lambda,z, base ClassSumsSpan, Gaussian/reproduction, short-rule dictionaries and ramified jets, retaining both-level, coefficient-field, Fits and eta agreements.



\paragraph{24. SLR\_\allowbreak BASE\_\allowbreak ROWS\_\allowbreak PRODUCER — C — Actual base row identities on that same bundle}

For the SAME selected Bundle, common original fields and readings used by the other three BTH12 inputs, supply actual RowsHold0\_\allowbreak 3 uniformly eventually for every family. Its original primitive row realization, spectrum/type and basis identifications must be retained. BTH supplies a theta consumer only; it does not construct the base row identity. No separate witness may be chosen merely to satisfy this row requirement.



\clearpage
\section{From ordinary end data to every coefficient-algebra point}
\label{sec:unramified-coefficient-extension-v49}

This chapter reconstructs the original curve and specialization arguments
between ordinary end data and the joint unramified-coefficient automorphy
assertion.  The conclusions are the literal \texttt{Fin\allowbreak Crit},
\texttt{Const\allowbreak Of\allowbreak Char}, \texttt{ord\allowbreak Pt\_\allowbreak of\_\allowbreak gd},
\texttt{fin\allowbreak A\_\allowbreak of\_\allowbreak end}, \texttt{exists\_\allowbreak ord\allowbreak Pt\_\allowbreak of\_\allowbreak point},
\texttt{mat\allowbreak Irr\_\allowbreak open\_\allowbreak gd}, and the same-marked-element implication
\[
 \mathrm{GdEnd}\Rightarrow\mathrm{AutomorphyUnramifiedCoefficients}.
\]
The common ordinary end datum remains an explicit input to the last
implication.  In particular, the chapter does not assert the existence of
its finite-characteristic or automorphic end conditions.

Here are the objects and the exact conventions.  Fix the original
\(D:\mathrm{Dat}(\Omega,K,c,e,\psi)\), put
\(\mathcal O=c.O\), and write
\[
 G=G_{\mathbf Q},\qquad H=\ker\eta=G_K,\qquad
 \Psi=\epsilon^e\psi,\qquad (w_0,w_1,w_2)=(1,0,-1).
 \tag{UCE.1}
\]
All algebraic closures, embeddings and the later bridge
\(\iota:\overline{\mathbf Q}_\ell\simeq\mathbf C\) are the supplied
ones.  The group is
\(\mathcal G_3=\mathrm{GL}_3\rtimes\{1,-1\}\), with
\(\theta(M)=JM^{-t}J\), where \(J\) is the antidiagonal permutation
matrix.  Put \(T(t)=\operatorname{diag}(t,1,t^{-1})\) and
\(r_0(g)=(T(\Psi(g)),\eta(g))\).  Entries have weights
\(w_i-w_j\).  The fixed complex conjugation is denoted by \(\sigma\).
A marked element satisfies
\[
 \sigma^2=1,\quad \eta(\sigma)=-1,\quad \Psi(\sigma)=1,
 \quad\eta(g_1)=1,\quad \Psi(g_1)^2-1\in\mathcal O^\times.
 \tag{UCE.2}
\]
Normalization means that the matrix of \(r(\sigma)\) is one and the
\((0,1)\) and \((1,0)\) entries of
\(r(g_1)-r(\sigma g_1^{-1}\sigma)\) vanish.  We always use the actual
local ideal \(I_t\) and the actual normalized universal ring \(S\)
of OCH7--OCH9.

The original \(GD\) assumptions remain unchanged: \(\ell\ne2\),
\(e\ge1\), even nonzero \(N\), prime-to-\(\ell\) order of
\(\psi\), \((\operatorname{comp}_\ell\psi)^2=1\), its prescribed
parity, imaginary quadratic \(K\),
\(\ell\nmid |\operatorname{disc}K|h_K\),
\((|\operatorname{disc}K|,N)=1\), the original residual kernel
noncontainment, and \(\ell-1\nmid2e\).  No lower bound on the prime
is added.  The local and finite-characteristic lemmas include prime
three in their stated scopes; the last \(GD\) condition itself excludes
it in the final arithmetic application.

For a valuation ring \(B\subset\Omega\) over \(\ell\), an
\emph{ordinary point} over a complete local algebra \(A\) has the
five original fields: the prescribed residual representation and
component character; continuity modulo every power of \(\mathfrak m_A\);
the normalization just described; triviality on
\(I_{B'}\cap\ker\eta\cap\ker\Psi\) at every place away from
\(\ell\); and, at \(B\), a basis congruent to the identity in which
the representation on the kernel of \(\eta\) in the decomposition
group is upper triangular, with inertia diagonal
\((\Psi,1,\Psi^{-1})\).  This is exactly \texttt{Ord\allowbreak Pt}.

The already reconstructed inputs used below are OCH4--OCH9 and the
local part of OCH15, with their named period and cohomology foundations.
The additional background is elementary complete local algebra,
valuation rings, finite extensions of complete discretely valued
fields, and the usual Galois decomposition/inertia dictionary.  The
proofs below give the curve detection, finite generation, ordered
matrix arguments and extension of arbitrary algebra maps explicitly.

\subsection{Ordinary points over the original coefficient rings}

% BEGIN UCE1
\paragraph{UCE1. The actual inertia separator and marked elements.}
\label{claim:UCE1}
For every valuation ring \(B\) of \(\Omega\) above \(\ell\), there is
\(x\in I_B\) such that
\(\eta(x)=1\) and \(\Psi(x)^2-1\in\mathcal O^\times\).
In particular, taking \(g_1=x\) gives (UCE.2).  This conclusion alone
does not assign the end conditions to that choice.

\emph{Proof.}
The cyclic group \(\mathbf F_\ell^\times\) has an element \(a\)
with \(a^{2e}\ne1\), since \(\ell-1\nmid2e\).
Inertia surjects onto the cyclotomic character modulo \(\ell\):
\(\mathbf Q_\ell(\mu_\ell)/\mathbf Q_\ell\) is totally ramified.
Choose an inertia element with cyclotomic reduction \(a\).
The character \(\eta\) is unramified at \(\ell\), and the square
of the finite inertia part of \(\Psi\) is one by the original
\(\operatorname{comp}_\ell\psi\) condition.  Consequently
\(\overline\Psi(x)^2=a^{2e}\ne1\).  This is precisely the asserted
unit condition.  Changing the prolongation of the place conjugates
inertia; both \(\eta\) and \(\Psi\) are characters and retain their
values.  The other marked identities are the supplied complex
conjugation and parity identities.  This proves the assertion on the
given algebraic closure and for every original \(B\).
% END UCE1

% BEGIN UCE2
\paragraph{UCE2. Ordered flags specialize after a finite field extension.}
\label{claim:UCE2}
Let \(R\) be a domain, let matrices \(M_s\) and prescribed diagonal
entries \(d_{s,i}\) be defined over \(R\), and suppose that over a
field containing \(\operatorname{Frac}R\) they admit one ordered
upper triangular flag, with those diagonals for every index in the
specified inertia subset.  For every map \(R\to k\) to a field,
the specialized matrices admit the same ordered triangular and
diagonal conditions over a finite extension of \(k\).
No restriction is imposed on the characteristic of \(k\).

\emph{Proof.}
Let \(\mathfrak p\) be the kernel of the map.  Choose a valuation
ring in the generic field dominating \(R_{\mathfrak p}\).
To recall the existence argument, order the local subrings of the
field dominating \(R_{\mathfrak p}\) by domination.  A union of a
chain is again such a ring.  A maximal one is a valuation ring:
if neither \(z\) nor \(z^{-1}\) belongs to it, adjoining one and
localizing at a prime above its maximal ideal gives a larger
dominating local ring, a contradiction.  Denote it by \(V\).

Every flag in the fraction field of \(V\) has a basis matrix
\(k b\), where \(k\in\mathrm{GL}_n(V)\) and \(b\) is upper
triangular over the field.  Indeed, in the first column choose an
entry dividing all its entries; such an entry exists because the
finite set of principal ideals is totally ordered.  Scale the
column to a primitive vector, carry it to the first basis vector by
invertible row operations over \(V\), and continue in the quotient.
Induction gives the assertion.  Changing the generic flag basis by
\(b\) preserves upper triangularity and all ordered diagonal entries.
In the basis \(k\), the matrices are integral because the original
matrices lie in \(V\).  Their asserted entries therefore remain zero,
or equal to the prescribed diagonals, on reduction to \(\kappa(V)\).

Domination identifies the kernel of \(R\to\kappa(V)\) with
\(\mathfrak p\).  The two fields \(k\) and \(\kappa(V)\) have a
common extension over \(\operatorname{Frac}(R/\mathfrak p)\): take
a prime quotient of their nonzero tensor product and then its
fraction field.  This gives a flag over some extension of \(k\).
Finally consider the affine scheme of invertible basis matrices
with the required triangular and diagonal entries.  Its equations
are polynomial, after adding an inverse for the determinant.
Although their indices may be infinite, their ideal in the
finite-variable polynomial ring is finitely generated.  It is proper
because of the flag just constructed.  A maximal ideal containing it
has residue field finite over \(k\), by the finite-type field lemma.
That point supplies all the original equations simultaneously and
proves the required finite-extension assertion.
% END UCE2

% BEGIN UCE3
\paragraph{UCE3. The full local ordered-flag assertion in every characteristic.}
\label{claim:UCE3}
For every original local datum \(D_\ell:\mathrm{Dat3}\), its
universal ring \(U\), every field \(k\) over \(\mathcal O\),
and every actual representation obtained from a map
\(U/I_t\to k\), there is a finite extension \(k'/k\) and a
basis over \(k'\) giving the full original \texttt{Tri\allowbreak In} condition.
Thus the original \texttt{lemma2c} is supplied, including targets
of characteristic \(\ell\).

\emph{Proof.}
OCH7 gives the domain \(R=U/I_t\), its injective characteristic-zero
formal chart, and its base point.  In every finite truncation of that
chart, OCH5 makes the representation crystalline and OCH4 makes it
ordinary.  Its diagonal infinitesimal characters, after division by
the fixed characters, are unramified: the weight-zero crystalline
\(H^1_f\) calculation used in OCH4 and OCH15 has precisely that
conclusion.  Induction on the Artin length therefore keeps the
\emph{full} inertia diagonals equal to
\(\lambda_i(s)=\Psi(s)^{w_i}\).
Consequently the following polynomial identities hold in every
truncation and hence in the injective chart and in \(R\):
\[
 \det(X-r(s))=\prod_{i=0}^2(X-\lambda_i(s)),\qquad
 \prod_{i=0}^2\bigl(r(s_i)-\lambda_i(s_i)I\bigr)=0
       \quad(s,s_i\in I_\ell),
 \tag{UCE.3}
\]
where the second product has the displayed order.

Over \(\operatorname{Frac}R\), choose an inertia element \(u\)
in the kernel of the finite character part whose cyclotomic value
has infinite order.  Such an element exists because an open subgroup
of inertia still has infinite cyclotomic image.  Since \(e\ge1\),
the three values \(\lambda_i(u)\) are distinct in characteristic
zero.  The first identity in (UCE.3) diagonalizes \(r(u)\) in that
order.  In that basis apply the second identity successively to
\((s,u,u),(u,s,u),(u,u,s)\).  Division by the nonzero eigenvalue
differences shows that \(r(s)\) is upper triangular and its
\(i\)-th diagonal is \(\lambda_i(s)\).  For example, the first
product kills the entries below the first diagonal position and
identifies its diagonal entry; the last product does the same for
the last row, and the middle product identifies the middle diagonal.

If \(g\) is in the kernel of \(\eta\) in the decomposition
group, then \(gug^{-1}\) is in inertia and has the same ordered
character values.  The equation
\(r(gug^{-1})r(g)=r(g)r(u)\), solved from the bottom row upward,
forces \(r(g)\) to be upper triangular.  At each off-diagonal
step the coefficient is one of the same nonzero eigenvalue
differences.  We have obtained the full ordered generic flag,
with all decomposition and inertia equations, over
\(\operatorname{Frac}R\).  Apply UCE2 to the given map \(R\to k\).
The multiplier is unchanged throughout.  Its output is exactly the
original finite-extension \texttt{Tri\allowbreak In} assertion.
% END UCE3

% BEGIN UCE4
\paragraph{UCE4. Descent of the ordered flag to the actual integral basis.}
\label{claim:UCE4}
Let \(A\) be a complete local domain over \(\mathcal O\), and let
\(r\) be a residual lift of the actual local rank-three datum over
\(A\).  Suppose its fraction-field representation has the ordered
\texttt{Tri\allowbreak In} flag after a field extension.  If inertia contains
\(x\) with \(\Psi(x)^2-1\in\mathcal O^\times\), then \(r\)
satisfies the original integral \texttt{Cond\allowbreak T}: its basis can be
chosen over \(A\), congruent to the identity.

\emph{Proof.}
Put \(d_i=\Psi(x)^{w_i}\).  Every \(d_i-d_j\), \(i\ne j\),
is a unit: each differs by a unit factor from \(\Psi(x)-1\),
\(\Psi(x)+1\), or \(\Psi(x)^2-1\).
The generic ordered flag gives
\(\det(X-r(x))=\prod_i(X-d_i)\); injectivity into the field
puts the same identity in \(A[X]\).  Form the spectral projectors
\[
 P_i=\prod_{j\ne i}\frac{r(x)-d_jI}{d_i-d_j}\quad\in M_3(A).
 \tag{UCE.4}
\]
Cayley--Hamilton makes them orthogonal idempotents of sum one.
Their reductions are the three coordinate projectors.  Thus
\(P_i e_i\) has \(i\)-th coordinate a unit; divide by that unit.
These three vectors form a matrix \(h\equiv I\pmod{\mathfrak m_A}\)
and \(h^{-1}r(x)h=\operatorname{diag}(d_i)\).

Let \(k\) be a generic ordered-flag matrix.  The equation
\[
       (k^{-1}r(x)k)(k^{-1}h)=(k^{-1}h)\operatorname{diag}(d_i)
\]
forces \(k^{-1}h\) to be upper triangular, by solving its entries
from the bottom row upward.  The different \(d_i\) have already
been matched in their prescribed order.  Conjugation by an upper
triangular matrix preserves upper triangularity and its diagonal
characters.  The required equations therefore hold in the \(h\)
basis over the generic field, and then over \(A\) by injectivity.
This proves every decomposition-group and inertia clause of
\texttt{Cond\allowbreak T}, with the actual integral near-identity matrix.
% END UCE4

% BEGIN UCE5
\paragraph{UCE5. Every original local domain specialization is an ordinary point.}
\label{claim:UCE5}
Let \(D,g_1,U,S\) be original data with \(g_1\) marked, and let
\(B\subset\Omega\) be the actual valuation ring defined by
\(x\in B\Longleftrightarrow\|D.j(x)\|\le1\).
For every complete local domain \(A\) over \(\mathcal O\) and
every local algebra map \(f:S\to A\), the actual push
\(f_*r_S\) is an \(\mathrm{OrdPt}\) for the same
\(D,g_1,B\).  This is the full original \texttt{ord\allowbreak Pt\_\allowbreak of\_\allowbreak gd}.

\emph{Proof.}
A local map sends \(\mathfrak m_S^a\) into
\(\mathfrak m_A^a\) for every \(a\).  Hence the residual matrix
identities and all locally constant reductions of \(r_S\) push
to the exact residual and adic-continuity fields.  The normalization
equations also push literally.
At a place away from \(\ell\), an element with
\(\eta(s)=\Psi(s)=1\) belongs to the original killed-inertia
set, so its universal image, and therefore its image under \(f\),
is one.  This proves the original all-place away condition.

The specified local map has image the decomposition group of \(B\)
and takes its local inertia onto \(I_B\); this is the fixed
embedding dictionary in OCH9.  The local factorization of \(S\)
kills the same \(I_t\).  Apply UCE3 over \(\operatorname{Frac}A\),
then UCE1 and UCE4.  The resulting integral basis gives exactly
the last field of \(\mathrm{OrdPt}\).  This works even when
\(\mathcal O\to A\) kills its maximal ideal, and uses neither
\(\mathrm{FinA}\) nor an automorphy hypothesis.
% END UCE5

\subsection{The actual weighted ring and its characteristic curves}

% BEGIN UCE6
\paragraph{UCE6. The integral weighted presentation by actual entries.}
\label{claim:UCE6}
For every original normalized ring \(S\), there is a finite
weighted presentation
\[
 P=\mathcal O[[X_1,\ldots,X_d]]\twoheadrightarrow S,
 \qquad X_i\longmapsto z_i,
 \tag{UCE.5}
\]
where each \(z_i\in\mathfrak m_S\) is an actual matrix-entry
deviation from \(r_0\), its weight \(a_i\) is the corresponding
\(w_b-w_c\), and the kernel is homogeneous under every formal
weight projection.  Every universal matrix entry belongs to its
asserted homogeneous piece.  These are the literal \texttt{WPres}
properties needed in both original coefficient arguments.

\emph{Proof.}
The entry deviations span the relative cotangent space.  Otherwise
a nonzero functional on that space vanishing on all the deviations
would give two different maps to the residue dual-number algebra
with identical pushed representations, contradicting the uniqueness
in the actual universal property.  Choose finitely many deviations
spanning that finite-dimensional space.  Successive approximation
and completeness, beginning with Nakayama's lemma, give the
surjection (UCE.5).

Here is a proof of integral homogeneity that does not divide by
differences of torus eigenvalues.  Over \(S[[T]]\), conjugate the
universal representation by \(T(1+T)\).  This preserves the
residual representation, killed inertia and normalization; OCH8
preserves its literal local \(I_t\)-condition.  Universality gives
a map \(S\to S[[T]]\) that multiplies every entry of weight \(j\)
by \((1+T)^j\), and specializes at \(T=0\) to the identity.
For \(F\in\ker(P\to S)\), it follows that
\(F((1+T)^{a_i}z_i)=0\).
Modulo \(\mathfrak m_S^b\), only finitely many monomials in the
\(z_i\) survive.  Replace \(1+T\) by a Laurent indeterminate
\(Z\).  The substitution
\[
 (S/\mathfrak m_S^b)[Z,Z^{-1}]
       \longrightarrow(S/\mathfrak m_S^b)[[T]],\qquad Z\mapsto1+T
\]
is injective: after multiplication by a power of \(Z\), this is
the invertible integral change of variable in a polynomial.
Each coefficient of \(Z^j\) is therefore zero.  Letting \(b\)
increase shows that every weight projection \(F_j\) is in the
original kernel.  The same argument applied to a lift of any
universal entry shows that all its other weight projections vanish.
Thus the pieces \(S_j\), defined as images of homogeneous formal
series, have independent finite sums, multiply in the prescribed
weights, and contain the actual entries with their original weights.
The chosen deviations also lie in the span of entry monomials of
their weight, because their base constants vanish off weight zero.
% END UCE6

% BEGIN UCE7
\paragraph{UCE7. Weight zero and every fixed weight have finite monomial presentations.}
\label{claim:UCE7}
For any finite weighted presentation as in (UCE.5), with an integral
homogeneous kernel and variables in the maximal ideal, \(S_0\) is
a complete Noetherian local \(\mathcal O\)-algebra with the same
residue field.  Each \(S_j\) is a finite \(S_0\)-module.  The
maximal ideal of \(S_0\) is generated by \(\mathfrak m_{\mathcal O}\)
and finitely many positive-degree weight-zero monomials.  Moreover
\[
               S_0\cap\pi S=\pi S_0
 \tag{UCE.6}
\]
for a uniformizer \(\pi\) of \(\mathcal O\).

\emph{Proof.}
The nonzero minimal exponent vectors of weight zero under
componentwise comparison are finite by the ascending-chain property
of monomial ideals in a finite-variable polynomial ring.  Every
nonzero weight-zero exponent vector dominates one of them; subtract
it and repeat.  Thus they generate the weight-zero monoid.  For
a fixed weight \(j\), the minimal vectors of that weight are
again finite, and every vector of weight \(j\) is a minimal one
plus a weight-zero vector.  Grouping formal monomials by these
choices shows that the weight-zero power-series ring is a quotient
of a finite-variable power-series ring over \(\mathcal O\), and
that each fixed-weight series module is finite over it.  The
grouping converges: the lengths of the chosen factorizations tend
to infinity with total degree, since the finite generator degrees
are bounded.

The homogeneous kernel in (UCE.5) gives the same statements after
quotienting.  The kernels in these finite-module presentations are
closed, so their images are complete.  A weight-zero element is
a constant from \(\mathcal O\) plus an element generated by the
positive-degree weight-zero monomials.  If the residue of the
constant is nonzero, its inverse is the convergent geometric series,
still of weight zero.  This identifies the local ring, its residue
field and its asserted maximal ideal.

Finally, if \(s\in S_0\) and \(s=\pi y\) in \(S\), choose
a homogeneous weight-zero lift \(F\) of \(s\) and an arbitrary
lift \(Y\) of \(y\).  Project \(F-\pi Y\), which lies in
the homogeneous kernel, to weight zero.  It gives
\(s=\pi\overline{Y_0}\) in \(S_0\).  This proves (UCE.6),
including rings with \(\pi\)-torsion.
% END UCE7

% BEGIN UCE8
\paragraph{UCE8. Detection of a surviving element by a characteristic-prime DVR curve.}
\label{claim:UCE8}
Let \(R\) be a complete Noetherian local ring with finite residue
field of characteristic \(p\).  Let \(Q\) be a prime containing
\(p\), and let \(b\in\mathfrak m_R\setminus Q\).  There is
a local map \(R\to A\) to a complete discrete valuation ring of
characteristic \(p\), with finite residue field, whose kernel
contains \(Q\) and does not contain \(b\).

\emph{Proof.}
Choose a prime \(\mathfrak p\supset Q\), maximal among primes
not containing \(b\), using the ascending-chain condition.
Then \(D=R/\mathfrak p\) is a complete local domain and the
nonzero image of \(b\) lies in its maximal ideal.  Its dimension
is one.  Here is the dimension argument.  Every nonzero prime of
\(D\) contains \(b\), by maximality.  For a nonzero
\(x\in\mathfrak m_D\), a prime minimal over \((x)\) has
height one by the principal ideal theorem.  It contains \(b\),
and hence is also minimal over \((b)\), since a strict intermediate
nonzero prime would give height at least two.  Consequently
\(\mathfrak m_D\) is contained in the union of the finitely many
primes minimal over \((b)\).  Finite prime avoidance puts it in
one of them.  It has height one, as asserted.

For completeness, the required discrete valuation ring can be
constructed directly in this equal-characteristic situation.
Let \(k_D\) be the finite residue field of \(D\), of order \(q\).
The roots of \(X^q-X\) lift uniquely by completeness and its unit
derivative \(-1\).  In characteristic \(p\) they are closed
under addition, multiplication and inverse, and give a coefficient
field \(k_D\subset D\).  Choose any nonzero
\(t\in\mathfrak m_D\).  Dimension one makes \((t)\)
\(\mathfrak m_D\)-primary, so the \(t\)-adic and maximal-ideal
topologies are equivalent.  The map \(k_D[[T]]\to D\),
\(T\mapsto t\), is injective: a nonzero series is a power of
\(T\) times a unit, and neither factor becomes zero.  Lifting a
finite basis of \(D/tD\) and successively subtracting its constant
terms proves that \(D\) is finite over \(k_D[[t]]\).

Thus \(\operatorname{Frac}D\) is a finite extension of
\(k_D((t))\).  The extension of its complete discrete valuation
has a complete DVR of integers \(A\), finite over \(k_D[[t]]\)
with finite residue field.  Every element of \(D\) is integral
over \(k_D[[t]]\), so belongs to \(A\).  The inclusion is local:
a power of \(\mathfrak m_D\) lies in \(tD\), whereas units
stay units.  Its injectivity keeps \(b\) nonzero.  Composing the
quotient map with this inclusion proves the assertion.  The same
argument applies when the specified prime contains a uniformizer
of a coefficient DVR: the quotient then has the residue
characteristic, which is the case used below.
% END UCE8

% BEGIN UCE9
\paragraph{UCE9. The complete original finite weight-zero criterion.}
\label{claim:UCE9}
Let \(\mathcal O\) be any complete DVR with finite residue field,
let \(R\) be an original complete Noetherian local
\(\mathcal O\)-algebra with the same residue field, and let
\(R\) have a finite integral weighted presentation with all
variables in its maximal ideal.  Suppose that for every local map
\(f:R\to A\) to a complete DVR with finite residue field, for
which \(f(\mathfrak m_{\mathcal O})=0\), all weight-zero
variables vanish and either all negative-weight variables or all
positive-weight variables vanish.  Then \(R_0\) is finite over
\(\mathcal O\).  This is the literal \texttt{Fin\allowbreak Crit}, with all
its original test rings and arbitrary integer variable weights.

\emph{Proof.}
Take one of UCE7's finitely many positive-degree weight-zero
monomial generators \(b\).  If its image in \(R\) were not
nilpotent modulo \(\pi R\), a prime \(Q\supset\pi R\)
would avoid it.  UCE8 supplies an original test map with
\(f(b)\ne0\).  In a nonempty weight-zero monomial, either a
weight-zero variable occurs, or both a negative and a positive
weight occur.  The assumed alternative therefore kills at least
one of its factors.  This contradiction proves
\(b^m\in\pi R\) for some \(m\).
Equation (UCE.6) puts \(b^m\) in \(\pi R_0\).

All generators of the maximal ideal of \(R_0\) are consequently
in \(\sqrt{\pi R_0}\).  Since there are finitely many, some
power of that maximal ideal lies in \(\pi R_0\).  The quotient
\(R_0/\pi R_0\) is therefore Artinian with finite residue field,
and is finite-dimensional over \(\mathcal O/\pi\).
The two adic topologies on \(R_0\) are equivalent.  Lift a finite
set of generators of this quotient.  At each successive
\(\pi\)-power subtract their \(\mathcal O\)-linear combination;
completeness sums the resulting coefficients in \(\mathcal O\).
The lifts generate \(R_0\) as an \(\mathcal O\)-module.
This proves finiteness, without assuming that \(R\) is reduced,
flat over \(\mathcal O\), or characteristic zero.
% END UCE9

% BEGIN UCE10
\paragraph{UCE10. Finite-index characteristic polynomials determine the constants everywhere.}
\label{claim:UCE10}
Let \(k\) be a finite field, let \(A\) be a local domain over
\(k\), and let a representation over \(A\) reduce to the supplied
diagonal character representation over \(k\).  If their
characteristic polynomials agree on a finite-index subgroup, then
they agree at every group element.  For a disconnected rank-three
representation the assertion applies on \(H=\ker\eta\), exactly
as in the original \texttt{Const\allowbreak Of\allowbreak Char} contract.

\emph{Proof.}
For each \(g\), some positive power \(g^a\) lies in the
given finite-index subgroup.  Every eigenvalue of \(r(g)^a\)
is algebraic over \(k\), so every eigenvalue of \(r(g)\) is
algebraic over \(k\).  Hence all coefficients of its characteristic
polynomial are algebraic over \(k\).
An element \(y\in A\) algebraic over \(k\) and congruent to
\(c\in k\) is equal to \(c\): the domain \(k[y]\) is a
finite-dimensional field over \(k\); if \(y-c\ne0\), its
inverse belongs to \(k[y]\subset A\), contrary to membership
in the maximal ideal.  Applying this argument to every coefficient
gives the asserted equality with the residual polynomial.
On \(H\), the matrix of \(r(g^a)\) is the ordinary power
of the matrix of \(r(g)\), so the same proof applies literally.
It uses neither continuity nor a prime-to-characteristic index.
% END UCE10

% BEGIN UCE11
\paragraph{UCE11. Constant characteristic polynomials force the original triangular alternative.}
\label{claim:UCE11}
The full original \texttt{Const\allowbreak Of\allowbreak Char} holds.  Namely, over a local
domain \(A\) over a finite field \(k\) of characteristic different
from two, take the supplied characters \(\eta,\Psi\), marked
\(\sigma,g_1\), normalized residual lift \(r\), and any
finite-index subgroup on which the characteristic polynomials are
those of \(r_0\).  Then all diagonal entries of \(r\) are the
given base entries, and \(r\) is either upper triangular on \(H\)
or lower triangular on \(H\).  No characteristic-three exception
is imposed.

\emph{Proof.}
UCE10 first gives the base characteristic polynomial at every
element of \(H\).  Normalization and the disconnected transpose
identity give, over \(A\),
\[
 r(g_1)-r(\sigma g_1^{-1}\sigma)
       =\zeta\operatorname{diag}(1,0,-1).
 \tag{UCE.7}
\]
Indeed the difference is minus its antidiagonal transpose, and its
two prescribed adjacent entries are zero.  The entries fixed by
that transpose are zero because two is invertible.  The residual
value of \(\zeta\) is \(\Psi(g_1)-\Psi(g_1)^{-1}\ne0\),
so \(\zeta\) is a unit.

Work over \(F=\operatorname{Frac}A\).  Every invariant subspace
is stable under (UCE.7), and hence is a coordinate subspace by
its three spectral projectors.  A zero \((i,j)\) entry forced
by such a subspace also forces the zero \((2-j,2-i)\) entry,
because
\(r(\sigma g^{-1}\sigma)=Jr(g)^tJ\) on \(H\).
The six proper nonempty coordinate subsets therefore give exactly
the following alternatives: the whole representation is irreducible;
it is upper triangular; it is lower triangular; or all four
odd-weight entries vanish and its outer two-dimensional block is
irreducible.  If that outer block were reducible, its two distinct
marked projectors would give a coordinate line and one of the
triangular alternatives.

We rule out both irreducible alternatives by an explicit trace
Gram argument.  For an irreducible block of size \(n=3\), or
the outer block of size \(n=2\), its marked diagonal has distinct
eigenvalues.  The coordinate matrix projectors belong to the
\(F\)-algebra spanned by its group matrices.  The directed graph
whose edges are nonzero matrix entries is strongly connected:
the coordinate span reachable from any vertex would otherwise be
a proper invariant subspace.  Multiplying a matrix on both sides
by coordinate projectors gives its matrix unit times a nonzero
scalar.  Multiplying along paths therefore supplies every matrix
unit.  The span of the group matrices is thus all \(M_n(F)\).

Choose \(n^2\) group matrices \(M_i\) forming an \(F\)-basis.
The matrix
\[
                     G_{ij}=\operatorname{tr}(M_iM_j)
 \tag{UCE.8}
\]
is nonsingular because the trace pairing on \(M_n(F)\) is
nondegenerate, in every characteristic.  For the full block its
entries lie in \(k\), by UCE10.  For the outer block, first note
that the middle entry is an eigenvalue of the full block matrix;
it is algebraic over \(k\) and congruent to one, so is one by
UCE10's elementary field argument.  The outer trace is then the
full trace minus one, and again (UCE.8) has entries in \(k\).
Its nonzero determinant is a nonzero constant, and hence remains
nonzero in the residue field of \(A\).  But every reduced \(M_i\)
is diagonal, so the reduced matrix (UCE.8) factors through the
\(n\)-dimensional space of diagonal matrices.  Its rank is at
most \(n<n^2\), a contradiction.  This argument does not divide
by \(n\) or \(n!\).

We have proved the triangular alternative.  Each diagonal entry
is now a root of the constant characteristic polynomial; it is
algebraic over \(k\) and has its prescribed residual value.
UCE10's field argument identifies it with that value.  Finally,
for \(g\notin H\), the element \(g\sigma\) belongs to
\(H\) and has the same matrix as \(g\), because the matrix
of \(r(\sigma)\) is one.  The base matrices also agree there.
This proves all the original diagonal assertions on the full group.
% END UCE11

% BEGIN UCE12
\paragraph{UCE12. The literal finite-end condition gives finite weight zero.}
\label{claim:UCE12}
Fix the original \(D,g_1,B\), with \(g_1\) marked and \(B\)
the valuation ring supplied by the local embedding.  Suppose the
original \(\mathrm{FinEnd}\) holds: there is one finite-index
subgroup \(V\subset G\) such that every \(\mathrm{OrdPt}\)
over every complete DVR \(A\) with finite residue field and
\(\mathfrak m_{\mathcal O}A=0\) has the base characteristic
polynomials at all \(g\in V\cap H\).
Then \(\mathrm{FinA}(D,g_1)\) holds for every original local
universal ring and every original normalized global ring.

\emph{Proof.}
Take any such \(U,S\), and use UCE6's weighted presentation by
actual entry deviations.  Every local map from \(S\) to a test
DVR in the hypothesis is an actual ordinary point by UCE5.
The single given subgroup \(V\) therefore supplies the
characteristic-polynomial hypothesis of UCE11.  Since the map
kills \(\mathfrak m_{\mathcal O}\), its coefficient action
factors through the finite residue field, exactly as required there.
UCE11 kills all diagonal deviations and either all positive- or
all negative-weight entry deviations.  Thus every one of the
test maps in UCE9 satisfies its stated alternative.  UCE9 proves
that \(S_0\) is finite over \(\mathcal O\).

Every product of universal entries of total weight zero lies in
\(S_0\).  The original weight-zero entry subalgebra, being an
\(\mathcal O\)-submodule of this finite module, is finite because
\(\mathcal O\) is Noetherian.  This is precisely the defining
conclusion of \(\mathrm{FinA}\).  Neither the test subgroup nor
the marked element was changed or made dependent on \(U,S\).
This proves the full original \texttt{fin\allowbreak A\_\allowbreak of\_\allowbreak end} implication.
% END UCE12

\subsection{Arbitrary entry-algebra points and finite coefficient realizations}

% BEGIN UCE13
\paragraph{UCE13. Under finite weight zero, the actual entry algebra is the graded algebra.}
\label{claim:UCE13}
Assume \(\mathrm{FinA}(D,g_1)\), and fix any original \(U,S\).
Let \(E_S\) be the actual entry subalgebra of \(S\), let
\(B_0\) be its original weight-zero entry subalgebra, and let
\(S_j\) be UCE6's homogeneous pieces.  Then
\[
                 B_0=S_0,\qquad E_S=\bigoplus_{j\in\mathbf Z}S_j.
 \tag{UCE.9}
\]
Each \(S_j\) is finite over \(\mathcal O\).  Its finite-module
topology agrees with its topology as a subspace of \(S\).

\emph{Proof.}
The inclusion \(B_0\subset S_0\) follows from the entry weights.
Every homogeneous weight-zero formal series is approximated, to
arbitrary maximal-ideal order, by polynomials in the entry deviations
of weight zero.  Each such polynomial belongs to \(B_0\).
By \(\mathrm{FinA}\), \(B_0\) is a finite \(\mathcal O\)-module.
It is the continuous image of a finite power of the compact ring
\(\mathcal O\) in the Hausdorff ring \(S\), and is therefore
closed.  This proves \(S_0\subset B_0\).

UCE7 expresses each \(S_j\) as a finite sum of monomials in the
chosen deviations with coefficients in \(S_0\).  Those monomials
belong to \(E_S\), so \(\bigoplus_jS_j\subset E_S\).
The reverse inclusion holds because each generating universal
entry is homogeneous.  This also shows that
\(E_S=S_0[z_1,\ldots,z_d]\), with the actual deviations.
Finiteness of each \(S_j\) follows from UCE7 and the established
finiteness of \(S_0\).  A finite free \(\mathcal O\)-module
surjects onto \(S_j\); its quotient is compact and its injection
into \(S\) is continuous and Hausdorff.  The resulting continuous
bijection is a homeomorphism.  This proves the final assertion,
including all torsion and nonreduced cases.
% END UCE13

% BEGIN UCE14
\paragraph{UCE14. Every entry-algebra point has a strictly small torus representative.}
\label{claim:UCE14}
Under the hypotheses of UCE13, let
\(\xi:E_S\to\overline{\mathbf Q}_\ell\) be any
\(\mathcal O\)-algebra map.  There is a single
\(t\in\overline{\mathbf Q}_\ell^\times\) such that the graded
map
\[
                  \xi_t(s)=t^j\xi(s)\quad(s\in S_j)
 \tag{UCE.10}
\]
takes every chosen deviation \(z_i\) to an element of strictly
positive valuation.  No continuity or localness of \(\xi\)
is assumed.

\emph{Proof.}
The finite local algebra \(S_0\) is integral over \(\mathcal O\).
Hence all its images under \(\xi\) are integral.  Moreover every
element of \(\mathfrak m_{S_0}\) has a power in \(\pi S_0\),
because the finite local algebra \(S_0/\pi S_0\) is Artinian.
Such images have strictly positive valuation.  In particular every
positive-degree weight-zero monomial in the \(z_i\) has strictly
positive valuation unless its image is zero.

For indices with \(\xi(z_i)\ne0\), write
\(b_i=v_\ell(\xi(z_i))\in\mathbf Q\).  We seek a rational
\(a\) with \(b_i+a_i a>0\) for all these indices.  Weight-zero
indices already satisfy the inequality.  If \(a_i>0>a_j\), the
balanced monomial
\(z_i^{-a_j/d}z_j^{a_i/d}\), where \(d=\gcd(a_i,-a_j)\),
has weight zero and positive degree.  Its strictly positive
valuation gives
\[
                       -b_i/a_i<-b_j/a_j.
 \tag{UCE.11}
\]
Thus the largest of the finitely many lower bounds is strictly
smaller than the smallest upper bound.  If one set is empty there
is only a one-sided bound.  Choose rational \(a\) in the resulting
open interval and choose \(t\) of valuation \(a\) in the given
algebraic closure.  Images equal to zero cause no constraint.
The direct grading in (UCE.9) makes (UCE.10) an actual algebra
map, and it has the required values.  This proves all cases,
including one-sided and zero images.
% END UCE14

% BEGIN UCE15
\paragraph{UCE15. The rescaled point extends continuously into one finite coefficient DVR.}
\label{claim:UCE15}
The map \(\xi_t\) of UCE14 extends to a local algebra map
\(f:S\to A\), where \(A\) is a subalgebra of the supplied
\(\overline{\mathbf Q}_\ell\), a complete DVR finite over
\(\mathcal O\).  Its inclusion into the algebraic closure agrees
with \(\xi_t\) on \emph{every} element of \(E_S\).

\emph{Proof.}
Choose a finite extension \(L/\operatorname{Frac}\mathcal O\)
inside the supplied algebraic closure containing \(t\), the
finitely many \(\xi(z_i)\), and the images of a finite set of
\(\mathcal O\)-module generators of \(S_0\).  By UCE13 all
of \(\xi_t(E_S)\) lies in \(L\).  Its generators over
\(\xi(S_0)\) have positive valuation, and \(\xi(S_0)\)
is integral, so this image lies in the ring of integers \(A\)
of \(L\).  This is a complete DVR finite over \(\mathcal O\)
and contains its original image.

Put \(b_i=\xi_t(z_i)\in\mathfrak m_A\).  Substitution
\(X_i\mapsto b_i\) gives convergent evaluation on the full
power-series ring \(P\) in (UCE.5).  To see that it factors
through \(S\), first take a homogeneous series \(F_j\).
Its truncations converge to its image in \(S_j\).  By UCE13
this is convergence in a finite \(\mathcal O\)-module, and
the \(\mathcal O\)-linear map \(\xi_t|_{S_j}\) is continuous.
Evaluation of \(F_j\) therefore agrees with that map, not only
with its values on polynomial truncations.
If \(F\in\ker(P\to S)\), every homogeneous projection
\(F_j\) is in the kernel by UCE6, so evaluates to zero.
The evaluation of \(F\) is the sum of these evaluations:
the variables all have norms bounded by one fixed \(q<1\),
so terms of total degree at least \(m\) have norm at most
\(q^m\), independently of their weight.  Grouping the convergent
sum by weights is thus legitimate.  Hence \(F\) evaluates to
zero and evaluation gives an actual map \(f:S\to A\).

The maximal ideal of \(S\) is generated by \(\pi,z_1,\ldots,z_d\),
and all their images lie in \(\mathfrak m_A\).  Thus \(f\)
is local and continuous.  The homogeneous argument above gives
the asserted equality on every \(S_j\), and (UCE.9) gives it
on the entire original entry algebra.  This establishes finite
coefficient realization without assigning an unjustified
continuity assumption to the initial map \(\xi\).
% END UCE15

% BEGIN UCE16
\paragraph{UCE16. The full original ordinary representative of every entry point.}
\label{claim:UCE16}
Assume \(\mathrm{FinA}(D,g_1)\).  For every original \(U,S\),
every map \(\xi:E_S\to\overline{\mathbf Q}_\ell\), and every
actual \(r:G\to\mathcal G_3(\overline{\mathbf Q}_\ell)\)
with \(\mathrm{IsPush}(\xi,r_S,r)\), there are a finite
coefficient DVR \(A\) inside the same algebraic closure, a unit
\(t\in\overline{\mathbf Q}_\ell^\times\), and an actual
ordinary point \(r':G\to\mathcal G_3(A)\) such that
\[
                    r'_\mathrm{incl}=T(t)\,r\,T(t)^{-1}
                        \quad\hbox{in }\mathcal G_3.
 \tag{UCE.12}
\]
This is the full literal \texttt{exists\_\allowbreak ord\allowbreak Pt\_\allowbreak of\_\allowbreak point}.

\emph{Proof.}
Use UCE14 and UCE15, and set \(r'=f_*r_S\).  UCE5 proves
all its ordinary-point fields with the same marked \(g_1\)
and the same valuation ring \(B\).  Every entry of weight
\(w_i-w_j\) is multiplied by \(t^{w_i-w_j}\) under
\(\xi_t\).  These are exactly the entries of the conjugate
in (UCE.12).  Its component character is unchanged by
\(\mathrm{IsPush}\).  Finally \(\theta(T(t))=T(t)\), so
the same entry formula holds on the disconnected component as
well as on \(H\).  Thus (UCE.12) is equality of the actual
full group homomorphisms, not merely an isomorphism of their
semisimplifications or a trace comparison.
% END UCE16

\subsection{The exact automorphy consumers on the common marked datum}

% BEGIN UCE17
\paragraph{UCE17. Every irreducible ordinary point stays irreducible on every open subgroup.}
\label{claim:UCE17}
Let \(A\subset\overline{\mathbf Q}_\ell\) be an original finite
coefficient DVR and let \(r'\) be an \(\mathrm{OrdPt}\) for
\(D,g_1,B\).  If its restriction to \(H\) is irreducible,
then its restriction to every open subgroup \(V\subset H\)
is irreducible.  This is precisely \texttt{mat\allowbreak Irr\_\allowbreak open\_\allowbreak gd}.

\emph{Proof.}
Replace \(V\) by its normal core \(V_0\) in \(G\); it
has finite index since \(G\) is profinite.  Choose an inertia
element with cyclotomic value \(1+\ell\), using total ramification
of every \(\ell\)-power cyclotomic extension and compactness.
Raise it to a positive multiple of \([G:V_0]\) and of the
order of the finite character part.  The resulting
\(u\in I_B\cap V_0\) has \(\Psi(u)\) of infinite order.
The three inertia eigenvalues \(\Psi(u),1,\Psi(u)^{-1}\)
are therefore distinct.

For clarity, the required finite-index Clifford argument is
elementary.  An irreducible \(H\)-module restricted to the normal
subgroup \(V_0\) is a sum of simple modules: take a minimal
nonzero \(V_0\)-submodule and sum its finitely many \(H\)-translates.
The sum is the whole module by irreducibility, and a sum of simple
submodules is semisimple.  Its isotypic components are permuted
transitively by \(H\).  Since the dimension is the prime three,
a reducible restriction is either scalar, or the sum of three
distinct character lines permuted transitively.  The distinct
eigenvalues of \(u\) rule out the scalar case.

Diagonalize \(r'(u)\) inside the ordered upper triangular
ordinary basis.  The diagonalizing change can be upper triangular,
by solving its three eigenvector equations successively.  The
three lines just found are its three eigenlines.  Every element
of \(H\) acts by a monomial matrix in this basis.
Take UCE1's \(x\in I_B\) with
\(\overline\Psi(x)^2\ne1\).  Its matrix is also upper
triangular, hence its permutation of the three lines is the
identity.  A transitive permutation group on three letters contains
a three-cycle.  Choose \(g\in H\) inducing one.  Thus
\[
                        \operatorname{tr}r'(gx^m)=0
                                      \quad(m\ge0).
 \tag{UCE.13}
\]
Reduce this equality in the original coefficient DVR.  With
\(a=\overline\Psi(g)\) and \(b=\overline\Psi(x)\), it becomes
\(ab^m+1+(ab^m)^{-1}=0\), or
\(a^2b^{2m}+ab^m+1=0\).  The equations for \(m=0,1\)
give \(a(b+1)+1=0\), since \(a(b-1)\ne0\).
Those for \(m=0,2\) give \(a(b^2+1)+1=0\), since
\(a(b^2-1)\ne0\).  Subtraction gives
\(ab(b-1)=0\), a contradiction.  Therefore the restriction
to \(V_0\), and hence to \(V\), is irreducible.
No additional small-prime restriction or cyclotomic-kernel
replacement was used.
% END UCE17

% BEGIN UCE18
\paragraph{UCE18. Both automorphic conclusions return along the actual torus conjugacy.}
\label{claim:UCE18}
Suppose (UCE.12) holds.  Type I is preserved, and an
\(\iota\)-automorphic conjugate-self-dual conclusion for
\(\rho_K(r')\epsilon^{-1}\) gives the same conclusion for
\(\rho_K(r)\epsilon^{-1}\).  Type II and its middle-entry-one
condition are preserved as well.  For every original extension
\(W\) of the outer corner of \(r\), set
\[
 D(t)=\operatorname{diag}(t,t^{-1}),\qquad
                         W'=D(t)WD(t)^{-1}.
 \tag{UCE.14}
\]
Then \(W'\) is the corresponding original extension of the
outer corner of \(r'\).  If it is attached, after the required
\(\epsilon^{-e}\) twist, to a newform of weight \(2e+1\),
the same newform and the same \(\iota\) attach to \(W\).

\emph{Proof.}
On \(H\), (UCE.12) is ordinary matrix conjugacy by \(T(t)\).
It identifies invariant subspaces and the actual representations
\(\rho_K\), and commutes with their scalar cyclotomic twists.
The defining automorphic datum is transported without change:
keep its cuspidal representation, regular algebraic infinity type
and conjugate-self-duality; transport its invariant complements
to preserve semisimplicity; and use invariance of characteristic
polynomials to preserve every member of the specified cofinite
Frobenius attachment set.  These are all the original
\(\mathrm{IsAutCsd}\) fields.

For Type II, all four odd-weight zero entries remain zero, and
the outer block is conjugated by \(D(t)\).  The middle entry
has weight zero and remains one.  The outer-corner equality in
the definition of \(\mathrm{IsExt}\) is thus satisfied by
\(W'\) at every element of \(H\), and its trace at the
specified complex conjugation is still zero by invariance of trace.
Conjugacy by \(D(t)\) is also an isomorphism of the actual full
two-dimensional group homomorphisms.  It commutes with
\(\epsilon^{-e}\).  Transporting invariant complements,
trivial inertia and every Frobenius characteristic polynomial
in the newform attachment preserves that attachment at the same
\(\iota\).  This argument applies to every supplied \(W\);
it does not choose a preferred extension or omit its other possible
quadratic twist.
% END UCE18

% BEGIN UCE19
\paragraph{UCE19. The original joint package from one common ordinary end datum.}
\label{claim:UCE19}
Suppose the original \(\mathrm{GdEnd}\) is supplied.  Then the
original \(\mathrm{AutomorphyUnramifiedCoefficients}\) holds.
More precisely, for a fixed original \(D\), fixed \(\iota\),
and actual local valuation ring \(B\), a single marked \(g_1\)
with the three original fields
\[
        \mathrm{FinEnd},\qquad
        \mathrm{AutEnd3}(K,\iota),\qquad
        \mathrm{AutEnd2}(e,\iota)
 \tag{UCE.15}
\]
supplies \(\mathrm{FinA}\), \(\mathrm{AutB}\), and
\(\mathrm{LemW}\) together, with exactly the original quantifiers.

\emph{Proof.}
Define \(B\) by the supplied local embedding of \(\Omega\).
It lies over \(\ell\).  Invoke \(\mathrm{GdEnd}\) once, with
the actual \(D\), \(B\), and fixed \(\iota\), obtaining the
single \(g_1\) in (UCE.15).  UCE12 supplies
\(\mathrm{FinA}(D,g_1)\) for all \(U,S\).

For \(\mathrm{AutB}\), now take any original \(U,S\), any
entry-algebra map \(\xi\), any actual pushed representation
\(r\) with its prescribed multiplier, and its Type-I hypothesis.
UCE16 gives an actual finite-coefficient ordinary point \(r'\)
conjugate to it by \(T(t)\).  UCE18 preserves Type I, and
UCE17 supplies irreducibility on every open subgroup of
\(H\), exactly the stronger antecedent in
\(\mathrm{AutEnd3}\).  Apply that field to the actual
\(\rho_K(r')\) obtained by restriction through the given
embedding \(K\hookrightarrow\Omega\).  It returns the original
\(\mathrm{IsAutCsd}\) assertion for
\(\rho_K(r')\epsilon^{-1}\).  UCE18 returns it to
\(\rho_K(r)\epsilon^{-1}\), which is precisely
\(D.\mathrm{IsAutI}(\iota,r)\).

For \(\mathrm{LemW}\), take the full original Type-II data,
including the middle-entry-one hypothesis and an arbitrary
\(W\) satisfying \(D.\mathrm{IsExt}(r,W)\).  Apply UCE16
to the same entry point.  UCE18 gives the Type-II ordinary
point with the same middle entry, and the actual full extension
\(W'\) in (UCE.14).  These are every hypothesis of
\(\mathrm{AutEnd2}\), including the trace-zero condition at
the original \(\sigma\).  That field supplies a newform of
weight \(2e+1\) attached to \(W'\epsilon^{-e}\).
UCE18 gives its attachment to \(W\epsilon^{-e}\), the literal
\(D.\mathrm{IsModular}(\iota,W)\).  Since \(W\) was arbitrary,
the universal extension quantifier in \(\mathrm{LemW}\) is retained.

Neither \(g_1\) nor \(\iota\) was changed in these arguments.
Only after their common choice were the local rings, normalized
rings, entry maps and pushed representations quantified.  Thus
the three fields hold simultaneously.  In particular they hold
when the supplied coefficient datum has \(e_{\mathcal O}=1\),
as required by the original joint target.  The implication uses
no extra ramification restriction on its intermediate finite DVRs.
% END UCE19

\subsection{Exact remaining producer and source boundary}

The closed results of this chapter are the local, ring, curve and
ordinary-representative assertions above, and their full original
same-object consumer.  A common \(g_1\) satisfying (UCE.15) has
not been constructed here.  Its three missing fields have different
roles: \(\mathrm{FinEnd}\) controls every characteristic-prime DVR
ordinary point using one finite-index subgroup; \(\mathrm{AutEnd3}\)
supplies actual polarized rank-three automorphy; and
\(\mathrm{AutEnd2}\) supplies actual newforms for every prescribed
outer-corner extension.  No result in this chapter obtains any of
those fields from finite weight zero or from an automorphic conclusion
that itself needs the joint package.

In the pinned original source the ordinary producer proceeds through
\texttt{Gd\allowbreak Frame\allowbreak Ex}, a common abelian CM frame, and
\texttt{Linked\allowbreak End\allowbreak E}, with its simultaneous rank-two and rank-three
runs, followed by the stated local and automorphic descent arguments.
The revised source route uses \texttt{Mod\allowbreak UAlt}, the two actual seed
producers, and the corresponding frame and descent assertions.
The previously completed selected rank-four Hida support does not
assert the general \texttt{Pfrom\allowbreak Aut} premise used to install that
\texttt{Mod\allowbreak UAlt}.  Those remaining producers are recorded separately
in this component's contracts and source manifest.  Consequently
\texttt{AUT\_\allowbreak UNRAM\_\allowbreak COEF} remains open, and UCE19 is not recorded as
its whole-contract closure.


\section{The original local wave packet and its global normalized product}
\label{sec:bessel-wave-packets-v49}

\status{This chapter proves the whole original \texttt{Coef\allowbreak Weighted\allowbreak Loc}
and \texttt{Wloc\allowbreak Exists}.  The latter has the original
\texttt{Bessel\allowbreak NVGap} antecedent.  It also proves the original
\texttt{WFrom\allowbreak Loc} and, using the completed setting blocks on their
original supplied \texttt{Forms0} scope, constructs
\texttt{Concrete\allowbreak Setting\allowbreak Exists} conditional on that
nonvanishing input.  No proof of \texttt{Bessel\allowbreak NVGap}, Bessel
multiplicity one, or the remaining Siegel--Weil comparisons is asserted.
The construction keeps every supplied Haar and quotient measure.
Its global vector is chosen before the later measured block.}

All Hilbert inner products in this chapter are conjugate linear in the
first argument.  Put \(n=2r+k+1\), \(R=r+1\), and
\[
 N=N(k,r)=4n^2\sqrt{2R}.
 \tag{BWP.1}
\]
The source's \(H_{k,r}\), after reversing the order of the final
isotropic basis, is the matrix
\[
 H_{k,r}=\begin{pmatrix}0&0&I_r\\0&J_V&0\\I_r&0&0\end{pmatrix},
 \qquad J_V=\operatorname{diag}(I_k,-1).
 \tag{BWP.2}
\]
This reordering is unitary, so it does not change the source operator
norm \(o(g)=\|g\|_{\mathrm{op}}\).  Write \(G=U(H_{k,r})\),
\(N_0=\texttt{Nset}(k,r)\), \(H_z=\texttt{UWset}(k,r,z)\),
and \(S=H_zN_0\).  The symbol \(N_0\) denotes the group and
\(N\) the exponent.  The original \(z\) belongs to the middle
space and satisfies \(-z^*J_Vz>0\).

The analytic foundations used below are Haar existence and uniqueness,
the quotient Haar theorem with its Borel-section formulation, Tonelli
and Bochner integration, and finite-dimensional compact representation
theory.  We also retain the same Harish--Chandra admissibility theorem
used in RFD1: every irreducible unitary representation of a real
reductive group has finite-dimensional compact isotypic spaces.
Borel's \emph{Automorphic forms on reductive groups}, printed
Section 10, paragraph labelled 8.5, states this precise input.
The reproducing function itself is proved below, not taken as a further
analytic theorem.  Source predicates and numerical assumptions are
recorded in the accompanying contracts and source excerpts.

\subsection{Explicit unipotent coordinates and compact reproduction}

% BEGIN BWP1
\paragraph{BWP1. A global chart on the actual unipotent subgroup.}
\label{claim:BWP1}
Let \(A\) run through upper unitriangular complex \(r\)-by-\(r\)
matrices, \(B\) through complex \(r\)-by-\((k+1)\) matrices,
and \(Z\) through skew-Hermitian \(r\)-by-\(r\) matrices.
Then
\[
 \Phi(A,B,Z)=
 \begin{pmatrix}
 A&B&(Z-\tfrac12BJ_V^{-1}B^*)A^{-*}\\
 0&I&-J_V^{-1}B^*A^{-*}\\
 0&0&A^{-*}
 \end{pmatrix}
 \tag{BWP.3}
\]
is a polynomial homeomorphism of a real vector space of dimension
\(d=2r^2+2rk+r\) onto the literal \(N_0\).

\begin{proof}
The source fixes the successive isotropic flag quotients pointwise,
and fixes the quotient of its orthogonal complement by that flag.
Its defining column equations therefore give the block shape
\([A,B,C;0,I,E;0,0,F]\), with \(A\) upper unitriangular.
The equation \(g^*H_{k,r}g=H_{k,r}\) first gives
\(F=A^{-*}\) and \(E=-J_V^{-1}B^*A^{-*}\).
The remaining equation is
\(CA^*+AC^*=-BJ_V^{-1}B^*\).
Thus \(Z=CA^*+\tfrac12BJ_V^{-1}B^*\) is skew-Hermitian,
and solving for \(C\) gives (BWP.3).  Conversely those equations
verify unitarity and every defining column condition.
The inverse reads \(A,B\) from their blocks and uses the displayed
quadratic formula for \(Z\).  Since the inverse of a unitriangular
matrix is a finite polynomial, both maps are polynomial.
The three real dimensions are \(r(r-1),2r(k+1),r^2\).
Their sum is the stated \(d\).  When \(r=0\) the chart is the
one-point zero-dimensional chart, and the same proof applies.
\end{proof}
% END BWP1

% BEGIN BWP2
\paragraph{BWP2. Haar measure and reciprocal-norm integrability in that chart.}
\label{claim:BWP2}
Lebesgue measure in (BWP.3) is Haar measure on \(N_0\).
Every measure satisfying the source's weak \texttt{Is\allowbreak Haar\allowbreak On}
on \(N_0\) is a finite positive multiple of it, and
\[
 \int_{N_0}o(m)^{-N}\,d\mu_{N_0}(m)<\infty.
 \tag{BWP.4}
\]
This proves the original \texttt{NInt\allowbreak Decay} at all \(k,r\).

\begin{proof}
For a fixed left multiplier \((A_0,B_0,Z_0)\), direct block
multiplication gives
\[
 \begin{split}
 A'&=A_0A,\qquad B'=A_0B+B_0,\\
 Z'&=A_0ZA_0^*+Z_0+
 \tfrac12(A_0BJ_V^{-1}B_0^*-B_0J_V^{-1}B^*A_0^*).
 \end{split}
 \tag{BWP.5}
\]
The real Jacobian is block triangular.  The first diagonal block
has determinant one by ordering upper entries row by row; the second
has determinant \(|\det A_0|^{2(k+1)}=1\); congruence by \(A_0\)
on skew-Hermitian matrices has determinant of absolute value
\(|\det A_0|^{2r}=1\).  Thus (BWP.5) preserves Lebesgue measure.
AMD8 promotes any given weak invariant measure to Haar, so uniqueness
gives a finite positive scalar multiple of this measure.

The matrix in (BWP.2) is a self-adjoint unitary involution.  Hence
\(g^{-1}=H_{k,r}g^*H_{k,r}\), \(o(g^{-1})=o(g)\), and
\(o(g)\ge1\) on \(G\).  Every entry of a matrix has absolute
value at most its operator norm.  The inverse-chart formula in BWP1
therefore gives, for any fixed Euclidean coordinate norm,
\[
 1+\|(A-I,B,Z)\|\le C_{k,r}\,o(\Phi(A,B,Z))^2.
 \tag{BWP.6}
\]
For \(r>0\), \(d< n^2\), whereas \(N\ge4n^2\); in
particular \(N/2>d\).  Raising (BWP.6) to exponent \(-N/2\)
majorizes the integrand in (BWP.4) by a constant times
\((1+\|t\|)^{-N/2}\), which is integrable in real dimension
\(d\) by radial integration.  For \(r=0\), the Haar measure
of the one-point group is finite and the integral is immediate.
\end{proof}
% END BWP2

% BEGIN BWP3
\paragraph{BWP3. The actual Bessel subgroup and its unitary character.}
\label{claim:BWP3}
The literal groups satisfy \(S=N_0\rtimes H_z=H_zN_0\), with
unique factorization in either order.  The group \(H_z\) is compact,
\(S\) is closed and unimodular, and
\(\psi=\texttt{psi\allowbreak Gen}(k,r,z,c)\) is a unitary character of
\(N_0\) invariant under conjugation by \(H_z\).  If \(\pi\)
has the original \texttt{Is\allowbreak HCType\allowbreak W}, then
\[
 \rho(hm)=\psi(m)\pi(h)\quad(h\in H_z,\ m\in N_0)
 \tag{BWP.7}
\]
is a continuous unitary representation of \(S\).

\begin{proof}
An element of \(H_z\) fixes the two isotropic blocks, acts on the
middle space by an element of \(U(k,1)\), and fixes its negative
vector \(z\).  Its restriction to \(z^\perp\) identifies this
group with \(U(k)\).  Reading the induced middle block of an element
of \(S\) gives its unique \(H_z\) factor, and the remaining
factor has exactly the columns in BWP1.  These readings are continuous;
equivalently the flag and middle-action equations show closedness.
The same argument is the local form of AMD4--AMD6.  The unipotent
group is unimodular, and the action modulus of the compact group
\(H_z\) on its Haar measure is a continuous homomorphism into
\(\mathbb R_{>0}\), hence one.  The semidirect product is unimodular.

For indices below the last one, the source function \(\ell_a(m)\)
is the \((a,a+1)\) entry of \(A\); for the last index it is
the last coordinate of \(Bz\).  Formula (BWP.5) shows that all
these functions add under multiplication: the final row of \(A_0\)
is the final coordinate row.  Conjugation by \(H_z\) fixes these
entries because it fixes the isotropic basis and \(z\).
Consequently \(\exp(i\operatorname{Re}\sum_a c_a\ell_a)\)
is the stated invariant unitary character, including the empty sum
when \(r=0\).  The defining frame of \texttt{Is\allowbreak HCType\allowbreak W}
transports a unitary finite-dimensional \(U(k)\) representation
to the whole \(H_z\).  This proves every assertion about (BWP.7).
\end{proof}
% END BWP3

% BEGIN BWP4
\paragraph{BWP4. A compact reproducer in every prescribed neighborhood.}
\label{claim:BWP4}
Let \(\sigma\) be any original \texttt{Is\allowbreak DS} representation,
and let \(w\) belong to its lowest compact-type subspace.  For
every neighborhood \(U\) of the identity there is
\(\alpha\in C_c(G,\mathbb C)\), supported in \(U\), such that
\[
 \int_G\alpha(a)\sigma(a)w\,d\mu_G(a)=w.
 \tag{BWP.8}
\]
This proves the source's \texttt{HCReprod\allowbreak DS}, and its weaker
compact measurable versions, without any gap assumption.

\begin{proof}
The source \texttt{Is\allowbreak DS} includes a complete Hilbert space and a
continuous irreducible unitary representation.  The given vector is
compact finite.  Harish--Chandra admissibility therefore makes the sum
\(E\) of its full compact isotypic spaces finite dimensional.  This
uses the named inherited foundation stated above, not a consequence
of square integrability alone.  If \(w=0\), take \(\alpha=0\).
Otherwise put \(m=\dim E\).  Compactness of the maximal compact
group \(K\) gives a conjugation-invariant neighborhood \(V\) of
one, containing one, with \(V^m\subset U\).
Take a continuous nonnegative approximate-identity function supported
in a sufficiently small subset of \(V\), of integral one, and
average its conjugates over \(K\).  Conjugation preserves Haar
measure.  The resulting \(f\in C_c(G)\) is conjugation invariant,
has integral one and support in \(V\).  Strong continuity, uniformly
on the unit sphere of the finite-dimensional \(E\), makes
\(T=\sigma(f)|_E\) arbitrarily close to the identity.  It preserves
\(E\), because \(\sigma(f)\) commutes with \(K\), and it is
invertible for a sufficiently small choice.

Write its characteristic polynomial as
\(p(t)=a_0+a_1t+\cdots+a_mt^m\); here \(a_0\ne0\).
Cayley--Hamilton gives
\(I_E=-a_0^{-1}\sum_{j=1}^m a_jT^j\).
Thus \(\alpha=-a_0^{-1}\sum_{j=1}^m a_j f^{*j}\) has
the required action.  Each convolution is continuous and compactly
supported, with support in \(V^j\subset V^m\subset U\).
Bochner integrability follows from compact support and unitarity.
The same argument works for any supplied Haar normalization.
\end{proof}
% END BWP4

\subsection{The exact coefficient exponent}

% BEGIN BWP5
\paragraph{BWP5. The original cubic gap dominates every required analytic bound.}
\label{claim:BWP5}
Suppose \(\lambda=\texttt{lam48}(k,r,v,b)\) is an original
Harish--Chandra parameter and \(\texttt{Gap\allowbreak Big}(k,r,b)\) holds.
Choose its witness \(t\), and put
\[
 M=n+(2r+2)t,\qquad h=2M^3+t+3.
 \tag{BWP.9}
\]
Every opposite-sign parameter difference has absolute value at least
\(h\).  If \(q_0=\min(p,q)>0\), then
\[
 h>npq+1,\quad h>16nq_0+8\sqrt{2q_0},\quad
 h>16nq_0+8N\sqrt{2q_0}.
 \tag{BWP.10}
\]

\begin{proof}
The first \(k-1\) signs are positive; the remaining signs alternate.
Their parameters are \(-b_j\).  Each consecutive difference of
the \(b_j\)'s is strictly larger than \(h\), so each distinct
tail difference is larger than \(h\).  A negative tail index has
odd tail index and therefore lies after the zeroth tail entry.
Strict decrease of \(\lambda\) places every initial parameter
above that zeroth entry; its distance to a negative tail entry is
again greater than \(h\).  These are all opposite-sign pairs.
Counting negative tail indices gives \(q_0\le q\le R\).

Since \(t\ge n+5\), we have \(M\ge3n\) and \(M\ge10R\).
Consequently \(npq\le M^3\) and
\(96n^2R\le2(3n)^2(10R)\le2M^3<h\).
Using \(N=4n^2\sqrt{2R}\), we obtain
\[
 16nq_0+8N\sqrt{2q_0}
 \le16nR+64n^2R\le80n^2R<h.
\]
Also \(\sqrt{2q_0}\le2R\), so
\(16nq_0+8\sqrt{2q_0}\le16nR+16R\le32nR<h\).
Finally \(2M^3+t+3>M^3+1\).  These prove (BWP.10),
including all small values permitted by the original predicates.
\end{proof}
% END BWP5

% BEGIN BWP6
\paragraph{BWP6. The whole original weighted coefficient statement.}
\label{claim:BWP6}
For every original \(k,r,v,b,\sigma,w\) in
\texttt{Coef\allowbreak Weighted\allowbreak Loc}, one has, with its prescribed real Haar measure,
\[
 \int_G |\langle w,\sigma(g)w\rangle|\,o(g)^N\,dg<\infty.
 \tag{BWP.11}
\]
No additional regularity or model-identification premise is required.

\begin{proof}
If \(q_0=0\), the real unitary group is compact and the integrand
is continuous, so the assertion follows.  Otherwise BWP5 supplies
the full DC.2 bounds.  DC13 therefore identifies the given literal
\texttt{Is\allowbreak DS} with the constructed representation, including its
prescribed lowest compact type.  The operator estimate DC6 holds
for every unit vector in that type, hence for the given \(w\).
Writing \(g=k\exp X\), \(a=\|X\|\), it bounds the
coefficient by
\(\exp[-(h/8)\min(a^2,a/\sqrt{2q_0})]\).
The operator norm satisfies \(o(g)\le e^a\), and DC6 bounds the
polar Jacobian by \(e^{n\sqrt{2q_0}a}\), times a polynomial
radial shell factor.  For large \(a\) the total exponential rate is
\[
 -h/(8\sqrt{2q_0})+N+n\sqrt{2q_0}<0
\]
by the final inequality in (BWP.10).  Radial integration is finite;
near zero continuity suffices.  Haar rescaling changes the answer
by a finite positive constant.  This proves the literal universal
\texttt{Coef\allowbreak Weighted\allowbreak Loc}, not merely an estimate for one chosen model.
\end{proof}
% END BWP6

% BEGIN BWP7
\paragraph{BWP7. Transport to the given hyperbolic coordinates and exact kernel.}
\label{claim:BWP7}
Retain the original \texttt{lam48}/\texttt{eps48} parameters and
the \texttt{Is\allowbreak HCParam}/\texttt{Gap\allowbreak Big} hypotheses of BWP5--BWP6.
For the original \texttt{Is\allowbreak DSHpk} and \texttt{Is\allowbreak Low\allowbreak Unit\allowbreak H} data
at these parameters,
every supplied weak Haar measure \(\mu_G\) is admissible in
(BWP.8) and (BWP.11).  Put
\[
 d_\sigma=\frac{\|w\|^4}
 {\int_G|\langle w,\sigma(g)w\rangle|^2\,d\mu_G(g)},
 \qquad e(g)=d_\sigma\overline{\langle w,\sigma(g)w\rangle}.
 \tag{BWP.12}
\]
Then \(d_\sigma>0\), this is precisely the source
\texttt{fdeg\allowbreak H}/\texttt{e\allowbreak Fun}, and
\[
 \int_G e(g)\sigma(g)w\,d\mu_G(g)=w,
 \qquad \int_G|e(g)|o(g)^N\,d\mu_G(g)<\infty.
 \tag{BWP.13}
\]

\begin{proof}
AMD8 first promotes \(\mu_G\) to genuine Haar measure on \(G\).
The Cayley matrix in each defining predicate is Euclidean unitary.
Conjugation by it preserves operator norms exactly and sends \(G\)
to the specified diagonal unitary group.  If the two predicates supply
different Cayley matrices, their ratio is Euclidean unitary and
preserves the same diagonal signature, hence belongs to its maximal
compact group.  Inner conjugation by that element preserves every
compact isotypic space and its minimum-norm type.  Thus the given
\(w\) and \(\sigma\) can use the same Cayley coordinates.
The transported measure is a finite positive multiple of the prescribed
Haar measure.  BWP4 and BWP6 now transport back with the exact original
\(\mu_G\).

DC12, and DC13 for the given representation, give Schur orthogonality
for this measure.  With \(\|w\|=1\), its formal degree is exactly
the quotient in (BWP.12).  The weighted integrability follows from
BWP6.  The Bochner integral in (BWP.13) exists by unitarity and
unweighted integrability.  Pairing it with any vector \(v\), Schur
orthogonality evaluates the result to \(\langle v,w\rangle\).
The integral is therefore \(w\).  This checks the conjugation and
formal-degree normalization in the literal source kernel.
\end{proof}
% END BWP7

\subsection{The packet for the original quotient measure}

% BEGIN BWP8
\paragraph{BWP8. The original quotient height and support.}
\label{claim:BWP8}
Let \(\mu_G,\mu_S,\nu\) be exactly the measures in
\texttt{Wloc\allowbreak Exists}.  Then \(\nu\) is supported on \(G\).
On \(G\), the source height
\[
 \eta(x)=\inf_{s\in S}o(sx)=\texttt{nu\allowbreak H}(k,r,z,x)
 \tag{BWP.14}
\]
is Borel measurable, satisfies \(1\le\eta(x)\le o(x)\),
and is left \(S\)-invariant.  The original quotient identity can
therefore be used with this exact height, without choosing a new \(\nu\).

\begin{proof}
Both weak measures are genuine Haar by AMD8.  Apply the given
quotient identity to the indicator of the complement of \(G\).
Multiplication by \(s\in S\subset G\) preserves that complement.
The inner integral equals \(\mu_S(S)>0\) on it and zero on
\(G\), so its \(\nu\)-measure is zero.  The norm inequalities
follow from BWP2 and \(1\in S\).  The infimum of continuous
functions is upper semicontinuous: its strict sublevel sets are the
unions of their strict sublevel sets.  Finiteness and the lower bound
thus give Borel measurability.  Replacing \(s\) by \(ss_0\)
proves \(\eta(s_0x)=\eta(x)\).  All statements concern the
source ambient measure, even when it is given by some other choice
of representatives.
\end{proof}
% END BWP8

% BEGIN BWP9
\paragraph{BWP9. The exact product scalar for the two Bessel measures.}
\label{claim:BWP9}
Choose any Haar measures \(\mu_H,\mu_N\) on the actual
\(H_z,N_0\), viewed on the ambient matrix space.  Their product
under \((h,m)\mapsto hm\) is \(c_S\mu_S\) for a finite
positive real number \(c_S\).  Thus, for every nonnegative
measurable \(f\),
\[
 \int\!\int_{H_z\times N_0} f(hmx)\,d\mu_H\,d\mu_N\,d\nu(x)
 =c_S\int_G f(g)\,d\mu_G(g).
 \tag{BWP.15}
\]
The same conclusions hold for any weak Haar measures quantified by
the original nonvanishing premise.

\begin{proof}
BWP3 gives the semidirect product and its trivial action modulus.
Consequently the product is Haar, as in AMD7; uniqueness relative to
the supplied \(\mu_S\) gives exactly the scalar stated.  Substitute
that identity into the supplied quotient formula.  AMD8 applies to
weak Haar measures and makes \(\mu_H\) finite because \(H_z\)
is compact, and \(\mu_N\) finite on compact sets.  No rescaling of
the given \(\mu_G,\mu_S,\nu\) has been made.
\end{proof}
% END BWP9

For the next claims fix the complete original real numerical setting,
the given \(\sigma,\pi,w\), and the measures above.  Let
\(x_0\in G\), \(\xi\in V_\pi\), and set
\[
 F(g)=\langle\sigma(x_0)w,\sigma(g)w\rangle,\qquad
 A(x)=\int_{H_z\times N_0}
 F(hmx)\overline{\psi(m)}\,\pi(h)^{-1}\xi\,d\mu_H(h)d\mu_N(m).
 \tag{BWP.16}
\]
These are the source's \texttt{Af} and its exact Bessel integrand.
On matrices outside \(G\), choose the value zero for the constructed
functions; all measures in the assertions are supported on \(G\).

% BEGIN BWP10
\paragraph{BWP10. Absolute convergence, continuity, and a uniform polynomial bound.}
\label{claim:BWP10}
The integral (BWP.16) converges absolutely for every \(x\in G\),
is continuous there, and satisfies
\[
 \int_{H_z\times N_0}\|F(hmx)\overline{\psi(m)}\pi(h)^{-1}\xi\|
 \,d\mu_Hd\mu_N\le C\,o(x)^N
 \tag{BWP.17}
\]
for a finite constant depending on the fixed data and \(x_0,\xi\).

\begin{proof}
Unitarity gives \(F(g)=\langle w,\sigma(x_0^{-1}g)w\rangle\).
By BWP7 and left translation,
\(L=\int_G|F(y)|o(y)^N\,d\mu_G(y)<\infty\).
Choose the compact reproducer \(\alpha\) from BWP4 for \(w\),
and let \(K\) be its compact support.  Applying its vector identity
to the coefficient gives
\(F(y)=\int_K\alpha(a)F(ya)\,d\mu_G(a)\).
For \(y=hmx\), set \(u=hmxa\).  Since
\(m=h^{-1}ua^{-1}x^{-1}\), BWP2 gives
\[
 o(m)\le o(h^{-1})o(u)o(a^{-1})o(x).
\]
The first and third factors are bounded uniformly on \(H_z\times K\).
Bound \(|\alpha|\) by its supremum and change variable from
\(a\) to \(u\) using left Haar invariance.  This proves
\[
 |F(hmx)|\le C_0L\,o(x)^N o(m)^{-N}.
 \tag{BWP.18}
\]
The remaining factors in the integrand have norm \(\|\xi\|\).
BWP2 and finiteness of \(\mu_H\) prove (BWP.17).
On every compact set of \(x\)'s, (BWP.18) is a single integrable
majorant.  Strong continuity of the representations and dominated
convergence prove continuity of \(A\).  This argument justifies
absolute convergence before every later change of variables.
\end{proof}
% END BWP10

% BEGIN BWP11
\paragraph{BWP11. The literal equivariance of the packet.}
\label{claim:BWP11}
For all \(h\in H_z,m\in N_0,x\in G\),
\[
 A(hmx)=\psi(m)\pi(h)A(x).
 \tag{BWP.19}
\]
In particular \(\|A\|\) is left \(S\)-invariant.

\begin{proof}
By BWP9 the integral is the same fixed scalar times
\(\int_S F(sx)\rho(s)^{-1}\xi\,d\mu_S(s)\).
For \(s_0\in S\), replace \(s\) by \(s's_0^{-1}\).
Right invariance, supplied by BWP3, and
\(\rho(s's_0^{-1})^{-1}=\rho(s_0)\rho(s')^{-1}\)
give \(A(s_0x)=\rho(s_0)A(x)\).
Absolute convergence is BWP10, and unitarity in BWP3 gives the norm
assertion.  Formula (BWP.19) is precisely \texttt{Is\allowbreak SPsi\allowbreak Equiv}.
\end{proof}
% END BWP11

% BEGIN BWP12
\paragraph{BWP12. Weighted integrability for the supplied quotient measure.}
\label{claim:BWP12}
The packet satisfies
\[
 \int_G\|A(x)\|\eta(x)^N\,d\nu(x)<\infty,
 \qquad \int_G\|A(x)\|\,d\nu(x)<\infty.
 \tag{BWP.20}
\]

\begin{proof}
Bound the norm of (BWP.16) by the integral of its norm.  For every
\(h,m\), BWP8 gives \(\eta(x)\le o(hmx)\).
Tonelli, followed by (BWP.15), therefore bounds the first integral by
\[
 c_S\|\xi\|\int_G |F(g)|o(g)^N\,d\mu_G(g),
\]
which is finite by the translated estimate in BWP10.  This uses only
nonnegative integration, so no unproved integrability is used to
exchange the integrations.  The second assertion follows from
\(\eta\ge1\).  Continuity and BWP8 supply the measurability needed.
\end{proof}
% END BWP12

% BEGIN BWP13
\paragraph{BWP13. Square integrability from the quotient height.}
\label{claim:BWP13}
The source squared norm satisfies
\(\texttt{l2SG}(\nu,A)=\int\|A\|^2d\nu<\infty\).

\begin{proof}
For every \(s\in S\), BWP10 and BWP11 give
\(\|A(x)\|=\|A(sx)\|\le C o(sx)^N\).
Taking the infimum, and using continuity and monotonicity of the
positive power, gives \(\|A(x)\|\le C\eta(x)^N\).
Multiply by \(\|A(x)\|\), integrate, and use BWP12.
The conclusion does not require a uniform bound on \(A\) over
the quotient; the original exponent supplies exactly the needed bound.
\end{proof}
% END BWP13

% BEGIN BWP14
\paragraph{BWP14. What the original nonvanishing input supplies.}
\label{claim:BWP14}
Assume the original \texttt{Bessel\allowbreak NVGap}.  There are
\(x_0,x\in G\) and \(\xi\in V_\pi\) for which the packet
(BWP.16) satisfies
\[
 A(x)\ne0,\qquad 0<\texttt{l2SG}(\nu,A)<\infty.
 \tag{BWP.21}
\]
In explicit terms, the retained premise says, for every original
numerical setting, representations, lowest unit vector and pair of
weak Haar measures, that some \(x_0,x\in G\) make
\[
 \int_{H_z\times N_0}
 \langle\sigma(x_0)w,\sigma(hm)\sigma(x)w\rangle
 \overline{\psi(m)}\,\pi(h)^{-1}\,d\mu_Hd\mu_N\ne0
 \tag{BWP.30}
\]
as an endomorphism of the given \(V_\pi\).  This endomorphism
nonvanishing is not proved here; the asserted packet positivity is
its consequence.

\begin{proof}
The original real setting has \(1\le k\) and \(k\) even,
hence \(2\le k\).  All its remaining parameter, interlacing,
half-integral difference, gap, and nonzero-character conditions are
exactly those quantified by \texttt{Bessel\allowbreak NVGap}.  Apply it to the
given representations, given lowest unit vector, and the measures
chosen in BWP9.  It gives a nonzero endomorphism integral for some
\(x_0,x\).  A nonzero endomorphism has a vector \(\xi\) with
nonzero image.  Evaluation on that vector is (BWP.16); evaluation
commutes with its absolutely convergent finite-dimensional integral.

Finiteness of the squared norm is BWP13.  To prove positivity for the
actual \(\nu\), choose a nonnegative nonzero compactly supported
continuous function \(f\) on an open set where \(\|A\|>\delta>0\).
Such a set exists by BWP10, and \(\int_G f\,d\mu_G>0\)
by full support of Haar measure.  If the squared norm were zero,
then \(A=0\) almost everywhere for \(\nu\).  Equivariance
implies \(f(sy)=0\) for all \(s\) whenever \(A(y)=0\).
The supplied quotient identity would give \(\int_G f=0\), a
contradiction.  This proves positivity without replacing the quotient
measure or presupposing full support of a chosen section.
\end{proof}
% END BWP14

% BEGIN BWP15
\paragraph{BWP15. A justified convolution transfer for every base point.}
\label{claim:BWP15}
Suppose a measurable scalar function \(\beta\) satisfies
\[
 \int_G|\beta(g)|o(g)^N\,d\mu_G(g)<\infty,
 \qquad \int_G\beta(g)\sigma(g)w\,d\mu_G(g)=w.
 \tag{BWP.22}
\]
Then for every \(x\in G\) the following integral is absolutely
convergent and
\[
 \int_G\beta(g)A(xg)\,d\mu_G(g)=A(x).
 \tag{BWP.23}
\]

\begin{proof}
BWP10 and submultiplicativity give the bound on the integral of the
norm of the joint integrand
\[
 \int_G |\beta(g)|
 \int_{H_z\times N_0}|F(hmxg)|\,\|\xi\|\,d\mu_Hd\mu_Nd\mu_G
 \le C o(x)^N\int_G|\beta(g)|o(g)^N\,d\mu_G<\infty.
\]
All joint functions are measurable by continuity on the supported
groups and the assumed measurability of \(\beta\).
Bochner Fubini therefore applies.  For fixed \(h,m\), the inner
integral is obtained by applying the bounded linear map
\[
 v\longmapsto
 \langle\sigma(x_0)w,\sigma(hm)\sigma(x)v\rangle
 \overline{\psi(m)}\pi(h)^{-1}\xi
\]
to the vector integral in (BWP.22).  It gives the integrand of
\(A(x)\).  This proves (BWP.23) simultaneously at every
\(x\in G\), with the stated absolute integrability.
\end{proof}
% END BWP15

% BEGIN BWP16
\paragraph{BWP16. Reproduction by the literal formal-degree kernel.}
\label{claim:BWP16}
For every \(x\in G\), the source \texttt{e\allowbreak Fun} convolution
with \(A\) is integrable and equals \(A(x)\).

\begin{proof}
Take \(\beta=e\) from (BWP.12).  Both hypotheses of BWP15
are exactly BWP7.  Thus its conclusion is
\(\texttt{rconv}(\mu_G,\texttt{e\allowbreak Fun}(\mu_G,\sigma,w),A,x)=A(x)\),
with the full integrability clause.  In particular the same formal
degree is computed from the given \(\mu_G\), rather than borrowed
from a differently normalized group measure.
\end{proof}
% END BWP16

% BEGIN BWP17
\paragraph{BWP17. Compact reproduction of the packet.}
\label{claim:BWP17}
There is a bounded measurable \(\alpha\), supported in a compact
set \(\Omega\subset G\) containing one, such that for every
\(x\in G\) its convolution with \(A\) is integrable and
equals \(A(x)\).

\begin{proof}
Use BWP4 and BWP7 to choose the continuous compact reproducer of
\(w\) for the actual \(\mu_G\).  A continuous compactly supported
function is bounded and has finite integral after multiplication by
the continuous weight \(o(g)^N\).  Adjoin one to its compact support
to obtain \(\Omega\), and extend by zero on the ambient matrix
space.  Closedness of \(G\) gives measurability of that extension.
BWP15 proves both assertions about the convolution at every base point.
\end{proof}
% END BWP17

% BEGIN BWP18
\paragraph{BWP18. The entire original local packet producer.}
\label{claim:BWP18}
The whole original \texttt{Wloc\allowbreak Exists} holds.  Namely, under its
literal \texttt{Bessel\allowbreak NVGap} and \texttt{Coef\allowbreak Weighted\allowbreak Loc}
antecedents, and for every supplied local numerical setting,
representation, lowest unit vector, and triple of measures, there
exists \(W\) having all seven original conclusions.

\begin{proof}
Choose \(x_0,\xi\) by BWP14 and put
\[
 L_2=\int_G\|A(x)\|^2\,d\nu(x),\qquad W=L_2^{-1/2}A.
 \tag{BWP.24}
\]
BWP13--BWP14 give \(0<L_2<\infty\), so the positive real
normalizing scalar is defined.  Continuity is BWP10; literal
\(S,\psi,\pi\) equivariance is BWP11; integrability and reproduction
by the exact \texttt{e\allowbreak Fun} is BWP16.  BWP12 supplies both ordinary
and \(\texttt{nu\allowbreak H}^{N1H}\)-weighted \(L^1\) integrability.
The squared \(L^2\) integral is
\(L_2^{-1}L_2=1\), precisely the source \texttt{l2SG}, rather
than an unnormalized Hilbert norm.  BWP17 supplies the measurable,
bounded, compactly supported reproducer, with support enlarged to
contain one.  Scalar multiplication preserves every convolution and
equivariance equality.  These are all the original clauses, on the
same \(\mu_G,\mu_S,\nu\).  BWP6 in addition independently supplies
the coefficient antecedent, so nonvanishing is the only unproved
local input left by this construction.
\end{proof}
% END BWP18

\subsection{Assembly on the given global data}

% BEGIN BWP19
\paragraph{BWP19. The quotient product formula with its actual scalar.}
\label{claim:BWP19}
For a finite product \(G=\prod_vG_v\), \(S=\prod_vS_v\)
of the preceding unimodular groups, fix the original product group
measure and any supplied \(S\) Haar measure and quotient measure
\(\nu_Q\).  Choose local Haar measures \(\mu_{S,v}\) and
their normalized quotient measures \(\nu_v\).  There is one
finite positive scalar \(C\), independent of the test functions,
such that for every family of nonnegative measurable
\(S_v\)-invariant functions \(f_v\),
\[
 \int_G\prod_v f_v(x_v)\,d\nu_Q(x)
   = C\prod_v\int_{G_v} f_v(x_v)\,d\nu_v(x_v).
 \tag{BWP.25}
\]

\begin{proof}
Haar uniqueness gives \(\mu_S=a\prod_v\mu_{S,v}\),
\(0<a<\infty\).  The product of the local quotient measures
satisfies the quotient formula for \(\prod_v\mu_{S,v}\),
by Tonelli.  Thus the constant is \(C=a^{-1}\).
For clarity, the given ambient \(\nu_Q\) need not coincide
as a measure with the product of chosen section measures.
What is unique is its pushforward to \(S\backslash G\).
Here is the normalization argument.  The quotient Haar theorem
used in ASG17 gives, for this closed unimodular subgroup, a
nonnegative Borel cutoff \(b\) with
\(\int_S b(sx)\,d\mu_S(s)=1\).
It can also be constructed from a locally finite cover of the
second-countable quotient, compactly supported functions having
positive orbit averages on that cover, division by those averages,
and a partition of unity.  For any nonnegative invariant \(f\),
the supplied quotient formula applied to \(bf\) identifies
\(\int f\,d\nu_Q\) with \(\int_G bf\,d\mu_G\).
The product quotient formula gives
\(\int_G bf\,d\mu_G=a^{-1}\int f\,d(\prod_v\nu_v)\).
This proves equality on all such
functions, including infinite integrals.  Taking \(f=\prod_v f_v\)
and using Tonelli proves (BWP.25).
\end{proof}
% END BWP19

% BEGIN BWP20
\paragraph{BWP20. Every supplied global frame has the required local packet.}
\label{claim:BWP20}
Fix the original blocks \(A,L,P,R_f,B_u,\Pi,M_u\) and
\(\texttt{Gap\allowbreak Hyp}(D)\).  Suppose the local numerical gaps and
\texttt{Bessel\allowbreak NVGap} hold.  For each original word \(u\) and
real place \(v\), BWP18 gives a packet in the actual supplied
space \(\Pi.V_{\mathrm{loc}}(u,v)\), for its actual compact
representation and lowest vector, with quotient measure \(\nu_v\).
It has the original local character, convolution, weighted integral,
and compact reproduction after transport to the given frame.

\begin{proof}
Use the given vectors \(e_i,f_i\), positive orthonormal frame
\(w_j\), and the vector \(z_u/\sqrt{T_u}\) to form the
matrix \(Q_v\).  Their original pairing identities say exactly
\(Q_v^*J_{\epsilon_v}Q_v=H_{k,r}\), so \(Q_v\) is
invertible and conjugation identifies the two unitary groups.
The same computation identifies the literal flag stabilizer,
unipotent subgroup and compact factor in BWP3.  The supplied compact
representation becomes exactly \(\texttt{Is\allowbreak HCType\allowbreak W}\) in these
coordinates: its matrix on the positive frame is the middle positive
block of the conjugated compact element.

There are two different coordinate changes here.  Choose a
Euclidean unitary Cayley matrix \(P_0\) with
\(P_0^*J_{\epsilon_v}P_0=H_{k,r}\), obtained by replacing each
hyperbolic pair by its normalized sum and difference, followed by
the required permutation of signs.  Define
\[
 C_G(g)=P_0^{-1}gP_0,\qquad C_S(s)=Q_v^{-1}sQ_v,
 \qquad T(x)=Q_v^{-1}xP_0.
 \tag{BWP.28}
\]
The representation used for the local packet is
\(\sigma_B(M)=\sigma_v(P_0MP_0^{-1})\), on the given Hilbert
space with the unchanged \(w_v\).  By this unitary Cayley matrix
it has the literal \texttt{Is\allowbreak DSHpk} and \texttt{Is\allowbreak Low\allowbreak Unit\allowbreak H}.
The generally nonunitary frame matrix \(Q_v\) is used for the
subgroup and the base point, not to change that lowest vector.
Direct multiplication gives the two identities
\[
 T(sx)=C_S(s)T(x),\qquad T(xg)=T(x)C_G(g).
 \tag{BWP.29}
\]

More explicitly, the signature is \(\texttt{eps48}(k,r)\);
the parameter is \((v_0,\ldots,v_{k-2},-b_0,\ldots,-b_{2r+1})\).
The compact highest weight has the two original interlacing fields.
They are precisely \texttt{Interlacing\allowbreak Condition}.  Their parities
give \texttt{Half\allowbreak Integral\allowbreak Difference}.  The vector becomes
\(\sqrt{T_u}\) times the final negative standard vector.
For the character take every \(c_a=4\pi/g_D\).
The frame identities give
\(\ell_a(Q_v^{-1}sQ_v)=\mathrm{rootCoord}_a(s)\).
Consequently its exponent is
\[
 i\operatorname{Re}\sum_a(4\pi/g_D)\ell_a
 =\frac{2\pi i}{g_D}\sum_a
       (\mathrm{rootCoord}_a+\overline{\mathrm{rootCoord}_a}),
\]
which is exactly the original \texttt{psi\allowbreak R}.  Since \(g_D>0\),
every \(c_a\ne0\).  This verifies the actual numerical hypotheses
of the local producer, including its character normalization.

Choose local subgroup and quotient measures with the prescribed
\(\mu_{G,v}\), as in ASG17.  Push the group measure by \(C_G\),
the subgroup measure by \(C_S\), and the quotient measure by \(T\).
They satisfy the exact quotient formula: (BWP.29) gives its right
side, while \(T=L_{Q_v^{-1}P_0}\circ C_G\) gives its left side
by Haar left invariance.  Thus BWP18 applies to exactly this triple.
If its packet is \(W_B\), set \(W_v(x)=W_B(T(x))\).
The first identity in (BWP.29) gives equivariance and the second gives
the convolution formula with the given \(w_v\).  The formal-degree
integral under \((C_G)_*\mu_{G,v}\) is literally the original
one by change of variables, so its scalar is unchanged.
Both quotient integrals, including their normalization, pull back
under \(T\).  The compact reproducer pulls back under \(C_G\),
which is a group homeomorphism preserving one.

Finally, put \(c_v=\max\{1,\|Q_v\|\|P_0^{-1}\|\}\).
Writing \(y=T(x)\) and \(s_B=C_S(s)\) gives
\(sx=Q_vs_ByP_0^{-1}\), and hence
\(\eta_v(x)\le c_v\eta_B(T(x))\).
This finite positive constant transfers the exact weighted estimate
with exponent \(N\).  Continuity follows from the displayed maps.
No compact representation, lowest vector, Hilbert space, or supplied
group measure has been replaced.
\end{proof}
% END BWP20

% BEGIN BWP21
\paragraph{BWP21. The normalized global vector on the supplied tensor data.}
\label{claim:BWP21}
For the data of BWP20 there is a \(W_b:\texttt{WBlk}\) with
every literal \texttt{Is\allowbreak Wu} field, on the actual \(\nu_Q(u)\)
and the given tensor map, chosen independently of any later
\(M_b:\texttt{MBlk}\).

\begin{proof}
Fix a word and apply BWP19 to the supplied subgroup measure,
which is the product of its given \(N\) and \(H\) measures by
AMD7--AMD8.  Let \(C>0\) be its scalar in (BWP.25).  The field
has a real place; fix one \(v_0\).  Multiply the local packet at
\(v_0\) by \(C^{-1/2}\), leaving every other one unchanged.
Call these functions \(W_v\), and define
\[
 W_u(x)=\operatorname{tens}_u\bigl((W_v(x_v))_v\bigr).
 \tag{BWP.26}
\]
The supplied tensor-inner-product identity gives
\(\|W_u(x)\|=\prod_v\|W_v(x_v)\|\).
It also implies multilinearity of the given tensor map: pair an
additivity or scalar-linearity difference in one coordinate with
every pure tensor.  The inner-product product formula makes each
pairing zero, and the supplied spanning property makes the difference
zero.  The product norm identity then gives continuity of this
finite multilinear map.
The finite tensor product is complete and finite dimensional,
as in the original representation block.

Every \(s\) in the actual stabilizer factors place by place as
\(h_vm_v\).  The supplied global representation is trivial on
the unipotent factors and agrees with the tensor product of the given
compact representations.  Multiplying (BWP.19) over the places
therefore gives the literal global character-equivariance identity.
The formal-degree kernel is exactly the product of the local kernels.
Each local convolution is absolutely integrable by BWP18.  Tensor
multilinearity, its product norm identity, and finite-product Fubini
show that the global convolution is integrable and equals (BWP.26).
This checks the original kernel, including its formal degrees.

Continuity follows from continuity of each factor.  Formula (BWP.25)
applied to the product of their invariant norms gives finite \(L^1\).
Applied to their squared norms it gives
\(C\cdot C^{-1}\cdot\prod_v1=1\).
These are all five \texttt{Is\allowbreak Wu} fields.  The constructed local
functions also give the exact \texttt{Wloc}, continuity, and
\texttt{Wu\_\allowbreak tens} fields of \texttt{WBlk}.
Make these choices for each original word.  The construction used
only the original preceding blocks and their measures; it did not
use a later adelic measured block.
\end{proof}
% END BWP21

% BEGIN BWP22
\paragraph{BWP22. Decay of that same global vector for every later measured block.}
\label{claim:BWP22}
The \(W_b\) chosen in BWP21 satisfies the entire original
\texttt{Wu\allowbreak Decay} for every later \(M_b:\texttt{MBlk}\).
In particular no uniqueness theorem for Bessel models is needed
for this existential constructor.

\begin{proof}
Let \(\eta_v(x_v)=\inf_{s_v\in S_v}\|s_vx_v\|\), and
let \(\eta_\infty\) be the original \texttt{nu\allowbreak Inf}.
Independent choices of the finitely many \(s_v\)'s give
\(\eta_\infty(x)\le\prod_v\eta_v(x_v)\).
All factors are at least one in the diagonal unitary coordinates.
The original exponent is
\(N_X=4n^2\sqrt{2\max_v q_v}\le N\), since
\(q_v\le r+1\) by the sign count in BWP5.
The norm comparison in BWP20 therefore implies finite
\(\int\|W_v\|\eta_v^{N_X}d\nu_v\), using the local
\(N\)-weighted estimate and the harmless finite factors
\(c_v^{N_X}\).  Formula (BWP.25), the tensor norm identity, and
the preceding infimum inequality now give
\[
 \int\|W_u(x)\|\eta_\infty(x)^{N_X}\,d\nu_Q(x)<\infty.
 \tag{BWP.27}
\]
This is the full original \texttt{Wu\allowbreak Weighted} clause.

For each place take the compact reproducer supplied by BWP18 after
transport, and put \(\alpha_\infty((g_v))=\prod_v\alpha_v(g_v)\).
It is measurable, bounded, and supported in the compact product
\(\prod_v\Omega_v\), which contains one.  The local integrability
clauses and tensor norm identity justify Fubini; their reproducing
equalities give the global one at every \(x\).  This is exactly
\texttt{Wu\allowbreak Reprod}.  Neither clause uses \(M_b\): the original
\(\nu_Q\), \(W_u\), infinite group measure and exponents belong
to the already chosen blocks.  They consequently hold for every
later \(M_b\), in the required order \(\exists W_b\,\forall M_b\).
This proves decay for the constructed vector; it asserts nothing
about uniqueness or every arbitrary \texttt{Is\allowbreak Wu} vector.
\end{proof}
% END BWP22

% BEGIN BWP23
\paragraph{BWP23. The original global gap supplies every local cubic gap.}
\label{claim:BWP23}
The full original \texttt{Gap\allowbreak Big\allowbreak Of\allowbreak Gap\allowbreak Hyp}(D) holds, for every
original dimension identification and every real place.

\begin{proof}
Write \(n=D.R.r=2r+k+1\) and \(m=2r+2\).  The original
rank conditions give \(n\) odd and \(m\le n\).
The source \texttt{Gap\allowbreak Hyp} says \(n\ge1801\) and every
successive \(b\)-difference is at least \(16n^{19}\).
Set \(t=n+6\).  It is odd, is at least \(n+5\), and
\(n+t\ge1801\).  Since \(n\ge7\),
\(M=n+mt\le n+n(n+6)\le2n^2\).
It follows that
\[
 2M^3+t+3\le16n^6+n+9<16n^{19}.
\]
Thus each of the required strict cubic-gap inequalities follows from
the original weak lower bound.  The source \(b_{48}\) is the
same \(b\)-tuple reindexed by the supplied dimension equality;
transporting a finite index changes none of these inequalities.
This proves the entire original predicate, with no new gap assumption.
\end{proof}
% END BWP23

% BEGIN BWP24
\paragraph{BWP24. The exact remaining condition for the complete analytic setting.}
\label{claim:BWP24}
The original \texttt{WFrom\allowbreak Loc}(D) holds for every \(D\).
Using the already reconstructed blocks for an original \(D\)
with its supplied \(\texttt{Forms0}(D,D.u_0)\), the retained
\texttt{Bessel\allowbreak NVGap} implies
\texttt{Concrete\allowbreak Setting\allowbreak Exists}(D).  The original \texttt{Forms0}
scope of the lattice constructor is preserved.  The nonvanishing statement,
and the later spectral, orbital and arithmetic comparisons, remain
separate producers.

\begin{proof}
Under the original antecedents of \texttt{WFrom\allowbreak Loc}, BWP20--BWP22
construct the required \(W_b\) for every supplied sequence of
blocks, with decay for every later \(M_b\).  BWP23 supplies the
local gaps from the global one.  These are exactly the quantifiers
and conclusion of \texttt{WExists}; the local packet input is
BWP18, and its coefficient input is BWP6.

For the full setting on that supplied \texttt{Forms0} scope,
choose the group instance, lattice, additive
character, reference frame, transport, bundle, compact representation
and quotient measures by ASG1--ASG17, in their original order.
Apply the preceding constructor to choose \(W_b\).  Only then
apply AMD12 to obtain \(M_b\) and one bridge with the literal
\texttt{Meas\allowbreak Tie}.  BWP22 gives decay for precisely that measured
block.  Finally the infinite projector kernel is a finite product
of the kernels in (BWP.12).  BWP6 and \(o(g)\ge1\) give local
absolute integrability; product Fubini gives
\(\operatorname{Integrable}(X.\mathrm{ph},\mu_\infty)\).
Thus the source \texttt{Concrete\allowbreak Setting\allowbreak Exists} conclusion, including
one setting, one bridge, the exact ties, decay and kernel
integrability, follows.
No step proved the nonzero Bessel endomorphism required in BWP14.
The conclusion therefore retains \texttt{Bessel\allowbreak NVGap} and the
stated original \texttt{Forms0} input to the lattice construction.
Nor does a completed analytic setting alone
identify the orbital measures or prove the original Siegel--Weil
and common-family counting statements.
\end{proof}
% END BWP24


\section{The original finite holomorphic CM uniformization family}

\paragraph{Scope and conventions.}
This chapter proves the missing family theorem, rather than assuming it
from its individual-point consequence.  Throughout, \(E,H,D,l,X_o\)
are the original data, with \(X_o:\texttt{Moduli}\ D\ l\) supplied.
Put \(R=\mathcal O_E\), \(f=[E:\mathbb Q]/2\), \(r=p+q\geq1\),
\(g=fr\), \(d=fpq\), and \(N=l.M'=\operatorname{lcm}(M,A)\).
The fixed map \(\sigma:K_A(l)\to\mathbb C\) satisfies
\(\sigma(\rho)=\mathrm{rt}_0\).  The \emph{given ordered} type
\(\mathrm{ty}:\operatorname{Fin}(g)\to\operatorname{Hom}(E,\mathbb C)\)
has multiplicities \(p\) on \(\Psi\), and \(q\) on \(\bar\Psi\).
We retain every original level condition, including odd \(M\geq3\),
the two relatively prime half-ideals, and arbitrary \(A>0\).
Hermitian forms are complex linear in their first variable.

Write \(L=D.\mathrm{La}\), \(B=D.\mathrm{EF}\).
For another datum \(D'\) write \(L'\) and \(B'\).
The notation
\[
 \mathcal U(D',t,P;y,e)
 \tag{HUF.1}
\]
means the \emph{literal} \texttt{Unif} predicate in
\texttt{Family\allowbreak Unif\allowbreak Of\allowbreak Parts}: \(y\) is over \(\sigma\);
\(P:L'\to\mathbb C^g\) is an original \texttt{Is\allowbreak Period};
\(e:\mathbb C^g\to G_y(\mathbb C)\) is a surjective holomorphic
group map with kernel \(P(L')\), with both source dual-number
derivative bijections and the exact diagonal \(R\)-action; and
\[
 \begin{split}
 e_n(e(Pa/n),e(Pb/n))&=\exp(-2\pi i B'(a,b)/n)\quad(n>0),\\
 \eta(a)&=e(P(ta)/M)\quad(a\in\bar{\mathfrak b}L),\\
 \eta_A(a)&=e(P(ta)/A)\quad(a\in L).
 \end{split}                                                     \tag{HUF.2}
\]
The first line is the source's universally tested
\texttt{Out.\allowbreak Weil\allowbreak Is}, not an unspecified pairing convention.
All functions may be assigned arbitrary values outside the open sets
on which their conditions are imposed.

\paragraph{Inherited results and named foundations.}
We use PUN6, PUN8 and PUN15, and the actual individual-point,
line-normalization, equality and nonemptiness proofs PUN20--PUN24.
In particular the original primitive identity fixes the
Appell--Humbert semicharacter to
\(\alpha(Pa)=\exp(\pi i B'(a,w_0a))\).
The supplied \texttt{Moduli} record retains its algebraic representing
property on reduced schemes, smoothness of relative dimension \(d\),
quasiprojectivity, universal object and stated infinitesimal condition.
Existence of that record is not proved here.

The retained classical foundations are Appell--Humbert and its
positivity criterion, the Riemann algebraicity theorem, projective
GAGA, elementary holomorphic differential equations and the
implicit-function theorem.  Two additional precise references are
used below: Grauert's analytic direct-image and base-change theorem
for a proper map, a sheaf flat over a reduced base, and constant
fiber cohomology dimension; and finite connected-component
decomposition for semialgebraic sets.  These are external foundations,
not numbered claims of a newly written proof of those classical
theorems.  Exact primary-source locators appear after their applications
and in the component metadata.

\subsection{Reverse extraction on the original moduli space}

% BEGIN HUF1
\paragraph{HUF1. The source holomorphic and tangent tests are the analytic ones.}
\label{claim:HUF1}
Let \(Z\) be a finite-type scheme over a subfield \(k\subseteq\mathbb C\),
with a fixed embedding \(k\to\mathbb C\), and let
\(Z_\mathbb C\) denote the corresponding base change.
For a map \(y:B\to Z_\mathbb C(\mathbb C)\), where
\(B\subseteq\mathbb C^m\) is open, the original
\(\texttt{Is\allowbreak Hol\allowbreak To}(B,y)\) is equivalent to holomorphy into
\(Z_\mathbb C^{\mathrm{an}}\).
At a smooth point, the original \texttt{Is\allowbreak Deriv\allowbreak At0} is the
usual derivative interpreted as a map from the dual numbers.
Consequently a holomorphic local biholomorphism at \(0\) has
both of the exact unique-existence clauses in \texttt{ex\_\allowbreak diff}.

\begin{proof}
On an affine open choose finitely many algebra generators and its
closed immersion into affine space.  The source test says that the
inverse image of this open is open and that all those coordinate
functions are holomorphic.  Their polynomial equations hold pointwise,
so this is precisely a holomorphic map into the closed analytic
subspace.  Coefficients introduced by base change are complex
constants; they cause no additional test.  Conversely a holomorphic
map has holomorphic pullbacks of regular functions.  This remains true
on any scheme open: locally its regular functions are quotients of
polynomials with nonzero denominator.  These local statements glue
and give the quantifier over every open and every section in
\texttt{Is\allowbreak Hol\allowbreak To}.

A morphism from \(\operatorname{Spec}\mathbb C[\epsilon]/(\epsilon^2)\)
with specified closed point is, on its local ring, exactly
\(a\mapsto a(y)+\epsilon\,\delta(a)\), with \(\delta\) a complex
derivation.  Holomorphic directional differentiation gives that
derivation.  Smooth algebraic coordinates, obtained from a standard
smooth presentation and the analytic implicit-function theorem, make
this an isomorphism between analytic and algebraic tangent spaces.
Evaluating every local regular function is exactly the equality
required in \texttt{Is\allowbreak Deriv\allowbreak At0}; restricting an open does not change
the germ or its derivative.  If the analytic derivative is invertible,
each direction determines a unique dual-number point, and each
vertical dual-number point determines a unique direction.
For a group scheme over \(y\) the condition of being vertical is
exactly the source's \texttt{eps\allowbreak Pt} condition.
\end{proof}
% END HUF1

% BEGIN HUF2
\paragraph{HUF2. The actual fibers and their relative Lie eigenspaces.}
\label{claim:HUF2}
The complex base change of the supplied universal group is a proper
smooth family of connected abelian varieties of dimension \(g\).
Its Lie bundle is the direct sum of the original embedding
eigenbundles, of ranks \(p\) and \(q\) in the order prescribed
by \(\mathrm{ty}\).  Locally on an analytic base, those eigenbundles
have simultaneous holomorphic frames with any prescribed eigenbasis
at one point.  The universal rigidified line is relatively ample
and its analytic restriction to each fiber has a positive integral
Appell--Humbert form.

\begin{proof}
The \texttt{SAB.\allowbreak semi} field gives a smooth group with connected
semiabelian geometric fibers.  The additional properness field
makes its toric part trivial: a positive-dimensional torus cannot
be a closed subgroup of a proper group, since it is affine and
not proper.  Thus the fibers are abelian varieties.  The
\texttt{Obj.\allowbreak geom.\allowbreak Det\allowbreak Cond} clause computes their Lie multiplicities.
Summing them over the embeddings gives \(g=f(p+q)\);
smoothness with this fiber dimension gives the asserted relative
dimension.

Choose an integral primitive element for \(E/\mathbb Q\).
Its distinct conjugates allow the usual interpolation polynomials
in its Lie action to construct the embedding idempotents.
They are holomorphic bundle endomorphisms.  Their images are
holomorphic subbundles of the constant ranks given by
\texttt{Det\allowbreak Cond}.  Choose local frames and multiply them by constant
invertible matrices to match the prescribed basis at the initial
point.  Ordering these frames according to the supplied type gives
the diagonal \(R\)-action, without permuting the original coordinates.

The printed \texttt{Rel\allowbreak Ample} fields have a useful concrete meaning.
On every affine base open they give a power of the line, finitely
many generating sections, and affine nonzero loci whose coordinate
rings are generated by the section ratios.  The associated map
to projective space is therefore a closed immersion: on each
standard projective open its coordinate-ring map is surjective.
In particular its analytification is proper, since it is closed
in the product with the compact projective space.  The same
construction after base change gives the proper analytic family
used here.  The line's rigidification and primitive, symmetric
and adjoint identities are identities of the actual universal
object and hence persist.  On a fiber, Appell--Humbert and the
positivity criterion identify the ample line with a positive
Hermitian form and an integral alternating imaginary part.
\end{proof}
% END HUF2

% BEGIN HUF3
\paragraph{HUF3. Exponential, full period lattice and the original Weil sign.}
\label{claim:HUF3}
Let \(G\) be any of the complex fibers in HUF2.
There is a holomorphic exponential
\[
       \exp_G:\operatorname{Lie}G\longrightarrow G(\mathbb C)
 \tag{HUF.3}
\]
with derivative the identity, surjective with discrete full
lattice kernel \(\Lambda\).  For an Appell--Humbert form
\(H_G\) of its line, put \(B_G=\operatorname{Im}H_G|_\Lambda\).
Then \(B_G\) is a perfect integral alternating form.  In the
original commutator convention its Weil scalar is
\[
 e_n(\exp_G(\lambda/n),\exp_G(\mu/n))
                     =\exp(-2\pi iB_G(\lambda,\mu)/n).
 \tag{HUF.4}
\]
The conclusion includes the universal condition on generators
in \texttt{Weil\allowbreak Is}.

\begin{proof}
Left-invariant holomorphic vector fields commute because the group
is commutative.  Their local flows give a local holomorphic
homomorphism with identity derivative.  For an arbitrary vector
\(z\), choose a positive integer \(a\) with \(z/a\) in its
domain and define \(\exp_G(z)=[a]\exp_G(z/a)\).
The local homomorphism law makes the choices compatible, by
passing to a common multiple.  It also proves global
holomorphy and addition.  Its image is an open subgroup,
and connectedness forces it to be the whole group.
The derivative is invertible everywhere by translation, so the
kernel is discrete.  The induced local biholomorphism
\(\operatorname{Lie}G/\Lambda\to G(\mathbb C)\) is bijective
and hence an analytic isomorphism.  Compactness of \(G(\mathbb C)\)
makes \(\Lambda\) a full lattice: a discrete subgroup of a real
vector space has a basis for its real span, and a proper span
would leave a noncompact real-vector-space quotient.

Appell--Humbert supplies the integral form \(B_G\).
Its positivity makes it nonsingular.  The kernel of the line's
polarization consists analytically of the quotient of the
integral \(B_G\)-dual lattice by \(\Lambda\).  This follows
directly by comparing the factors for a translate of the line:
the translation difference is trivial exactly for a vector
pairing integrally with \(\Lambda\).  Such a trivial analytic
line is algebraically trivial by projective GAGA.
The source \texttt{KTriv}, applied with coefficient ring
\(\mathbb C\), forces the corresponding point to be zero.
Thus the dual lattice is \(\Lambda\), proving perfection.

For \(x=\lambda/n\), the section of the source's translation
quotient is represented by
\(a_x(z)=\exp(-\pi nH_G(z,x))\).
The semicharacter cancels from that quotient.  Its factors
are exactly those required to descend this nonvanishing
section.  As in the explicit calculation PUN21, the source
commutator, in its printed order, is
\[
 \frac{a_x(z)a_y(z+x)}{a_y(z)a_x(z+y)}
   =\exp[-\pi n(H_G(x,y)-H_G(y,x))]
   =\exp(-2\pi iB_G(\lambda,\mu)/n).
\]
GAGA makes these sections algebraic.  Every other generator
differs by a nonzero constant on the connected proper fiber,
so the same scalar holds for every generator.  This proves
the literal \texttt{Weil\allowbreak Is}, not only the value for the
exhibited analytic representatives.  HUF1 supplies the two
derivative clauses for (HUF.3).
\end{proof}
% END HUF3

% BEGIN HUF4
\paragraph{HUF4. Recovering the original integral CM datum from a fiber.}
\label{claim:HUF4}
For every \(y\) over the fixed \(\sigma\), the exponential
kernel of \(G_y\) carries an original datum \(D'\) of
signature \((p,q)\), and an original ordered period
\(P:L'\to\mathbb C^g\).  Its exponential has the exact
diagonal \(R\)-action, form \(B'\), tangent dictionary and
all Weil pairings in (HUF.1)--(HUF.2).

\begin{proof}
Differentiate each \(R\)-endomorphism.  Uniqueness of local
flows, followed by division by integers, gives
\[
          \exp_G(d[a]\,z)=[a]\exp_G(z).
 \tag{HUF.5}
\]
Thus the derivative action preserves the actual lattice
\(\Lambda\).  The adjoint identity of the actual line, read
in its mixed Appell--Humbert Hermitian forms, gives
\(B_G(a\lambda,\mu)=B_G(\lambda,\bar a\mu)\).
The lattice is finite over \(R\), since its finite
\(\mathbb Z\)-basis already generates it over \(R\).
It is torsion free: a nonzero \(a\in R\) acts on each Lie
eigenspace by the nonzero number \(\tau(a)\).
Over the Dedekind domain \(R\), a finite torsion-free module
is projective.  Its rank is \(r\), because its
\(\mathbb Z\)-rank is \(2g=[E:\mathbb Q]r\).

Use the given ordered type to identify the Lie algebra
complex linearly with \(\mathbb C^g\), and let \(P\) be
the image of the inclusion of \(\Lambda\).
Transport multiplication by \(i\) to \(J\) on
\(\mathbb R\otimes\Lambda\).  The form \(B_G\) is
\(J\)-invariant and \(B_G(Jv,v)=H_G(v,v)>0\) for
\(v\ne0\).  The exact \texttt{Datum.\allowbreak dom} trace identities
follow on each embedding eigenspace:
\[
 \operatorname{tr}_{\mathbb R}a
   =2\operatorname{Re}\sum_{\tau\in\Psi}
                       (p\tau(a)+q\bar\tau(a)),\qquad
 \operatorname{tr}_{\mathbb R}(Ja)
   =-2\operatorname{Im}\sum_{\tau\in\Psi}
                       (p\tau(a)+q\bar\tau(a)).
\]
Perfection is HUF3.  These verifications give every field
of \(D'\), including its original integral lattice, rather
than a lattice of finite index.  Discreteness, full real
span, positivity and the diagonal action prove
\texttt{Is\allowbreak Period}.  Equations (HUF.4)--(HUF.5) and HUF1
give the remaining stated exponential fields.
\end{proof}
% END HUF4

% BEGIN HUF5
\paragraph{HUF5. Exact comparison with the original geometric Tate marking.}
\label{claim:HUF5}
The construction of HUF4 admits an \(R\)-linear map
\(t:L\to L'\) for which
\(\texttt{In\allowbreak Genus}(D,l,\sigma,D',t)\) holds and the two
marking identities in (HUF.2) hold.  In particular all
the compatible maps \(f_n\) and multipliers \(\nu_n\)
in the original \texttt{In\allowbreak Genus} are produced.

\begin{proof}
Use the \texttt{Obj.\allowbreak geom} clause at \(y\), with coefficient
field \(\mathbb C\) and the fixed \(\sigma\).  It provides
the original \(\texttt{Tate\allowbreak Iso}\ I\), its chosen roots of
unity, and its two marking identities.
The exponential provides the second Tate description
\[
  I'_n(a)=\exp_G(Pa/n),\qquad
                      \zeta'_n=\exp(-2\pi i/n)\quad(n>0).
\]
Its kernel is \(nL'\); every \(n\)-torsion point has this
form, by surjectivity and the exact exponential kernel.
Its tower law and \(R\)-linearity follow from (HUF.5),
and its pairing is HUF3.

For each \(n>0\), comparison of the two descriptions gives
an \(R\)-linear isomorphism
\(\bar f_n:L/nL\to L'/nL'\).
The composed map \(L\to L'/nL'\) lifts to an
\(R\)-linear \(f_n:L\to L'\), because the \emph{given}
\(L\) is projective over \(R\).  This is an integral
lift through the quotient \(L'\to L'/nL'\); no choice of
an \(R\)-basis is required.  The torsion tower laws give
\(f_{nm}(a)-f_n(a)\in nL'\) for all \(n,m>0\).
Each \(\bar f_n\) is an isomorphism, giving exactly the
two source injectivity and surjectivity clauses modulo \(n\).

Choose an integer \(\nu_n\) such that
\(I.\zeta_n=\exp(-2\pi i\nu_n/n)\).  Primitivity makes
\(\nu_n\) invertible modulo \(n\).
Comparing the two Weil values, whose values are unique,
gives
\[
           B'(f_na,f_nb)-\nu_nB(a,b)\in n\mathbb Z.
\]
The specified root at \(A\) gives the required
\(\sigma(\zeta_A)=\exp(-2\pi i\nu_A/A)\).
Set \(t=f_N\).  Since \(M,A\mid N\), compatibility
gives \(f_N\equiv f_M\bmod M\) and
\(f_N\equiv f_A\bmod A\).  Substituting in the geometric
Tate marking identities gives exactly (HUF.2), on the
source's half-lattice for \(M\) and on all of \(L\)
for \(A\).  Values \(f_0=0,\nu_0=0\) are harmless:
all the relevant genus congruences have their original
positive-\(n,m\) guards.
\end{proof}
% END HUF5

% BEGIN HUF6
\paragraph{HUF6. The whole reverse-uniformization-at-a-point producer.}
\label{claim:HUF6}
For every original \(E,H,D,l,X_o,\sigma,g,\mathrm{ty}\)
in the source scope and every \(y\) over \(\sigma\),
there exist \(D',t,P,e\) such that
\[
 \texttt{In\allowbreak Genus}(D,l,\sigma,D',t)
                 \quad\hbox{and}\quad \mathcal U(D',t,P;y,e).
 \tag{HUF.6}
\]
This is the literal
\texttt{Point\allowbreak Unif.\allowbreak exists\_\allowbreak uniformisation\_\allowbreak at\_\allowbreak point\_\allowbreak closed}
conclusion, with the original moduli point unchanged.

\begin{proof}
Pull the supplied universal object back to \(y\).
HUF2 identifies its actual abelian variety and line.
HUF3 constructs the exponential with the required
holomorphic and derivative fields.
HUF4 extracts the integral datum and the period with the
given ordering of the type, and HUF5 supplies the exact
genus and both markings.  These are all conjuncts of
(HUF.6).  Neither existence of the moduli scheme nor
existence of a family uniformization has been used.
The fixed embedding \(\sigma\), including its chosen
root of unity, is the one used in HUF5.
\end{proof}
% END HUF6

\subsection{Prescribed local lifts and holomorphic dependence}

% BEGIN HUF7
\paragraph{HUF7. Relative exponentials and continuation of the full lattice.}
\label{claim:HUF7}
For a holomorphic family \(y:B\to X_{o,\sigma}^{\mathrm{an}}\)
and \(b_0\in B\), fix an exponential at \(b_0\) with its
given ordered Lie identification and an integral basis
of its full period lattice.  On a sufficiently small
connected open polydisc about \(b_0\) there is a jointly
holomorphic relative exponential, and the prescribed
period basis extends holomorphically to a basis of its
\emph{entire} kernel at every point.

\begin{proof}
Pull back the proper analytic group of HUF2 to \(B\).
Choose the local eigenframes of HUF2 to agree with the
prescribed initial derivative.  In these frames the
relative invariant vector fields depend holomorphically
on both their vector and base parameters.  The local
holomorphic differential-equation theorem gives a local
exponential near the zero section.  Define it at arbitrary
vectors by multiplication after division by a sufficiently
large positive integer, exactly as in HUF3.  On any
neighborhood of a given \((b,z)\) the same integer works.
The compatible formulas therefore give a jointly
holomorphic relative exponential on \(B'\times\mathbb C^g\).

Its derivative in the vector direction is invertible
everywhere.  The analytic implicit-function theorem,
applied to the equation \(\exp_b(\lambda)=0\) near each
of the finitely many prescribed basis vectors, continues
them as holomorphic kernel sections \(\lambda_j(b)\).
Shrink once so that their real coefficient matrix remains
invertible.  They generate a lattice \(\Lambda_b^0\).
The action of \(\mathbb Z^{2g}\) by these translations
on \(B'\times\mathbb C^g\) is properly discontinuous:
on every compact base subset the inverse real coefficient
matrices have a common bound.  The quotient
\[
       T^0=(B'\times\mathbb C^g)/\Lambda^0
                         \longrightarrow B'
\]
is a holomorphic torus family, topologically the product
with \(\mathbb R^{2g}/\mathbb Z^{2g}\), and is proper.
The exponential descends to a map \(T^0\to G_y\).
It is a local biholomorphism, since its vertical derivative
is invertible and it is the identity on the base.
On every fiber it is onto: its image is an open subgroup
of a connected compact group.

This map is proper.  Indeed the inverse image of a
compact subset is closed in the inverse image in \(T^0\)
of its compact projection to \(B'\).  A proper local
biholomorphism is a finite covering.  Its degree is
constant on the connected total space \(G_y|_{B'}\).
At \(b_0\) that degree is one, because the initial
period basis was the full kernel.  Thus it is an
isomorphism, and \(\Lambda_b^0\) is the full kernel
for every \(b\).  This proves fullness; continuation
of a finite collection of kernel elements alone would
not have sufficed.
\end{proof}
% END HUF7

% BEGIN HUF8
\paragraph{HUF8. One neighborhood preserves the form, action and both markings.}
\label{claim:HUF8}
Suppose the initial data in HUF7 satisfy
\(\mathcal U(D',t,P_0;y(b_0),e_0)\).
The continued lattice gives, on one open \(B'\ni b_0\),
\[
 P(b_0)=P_0,\qquad \texttt{Is\allowbreak Hol\allowbreak Periods}(D',\mathrm{ty},B',P),
 \qquad
 \mathcal U(D',t,P(b);y(b),e_b)\quad(b\in B').
 \tag{HUF.7}
\]
The maps \(e_b\) are jointly holomorphic.  This is the
whole original \texttt{exists\_\allowbreak local\_\allowbreak period\_\allowbreak lift}, including
preservation of its prescribed initial period.

\begin{proof}
Choose a \(\mathbb Z\)-basis of \(L'\) and continue its
images under \(P_0\) as in HUF7.  Extend by integral
linearity to obtain \(P(b)\).  The chosen eigenframes
ensure the desired diagonal \(R\)-action on the Lie
bundle.  Applying an element of \(R\) to a continued
period gives another continuous kernel section.  Its
integer coordinates in the continued basis are locally
constant, so are their original coordinates at \(b_0\).
This proves \(R\)-linearity for a finite integral basis
of \(R\), and hence for every element.  The period map
at \(b_0\) is the original \(P_0\), and uniqueness of
the exponential with a specified derivative gives
\(e_{b_0}=e_0\).

Pull the actual line back to the torus family of HUF7.
Its integral Appell--Humbert alternating form is locally
constant, because it is determined, with our fixed sign
convention, by the first Chern class on the continued
two-cycles.  One can see this
without choosing a connection: identify the underlying
torus topologically with a fixed torus on the polydisc;
the restrictions of a line bundle along two homotopic
fiber inclusions have the same integral Chern class.
The resulting alternating matrix is therefore \(B'\)
at every point.  At the initial point this follows
also from the full Weil identities: the difference
of its actual matrix and \(B'\) is divisible by
every positive integer.  Positivity of the ample line
supplies the positive Hermitian matrix at each fiber.
Thus every \(P(b)\) is precisely an \texttt{Is\allowbreak Period}.

The two markings are torsion sections, of orders dividing
\(M\) and \(A\).  Multiplication by either positive
integer on the analytic family is a local biholomorphism
with finite fibers, so its kernel is a finite covering.
Two continuous sections of this kernel that agree at
one point agree throughout a connected small base.
Apply this to a finite integral generating set of
\(\bar{\mathfrak b}L\) for the first marking, and of
\(L\) for the second.  Addition gives both equalities
for every required vector.  Finally HUF3 computes
all positive-\(n\) Weil pairings from the single
constant form \(B'\); it does not require one
neighborhood for each \(n\).  HUF1 gives the tangent
clauses.  These arguments establish all of (HUF.7)
on one neighborhood.  This proof needs no
\texttt{In\allowbreak Genus} hypothesis beyond the stated initial
\(\mathcal U\), exactly as in the source local-lift
theorem.
\end{proof}
% END HUF8

% BEGIN HUF9
\paragraph{HUF9. Local period coordinates have exactly dimension \(fpq\).}
\label{claim:HUF9}
For the original \(D'\), every ordered period \(P_0\)
has a finite-period neighborhood and a holomorphic
coordinate \(Z(P)\in\mathbb C^{fpq}\) with the following
property: if \(P,P'\) are in that neighborhood, then
\[
 Z(P)=Z(P')\ \Longrightarrow\
 \exists h\in\operatorname{GL}_g(\mathbb C):
          P'=hP,\quad h\text{ commutes with the diagonal }R\text{-action}.
 \tag{HUF.8}
\]
For every holomorphic period family in the neighborhood,
\(Z(P(b))\) is holomorphic.  This proves the exact
\texttt{Period\allowbreak Coord.\allowbreak exists\_\allowbreak period\_\allowbreak coordinates} interface.
As a standalone assertion it also includes \(p+q=0\):
then \(g=0\), the lattice is zero and the coordinate
space is a point.

\begin{proof}
Put \(W=L'\otimes_{\mathbb Z}\mathbb C\) and extend \(P\)
complex linearly to \(W\to\mathbb C^g\).
Its kernel \(Q(P)\) is a complex \(g\)-plane.
The full real span implies
\(Q(P)\oplus\overline{Q(P)}=W\).
The Riemann identities say that \(B'_\mathbb C\) vanishes
on \(Q(P)\), and positivity is an open condition.
Decompose \(W=\bigoplus_\tau W_\tau\).
Each \(W_\tau\) has dimension \(p+q\).
The adjoint relation makes \(B'_\mathbb C\) pair
\(W_\tau\) perfectly with \(W_{\bar\tau}\) and with
no other summand.

For \(\tau\in\Psi\), the kernel \(Q_\tau\) has dimension
\(q\).  Its partner \(Q_{\bar\tau}\), of dimension \(p\),
is its annihilator under this perfect pairing.
Thus \(Q\) is specified by the \(q\)-planes \(Q_\tau\)
for \(\tau\in\Psi\), subject only to the open
real-complementarity and positivity conditions.
It lies in an open subset of
\(\prod_{\tau\in\Psi}\operatorname{Gr}_q(\mathbb C^{p+q})\),
of complex dimension \(fpq\).
Choose complements at \(Q(P_0)\).  Nearby kernels are
graphs of \(q\)-by-\(p\) matrices, which give \(Z\).
Their entries are quotients of matrix minors of \(P\)
with fixed nonzero denominators at \(P_0\); hence they
are holomorphic functions of finitely many period
coordinates.

Choose a finite integral basis of \(L'\).
Continuity of those finitely many minors gives an
\(\varepsilon>0\) such that
\(\|P(a)-P_0(a)\|<\varepsilon\) on that basis places
\(P\) in the graph chart.  This supplies the finite
set and the common radius in the literal source
statement, for every parameter dimension \(m\).
Equal graph coordinates mean equal kernels.  The
two quotient isomorphisms \(W/Q\to\mathbb C^g\)
then differ by a unique \(h\), proving (HUF.8).
Both maps are \(R\)-linear, so \(h\) commutes with
the prescribed action.  If \(p=0\) or \(q=0\),
these Grassmannians are points and the same argument
gives the required zero-dimensional coordinate.
\end{proof}
% END HUF9

% BEGIN HUF10
\paragraph{HUF10. A marked exponential is unique.}
\label{claim:HUF10}
Suppose \(\texttt{In\allowbreak Genus}(D,l,\sigma,D',t)\) and
both \(\mathcal U(D',t,P;y,e)\) and
\(\mathcal U(D',t,P;y,e')\) hold.  Then \(e=e'\).
This is the literal \texttt{exp\_\allowbreak eq\_\allowbreak of\_\allowbreak periods\_\allowbreak of\_\allowbreak marking}.

\begin{proof}
By HUF1 and the derivative clauses, both maps induce
analytic isomorphisms \(\mathbb C^g/P(L')\to G_y\).
Their quotient gives an analytic group automorphism
of this torus.  Lifting to its simply connected
covering and differentiating gives a complex-linear
automorphism \(h\) and an integral lattice
automorphism \(\gamma\), with \(P\gamma=hP\).
The \(R\)-action implies that \(\gamma\) is \(R\)-linear.
Comparison of all the Weil identities gives
\(B'(\gamma a,\gamma b)-B'(a,b)\in n\mathbb Z\)
for every \(n>0\).  Hence \(\gamma\) is an exact
isometry.

The level identities give
\(\gamma(ta)-ta\in ML'\) for
\(a\in\bar{\mathfrak b}L\), and
\(\gamma(ta)-ta\in AL'\) for all \(a\in L\).
The genus map \(t\) is onto modulo \(N\).
Writing an arbitrary element of
\(\bar{\mathfrak b}L'\) as a sum of ideal multiples
and using this surjectivity shows that it is congruent
modulo \(N\bar{\mathfrak b}L'\) to \(ta\) for
some \(a\in\bar{\mathfrak b}L\).  The half-level
congruence therefore holds on that whole half-lattice.
PUN6 turns it into \(\gamma\equiv1\pmod M\)
on all of \(L'\).

The positive Hermitian form is determined by its
imaginary part and the complex structure.  Thus
\(\gamma\) preserves the positive real norm as well.
An integral isometry group of a positive norm is
finite: images of a fixed integral basis lie in
finite intersections of a lattice with bounded
sets.  Consequently \(\gamma\) has finite order.
PUN8's principal-congruence argument, for the
original odd \(M\geq3\), forces \(\gamma=1\).
Full real span then gives \(h=1\), and \(e=e'\).
No prime \(3\) exception is used.
\end{proof}
% END HUF10

% BEGIN HUF11
\paragraph{HUF11. The period coordinate on a moduli chart is a local biholomorphism.}
\label{claim:HUF11}
At any \(y_0\) with a chosen \(\mathcal U(D',t,P_0;y_0,e_0)\),
there is an analytic chart of the original moduli space
on which the coordinate \(Z\) of its continued periods
is a local biholomorphism to an open subset of
\(\mathbb C^d\).

\begin{proof}
The supplied smoothness field and HUF1 give a
holomorphic injective chart \(x:U\subseteq\mathbb C^d
\to X_{o,\sigma}^{\mathrm{an}}\), with \(x(0)=y_0\).
Apply HUF8 to it, with the prescribed \(P_0,e_0\),
and shrink into the finite-period neighborhood of
HUF9.  The map \(z\mapsto Z(P_z)\) is holomorphic.
If two of its values coincide, HUF9 gives a
complex-linear \(h\) taking one period to the other.
PUN23, with integral isometry \(\gamma=1\) and
the two identical marking maps \(t\), gives equality
of the two moduli points.  The chart is injective,
so the two parameters are equal.

An injective holomorphic map between open subsets
of \(\mathbb C^d\) has invertible derivative at
every point and is a biholomorphism onto its image.
We use exactly this classical theorem, as stated
and proved in Lebl, \emph{Tasty Bits of Several
Complex Variables}, Theorem 1.6.6, pp.\ 43--45.
Its equal-dimension hypothesis holds because
HUF9 gives \(fpq=D.\mathrm{dim}\), the dimension
of the supplied smooth moduli chart.
The holomorphic inverse-function theorem gives
the asserted inverse on an open image.
For \(d=0\), both charts are a single point
and the statement is immediate.
\end{proof}
% END HUF11

% BEGIN HUF12
\paragraph{HUF12. Arbitrarily chosen pointwise uniformizations vary holomorphically.}
\label{claim:HUF12}
Retain \(\texttt{In\allowbreak Genus}(D,l,\sigma,D',t)\).
Let \(B\subseteq\mathbb C^m\) be any open set and
\(P(b)\) any original \texttt{Is\allowbreak Hol\allowbreak Periods} family.
Suppose arbitrary functions \(y(b),e_b\) satisfy
\(\mathcal U(D',t,P(b);y(b),e_b)\) for \(b\in B\).
Then
\[
 \texttt{Is\allowbreak Hol\allowbreak To}(B,y),\qquad
 \texttt{Is\allowbreak Hol\allowbreak To}(B\times\mathbb C^g,(b,z)\mapsto e_b(z)).
 \tag{HUF.9}
\]
These are the full original
\texttt{is\allowbreak Hol\allowbreak To\_\allowbreak modulipoint\_\allowbreak exp\_\allowbreak of\_\allowbreak holomorphic\_\allowbreak periods}
conclusions; continuity of the chosen functions is not
an added hypothesis.

\begin{proof}
Fix \(b_0\).  Use \(y(b_0),P(b_0),e_{b_0}\) in HUF11.
For \(b\) close enough to \(b_0\), the holomorphic
coordinate \(Z(P(b))\) lies in the image of that chart.
Its holomorphic inverse produces a moduli point
\(\widetilde y(b)\) and the chart's jointly
holomorphic periods \(\widetilde P(b)\) and exponential
\(\widetilde e_b\).  Equal \(Z\)-coordinates give
an \(R\)-equivariant complex-linear isomorphism
\(h(b)\) with
\(\widetilde P(b)=h(b)P(b)\).
This matrix is holomorphic: choose \(g\) period
vectors independent over \(\mathbb C\) at \(b_0\);
locally it is the matrix of their tilded columns
times the inverse matrix of their untilded columns.

PUN23, with \(\gamma=1\), gives
\(\widetilde y(b)=y(b)\) for each \(b\), without
first requiring that \(y\) be continuous.
Moreover \(\widetilde e_b\circ h(b)\) and \(e_b\)
have the same complete \(\mathcal U\) data:
the period, action, kernel, tangent, Weil and
marking formulas transport directly under \(h(b)\).
HUF10 therefore identifies them.  The former is
jointly holomorphic in \((b,z)\), so the latter is
as well.  HUF1 gives the precise source predicates.
Doing this at every \(b_0\) proves (HUF.9) on
the whole possibly disconnected open \(B\).
\end{proof}
% END HUF12

\subsection{Finite disjoint families on the supplied moduli scheme}

% BEGIN HUF13
\paragraph{HUF13. Finiteness of the analytic component set.}
\label{claim:HUF13}
The analytic space \(X_{o,\sigma}^{\mathrm{an}}\) has
finitely many connected components.  In particular every
locally constant map from it to a discrete set has finite
image.

\begin{proof}
The supplied quasiprojectivity gives finite type and
quasicompactness, so choose a finite affine cover after
base change to \(\mathbb C\).  Each affine has a closed
immersion into some \(\mathbb C^a\).  In real and
imaginary coordinates its complex points are a real
semialgebraic set: a complex polynomial equation is
two real polynomial equations, with arbitrary real
coefficients allowed.  The finite connected-component
theorem for semialgebraic sets applies.
The finitely many connected pieces of the finitely
many affines cover \(X_{o,\sigma}^{\mathrm{an}}\).
Every such piece lies in one connected component,
so the number of components of the whole space is
at most the number of pieces.  On each connected
component a locally constant map to a discrete set
is constant, proving the assertion.

The named input is Coste, \emph{An Introduction to
Semialgebraic Geometry} (October 2002), Theorem 2.23,
p.\ 40, following the finite cylindrical decomposition
Theorem 2.20.  We use only finiteness of connected
components.  We do not assume that the moduli space
itself is compact, or that its period map is proper.
\end{proof}
% END HUF13

% BEGIN HUF14
\paragraph{HUF14. The exact marked-lattice classes have disjoint open images.}
\label{claim:HUF14}
Consider all pairs \((D',t)\) in the original genus.
Declare \((D',t)\sim(D'',t')\) if there is an
\(R\)-linear isometry \(\gamma:L'\simeq L''\) with
\[
 \begin{split}
 \gamma(ta)-t'a&\in ML''\quad(a\in\bar{\mathfrak b}L),\\
 \gamma(ta)-t'a&\in AL''\quad(a\in L).
 \end{split}                                                     \tag{HUF.10}
\]
The images in \(X_{o,\sigma}(\mathbb C)\) of all periods
for two such pairs are equal if the pairs are equivalent,
and disjoint otherwise.  Each image is nonempty and open,
and together they cover all of \(X_{o,\sigma}(\mathbb C)\).

\begin{proof}
Composition and inversion of \(R\)-linear integral
isometries preserve both congruences in (HUF.10);
thus this is an equivalence relation.
Every datum has periods of the given ordered type
by PUN15.  PUN22 constructs a point for each such
period and genus marking, so each image is nonempty.

If two images meet, apply PUN23 to the two
uniformizations of that same point.  Its integral
isometry has exactly the congruences (HUF.10).
Conversely, given \(\gamma\) as there, a period
\(P\) for \(D'\) transports to \(P\gamma^{-1}\)
for \(D''\).  It has the same exact pairing,
diagonal action, discreteness, fullness and positivity.
PUN22 constructs a point for it, and PUN23 with
complex map \(h=1\) identifies this point with the
one for \(P\).  Thus equivalent pairs have identical
images, not merely an intersection.

At a point in any image, take a small analytic
moduli chart.  HUF8, starting from that point's
uniformization, gives periods on the \emph{same}
\(D',t\) for every nearby point.  The image is
therefore open.  HUF6 proves that all the images
cover the moduli space.  They form a disjoint open
partition, hence define a locally constant class
map to the discrete set of classes.
There is no set-theoretic choice over a proper
class hidden here: one may transport each free
\(\mathbb Z\)-module \(L'\) to
\(\mathbb Z^{2g}\), and then take the set of
its integral actions, forms and marking matrices.
Transporting back gives the original arbitrary
carriers and loses no pair.
\end{proof}
% END HUF14

% BEGIN HUF15
\paragraph{HUF15. One finite positive family, chosen before every period.}
\label{claim:HUF15}
There are \(n>0\) original data \(D_i\) and genus
maps \(t_i:L\to L_i\), one for each of the classes
in HUF14, and total functions
\[
 u_i:(L_i\to^+\mathbb C^g)\to X_o(\mathbb C),\qquad
 e_{i,P}:\mathbb C^g\to G_{u_i(P)}(\mathbb C)
 \tag{HUF.11}
\]
such that \(\mathcal U(D_i,t_i,P;u_i(P),e_{i,P})\)
for every valid period.  These same functions have
the original \texttt{onto}, \texttt{disj},
\texttt{hol} and \texttt{ex\_\allowbreak diff} fields.

\begin{proof}
HUF13 and the locally constant class map of HUF14
make the class image finite.  PUN24 supplies an
actual point over the prescribed \(\sigma\);
HUF6 places it in that image, so its cardinality
\(n\) is positive.
Choose one \((D_i,t_i)\) for each distinct class.
The genus conditions give their correct signatures
and both injectivity and surjectivity of \(t_i\)
modulo \(N\), as required in \texttt{t\_\allowbreak inj}
and \texttt{t\_\allowbreak surj}.

For every valid \(P\), choose the point and
exponential of PUN22.  For invalid additive maps,
assign a fixed point of the same nonempty class
and the constant identity-valued map.
This defines the source's total functions on
\emph{all} additive maps.  Their axioms concern
valid periods only, so these extra values impose
no condition.  HUF14 gives both coverage and
disjointness for these choices.  To see coverage
literally, transport the period supplied by HUF6
through the integral isometry to the selected
representative; PUN23 then gives \(u_i(P)=y\).

For any open \(B\subseteq\mathbb C^m\) and any
holomorphic period family, the chosen pointwise
functions satisfy all the hypotheses of HUF12.
Thus they have the \emph{same} holomorphic maps
and joint exponential maps required by
\texttt{Family\allowbreak Unif.\allowbreak hol}; no separate choice for
each base or family is made.  Their derivative
clauses already belong to \(\mathcal U\), by HUF1.
The remaining pointwise exponential, action,
Weil and marking fields are likewise exactly
those of \(\mathcal U\).
\end{proof}
% END HUF15

% BEGIN HUF16
\paragraph{HUF16. The two exact remaining fiber and lift fields.}
\label{claim:HUF16}
For the single family chosen in HUF15, the
original \texttt{Family\allowbreak Unif.\allowbreak fibre} and
\texttt{Family\allowbreak Unif.\allowbreak lift} fields hold with their
printed half-lattice and full-\(A\) congruences.

\begin{proof}
PUN23 says that \(u_i(P)=u_i(P')\) is equivalent
to an \(R\)-linear isometry \(\gamma:L_i\to L_i\)
and a complex-linear \(h\) with
\(P'\gamma=hP\), and with
\(\gamma t_i a-t_i a\in ML_i\) on
\(\bar{\mathfrak b}L\) and in \(AL_i\) on
all of \(L\).
Surjectivity of \(t_i\) modulo \(N\) transfers
these statements to the exact source lattices.
For clarity, an element of
\(\bar{\mathfrak b}L_i\) is a finite sum
\(\sum_j b_jx_j\) with \(b_j\in\bar{\mathfrak b}\).
Write \(x_j=t_i a_j+N z_j\).
Then this element equals \(t_i(\sum_j b_ja_j)\)
modulo \(N\bar{\mathfrak b}L_i\).
Since \(M\mid N\), its displacement under
\(\gamma\) is zero modulo \(M\).
The analogous argument for all of \(L_i\)
uses \(A\mid N\).  Conversely \(t_i\) maps
each required source lattice into the indicated
target lattice, so the target congruences imply
the original congruences.  This proves both
directions of the literal \texttt{fibre} field.

For the lift field, take any holomorphic
\(y:B\to X_{o,\sigma}\) and \(b_0\in B\).
HUF15 gives \(i,P_0\) with \(u_i(P_0)=y(b_0)\),
together with its chosen exponential.
HUF8 continues these on an open
\(B'\ni b_0\), \(B'\subseteq B\).
At each \(b\), that continued uniformization
and the selected one at \(P(b)\) have identical
period and marked-lattice data.  PUN23 with
\(\gamma=h=1\) gives \(u_i(P(b))=y(b)\).
All the original open-set, base-point and
holomorphic-period requirements are retained.
\end{proof}
% END HUF16

\subsection{Simultaneous normalized theta frames}

% BEGIN HUF17
\paragraph{HUF17. Normalized periods force the original theta coordinates to be holomorphic.}
\label{claim:HUF17}
Assume the exact hypotheses of
\texttt{Family\allowbreak Unif.\allowbreak theta}, with an arbitrary supplied
integral symplectic basis, sign unit \(s\), and
matrices \(h(b),\tau(b)\).
On the nonempty open base they necessarily satisfy
\[
       s=-1,\qquad \tau=\tau^t,\qquad
       h(b)=A(b)^{-1},\qquad \tau(b)=A(b)^{-1}C(b),
 \tag{HUF.12}
\]
where \(A\) and \(C\) are the columns of the two
halves of the supplied period basis.
In particular both matrices are holomorphic,
although their holomorphy is not assumed in the
source theta field.

\begin{proof}
The source equations are \(hA=1\) and
\(hC=\tau\).  These are square \(g\)-by-\(g\)
matrices, so \(A,h\) are invertible and the
last two formulas follow.  Their entries are
holomorphic because the period columns are,
and matrix inversion is holomorphic where the
determinant is nonzero.

In the coordinates \(w=hz\), the positive
Hermitian form has a matrix \(S\).
Its imaginary part vanishes on the first half
of the integral basis, so \(S\) is real symmetric
and positive definite.  Set \(\tau=X+iY\).
The mixed pairing on that basis gives
\(-SY=s\,1\).  The stated positive definiteness
of \(Y\), and that of \(S\), imply \(s=-1\)
and \(S=Y^{-1}\); here \(g\geq1\), since the
original \(r\geq1\).  The pairing vanishes on
the second half too.  Its imaginary matrix is
\(\operatorname{Im}(\tau^tS\bar\tau)=X-X^t\),
so \(X=X^t\).  As \(Y=Y^t\), this proves
\(\tau=\tau^t\).
The sign is therefore the one appropriate to
our first-variable-linear Hermitian convention:
\(B'(a_i,b_j)=-\delta_{ij}\).
When the base is empty the desired theta
assertion has no conditions there, and needs
no sign conclusion.
\end{proof}
% END HUF17

% BEGIN HUF18
\paragraph{HUF18. One constant half-characteristic and its exact theta line.}
\label{claim:HUF18}
Under HUF17 there are constants
\(c_1,c_2\in\{0,\tfrac12\}^g\), depending on
the original marked lattice and the given
symplectic basis, such that the holomorphic
theta line on the period-torus family has
the same rigidified line as the actual
universal primitive line on each fiber.
Its section on the covering space is precisely
the source series
\(\operatorname{thetaChar}(\tau(b),c_1,c_2,h(b)z)\).

\begin{proof}
Write the symplectic basis as \(a_i,b_i\), and
use \(w=h(b)z\) as the covering coordinate.
The period lattice is
\(\mathbb Z^g+\tau(b)\mathbb Z^g\).
On the fixed integral lattice the primitive
identity, by the explicit Appell--Humbert
calculation in PUN20, forces
\[
          \alpha(x)=\exp(\pi iB'(x,w_0x))\in\{1,-1\}.
 \tag{HUF.13}
\]
This conclusion applies to the actual line:
its alternating form is \(B'\), since its
Weil formulas agree with (HUF.2) for every
positive integer.  The difference of two
integral forms with those values is divisible
by every such integer and hence zero.
The Hermitian form is then determined by
\(B'\) and the complex structure.

Choose \(c_{1i},c_{2i}\) in the stated two-element
set so that
\(\alpha(a_i)=e^{2\pi i c_{1i}}\) and
\(\alpha(b_i)=e^{-2\pi i c_{2i}}\).
The semicharacter law and the sign in HUF17
give, for \(m,n\in\mathbb Z^g\),
\[
 \alpha(m+\tau n)=
       (-1)^{m^tn}\exp(2\pi i(c_1^tm-c_2^tn)).
 \tag{HUF.14}
\]
All these signs are fixed by the lattice data;
none is chosen separately at different base points.

The holomorphic factors
\[
 j_\theta(m+\tau n,w)=
 \exp\!\left(2\pi i c_1^tm-\pi i n^t\tau n
                  -2\pi i n^t(w+c_2)\right)
 \tag{HUF.15}
\]
obey the lattice cocycle identity, by direct
multiplication and the integrality of the
cross terms.  They therefore define a
holomorphic line on the torus family, with
its canonical rigidification at \(w=0\).
On a fixed fiber put \(Y=\operatorname{Im}\tau\)
and \(H(w,v)=w^tY^{-1}\bar v\).
Multiplying the Appell--Humbert frame by the
quadratic gauge
\(\exp(-\pi w^tY^{-1}w/2)\) changes its factors
to (HUF.15).  Indeed the factor quotient is
\[
 \alpha(m+\tau n)\,
 \exp(-2\pi i n^tw-\pi i m^tn-\pi i n^t\tau n),
 \]
which is exactly (HUF.15) by (HUF.14).
This checks equality of the line itself and
its rigidification, not only of its Chern
class.  The gauge is used on a single fiber;
no holomorphy of \(Y\) in the base is asserted.

Reindexing the source theta series under
\(w\mapsto w+m+\tau n\) gives (HUF.15).
It is jointly holomorphic: on every compact
subset the smallest eigenvalue of \(Y\) is
bounded below by a positive constant, and
its summands and all differentiated summands
are bounded by a polynomial times
\(\exp(-c\|k+c_1\|^2+C\|k+c_1\|)\).
The resulting normally convergent series
therefore gives a holomorphic section of
this exact line.  It is not identically
zero on any fiber, as is also seen from
its nonzero Fourier coefficients.
\end{proof}
% END HUF18

% BEGIN HUF19
\paragraph{HUF19. Rigidified fiber isomorphisms glue over an arbitrary open base.}
\label{claim:HUF19}
Let \(\pi:T\to B\) be a proper holomorphic
submersion with connected compact fibers and
a zero section.  If two rigidified holomorphic
lines on \(T\) are isomorphic on every fiber,
then they have a unique global holomorphic
isomorphism respecting the rigidifications.
This applies to the theta line and actual
line of HUF18 on every original open
\(B\subseteq\mathbb C^m\).

\begin{proof}
Let \(Q\) be the Hom line between them.
On each fiber it is trivial, and its space
of holomorphic sections has dimension one,
because holomorphic functions on a compact
connected complex manifold are constant.
The line \(Q\) is flat over \(B\): locally
the submersion is a projection of complex
coordinate polydiscs, and \(Q\) is locally
free of rank one there.

Apply Grauert's constant-cohomology theorem
with degree zero.  It says that
\(E=\pi_*Q\) is a holomorphic line bundle
and that \(E\otimes\mathbb C(b)\simeq
H^0(T_b,Q_b)\) for every \(b\).
The evaluation map \(\pi^*E\to Q\) is
an isomorphism: on each fiber it evaluates
the constants in a trivial line and is
nonzero at every point.  Evaluating also
on the zero section gives an isomorphism
\(E\to e^*Q\).  The two rigidifications
identify \(e^*Q\) with \(\mathcal O_B\).
Their unit section therefore defines a
global nowhere-zero section of \(Q\),
normalized at zero.  It is the required
line isomorphism.  Another such isomorphism
has a ratio constant on each fiber and
equal to one at its zero point, so is
identical.

The precise named input is Grauert,
\emph{Ein Theorem der analytischen
Garbentheorie und die Modulr\"aume komplexer
Strukturen}, Publ.\ Math.\ IHES 5 (1960),
\S7, Satz 5, pp.\ 62--64, with the 1963
corrigendum in volume 16, pp.\ 35--36.
Its hypotheses here are a proper map,
a reduced base, a base-flat coherent sheaf
and constant \(h^0=1\), all just checked.
It gives local freeness, not a gratuitous
global trivialization: the latter is
provided by the rigidification argument
above.  No Stein, contractibility, or
simple-connectivity condition on \(B\)
is used.
\end{proof}
% END HUF19

% BEGIN HUF20
\paragraph{HUF20. The full simultaneous normalized theta-frame producer.}
\label{claim:HUF20}
For every open preconnected \(B\subseteq\mathbb C^m\),
every holomorphic period family, every original
symplectic basis and sign, and every \(h,\tau\)
in the printed theta hypotheses, suppose the
arbitrarily supplied \(y,e\) satisfy pointwise
\(\mathcal U\), \(\texttt{Is\allowbreak Hol\allowbreak To}(B,y)\), and
the printed joint holomorphy condition for \(e\).
Then the constants
of HUF18 satisfy the exact
\(\texttt{Is\allowbreak Theta\allowbreak Frame}\) predicate for the
actual universal rigidified line.
Thus the whole original
\texttt{exists\_\allowbreak theta\allowbreak Frame\_\allowbreak family} conclusion
is supplied.

\begin{proof}
The given joint exponential map (supplied by HUF12
when applied to the chosen family) identifies the holomorphic
torus family with the analytic pullback of
the universal abelian scheme.  Indeed it
descends through the exact lattice, and
on each fiber is a bijective local
biholomorphism; its total differential is
invertible because it is the identity on
the base.  Thus the descended map is an
analytic isomorphism.  The torus family is
proper: its underlying real torus bundles
are locally products with a compact torus,
using a real basis of the holomorphic
period lattice.

HUF17 derives the holomorphic coordinates,
and HUF18 constructs the theta line and
its constants.  Pull back the actual
universal line to the torus family.
HUF19 gives the unique global rigidified
holomorphic isomorphism from the theta
line to it.  On the universal covering
\(B\times\mathbb C^g\), the theta line
has its canonical nonzero frame.
Transport this frame by that isomorphism,
and write its coefficient \(fr_j(b,z)\)
in each of the source's algebraic bundle
charts.

These coefficients are holomorphic on
the exact inverse images of the chart
opens, nowhere zero there, and obey
\[
 fr_j(b,z)=g_{jj'}(e_b(z))\,fr_{j'}(b,z).
 \tag{HUF.16}
\]
Because the isomorphism is rigidified,
their values at \(z=0\) equal the
printed rigidification section in
those same charts.  These verify
the first four \texttt{Is\allowbreak Theta\allowbreak Frame}
clauses with their actual transition
orientation.

The theta section of HUF18 transports
to a holomorphic section of the actual
line on each compact projective fiber.
Projective GAGA makes it an algebraic
section on that fiber.  In every
pullback chart its value is exactly
\(\operatorname{thetaChar}(\tau(b),c_1,c_2,h(b)z)
fr_j(b,z)\); this is the last clause,
including every point of the specified
fiber product.  That clause asks for
an algebraic section separately on
each fiber, not for algebraization
of the entire analytic family.
We have therefore proved exactly its
scope.  For \(B=\varnothing\), take
both characteristics to be zero and
arbitrary frame functions; all clauses
are vacuous.  We retain the source's
preconnectedness hypothesis throughout.
\end{proof}
% END HUF20

\subsection{Installation of the literal finite-family theorem}

% BEGIN HUF21
\paragraph{HUF21. The whole original \texttt{N1c} and \texttt{Family\allowbreak Unif}.}
\label{claim:HUF21}
For every original \(E,H,D,l,X_o,\sigma,g,\mathrm{ty}\)
with \(p+q\geq1\), \(\sigma(\rho)=\mathrm{rt}_0\),
and the given \texttt{Is\allowbreak Type\allowbreak Fn}, there exists a
\[
                       \texttt{Family\allowbreak Unif}\ X_o\ \sigma\ \mathrm{ty}.
 \tag{HUF.17}
\]
This is the literal \texttt{N1c} conclusion.
It retains the supplied \texttt{Moduli} record
and the named classical foundations stated above;
no individual-point or supplied-family hypothesis
remains in place of the family conclusion.

\begin{proof}
Take the finite positive family and the single
total maps \(u_i,e_{i,P}\) of HUF15.
Their genus records supply \texttt{sig},
\texttt{t\_\allowbreak inj} and \texttt{t\_\allowbreak surj}.
The complete pointwise \(\mathcal U\) supplies
\texttt{over\_\allowbreak \(\sigma\)}, \texttt{ex\_\allowbreak add},
\texttt{ex\_\allowbreak onto}, \texttt{ex\_\allowbreak ker},
\texttt{ex\_\allowbreak lin}, \texttt{ex\_\allowbreak weil},
\texttt{ex\_\allowbreak eta}, \texttt{ex\_\allowbreak eta\allowbreak A}
and both \texttt{ex\_\allowbreak diff} clauses.
HUF15 supplies \texttt{onto}, \texttt{disj}
and the universal \texttt{hol} field.
HUF16 supplies the exact \texttt{fibre}
and \texttt{lift} fields.
HUF20 supplies \texttt{theta}, for all
parameter dimensions, all open preconnected
bases, every prescribed symplectic basis
and every pair of normalizing matrices
allowed by the original statement.
These exhaust the source record.

Every choice of components and total maps
was made once, before the universal
period, base, matrix, basis and torsion
quantifiers.  All Weil assertions use
every positive integer with the original
negative sign.  The lattice maps are
integral modulo the exact \(N\);
the half-level and full-\(A\) conditions
have not been replaced by rational
or finite-index substitutes.
The case \(d=0\) was included in HUF9
and HUF11, and arbitrary \(m\), including
zero, is allowed throughout.
Thus (HUF.17) has precisely the original
quantifier order and scope.
\end{proof}
% END HUF21

\paragraph{Present boundary.}
HUF21 supplies the finite holomorphic family
that PUN11 previously had to assume.
It does not itself construct the representing
\texttt{Moduli} scheme.
The adjoining principal-chart theorem additionally
requires its own simultaneous normalized framed
points, the exact \texttt{Is\allowbreak Unf} dictionary and
the reduced-test-ring \texttt{From\allowbreak Test} assertion.
Those are separate clauses, not consequences
of merely naming HUF21; their installation is
kept distinct from the family endpoint here.


\section{Compatibility on the actual built family and the full local gain}
\label{sec:witness-realization-v049}

\noindent\textbf{Scope and conclusion.}
The five-field compatibility record constructed in GCF12 and transported
in GCF16 has to concern the comparison map in the supplied arithmetic
object.  An independently chosen compatible comparison does not suffice.
This chapter obtains the required identification directly from the
original built-point class.  Its construction record already contains
one joint equality of the group ties, comparison, twisted forms and
kernel.  The final point changes the division lattice and its height;
we prove that the same joint equality survives that change.  This gives
the original \(\mathrm{GroupCompatibilityFamily}\) on every member of
the original built class, before imposing the selected-design conditions.
Together with the selected tops already treated in GCF18 and the full
local estimates of LGS13, it supplies the five original
\(\mathrm{LinkG}\) conclusions with their original class quantifiers.

The argument does not construct a nonempty design class.  Nor does it
construct the stronger \(\mathrm{RecT}\), its row pin, host record or
Weil-equivariance field.  Those objects are required elsewhere in the
proof.  They are unnecessary for this direct compatibility supplier.
An arbitrary family outside the built class still needs an actual
compatibility input.  A separate final statement also proves the
universal original \texttt{Global\allowbreak Theory.\allowbreak Top.\allowbreak link\_\allowbreak g} implication, at its
literal \(\mathrm{TSel}\) and duty premises.

\paragraph{Objects and notation.}
The original \(D\) contains the number field \(L=D.\mathrm{LL}\), its
specified complex embedding, the host, words and places, the column
spaces, their integral models and comparison domains.  A full object
\(O:\mathrm{Obj}(D)\) includes its size data, actual group instance,
local lattices and \(O.\mathrm{ex}:\mathrm{Extra}\).
The smaller object \(\mathrm{ofObj}(O):\mathrm{Obj2}(D)\) retains
\[
 \begin{gathered}
 G_t,\ \mathrm{th},\ c,\ \mathrm{tw},\ G_y,\
 \mathrm{errX},\ \mathrm{nch},\ \mathrm{xd},\ \mathrm{lev},\
 X,\ p_0,\ \mathrm{ker},\ \mathrm{amb},\ \mathrm{amb}_0,\
 \mathrm{gauge}.
 \end{gathered}
\]
In particular it retains the actual \(\mathrm{Extra}.\mathrm{cmp}\)
and \(\mathrm{Extra}.\mathrm{ker}\); forgetting size fields does not
forget either of them.  A prime on an object or map below means the
original \texttt{set\allowbreak Div} transport.  It never denotes an independent
choice of a local model or a change of coefficient field.

% BEGIN WRS1
\paragraph{WRS1. Division transport preserves the whole local kernel record.}
\label{claim:WRS1}
Let \(N_u\) be arbitrary \(\mathcal O_L\)-submodules of the original
division vector spaces and let \(h_u\in\mathbb R\) be arbitrary.
The original record
\(D'=\mathrm{setDiv}(D,N,h)\) replaces
\[
 \mathrm{Nu}(u).L\ \text{by}\ N_u,\qquad
 \mathrm{div}.h\ \text{by}\ h,\qquad
 \mathrm{HH}\ \text{by}\
 \max\{1,\mathrm{sFrame}(\mathrm{Nu}[L:=N],h)+\mathrm{RR}.\mu_{\max}\}.
 \tag{WRS.1}
\]
All other fields are retained.  For every supplied
\[
 k_{u,w}:\mathrm{LocalToGlobal}_2
       (D,G_t,c,\mathrm{tw},u,w;G_y.\mathrm{Ch}(u,w)),
\]
there is the original transported record on
\((D',G_t',c',\mathrm{tw}',G_y')\), with precisely the same
coefficient embeddings, local lattice maps, class maps, Taylor series
and comparisons.  This holds for every allowed \(u,w\) and the
original good-index guard.  There is no condition that the new
\(N_u\) be full, finitely generated or nonzero.

Here is a field-by-field verification.  Neither the residue prime, the
local Frobenius period, the chain length nor the local filling and level
data changes under (WRS.1).  The two identities identifying these
lengths are therefore unchanged.  Keep the original maps
\(\lambda:\overline{\mathbb Q}\to\mathbb C_p\),
\(i_K:K_f\to\mathbb C_p\), and all
\(i_{w,j,v}:E\to Q_{j,v}\).
Their coefficient-embedding identity remains
\[
 i_K(c.\mathrm{embW}(z))=\lambda(\mathrm{emb}_L z),
 \tag{WRS.2}
\]
with the same algebraicity guard on the right.
Keep the original global-lattice map, its local-span identity, its
class map, the equivalence relating equality of classes to the local
level action, and its surjectivity onto every class at every
\(j\le m\).  Those statements use the twisted column lattices and
local level tuples; they do not use the division submodules replaced
in (WRS.1).

Retain the original Taylor series at every adjacent pair of local
classes.  Finally keep every semilinear comparison \(k.\mathrm{cmp}_j\),
its injectivity for \(j\le m\), and the three exact laws
\[
 \begin{aligned}
 k.\mathrm{cmp}_0(\mathrm{tw}.e_0 x)&=c.\mathrm{cmp}(x),\\
 \mathcal O_{K_f}\operatorname{span}
   \bigl(k.\mathrm{cmp}_0(\mathrm{tw}.\mathrm{lat}_0)\bigr)
     &=\mathrm{Lv.Int}_0,\\
 \mathrm{emb}_{j+1}y=\mathcal K_j(\mathrm{emb}_j x)
   &\ \Longrightarrow\
 k.\mathrm{cmp}_{j+1}y=\mathrm{Ch}.K_{z,j}(k.\mathrm{cmp}_j x)
       \qquad(j<m).
 \end{aligned}                                                    \tag{WRS.3}
\]
All expressions in (WRS.2)--(WRS.3) are the same expressions on the
same vectors before and after transport.  This checks all seventeen
original fields of \(\mathrm{LocalToGlobal}_2\), including the last
chain-compatibility field, not just the underlying vector-space map.
Taking these records for all \(u,w\) gives the original
\texttt{ker\allowbreak Tr}.  Transporting a full \(\mathrm{Extra}\) or an
\(\mathrm{Obj2}\) uses this very record at its kernel field.
In particular no auxiliary zero-chain substitution is used here.
\hfill\(\square\)
% END WRS1

% BEGIN WRS2
\paragraph{WRS2. Taylor presentation and every kernel tie survive division.}
\label{claim:WRS2}
Fix original true-column data \(w,Z\), a column record
\(t:\mathrm{ColT}(w,Z)\), a divisor record \(X\), a local lattice
record \(G_y\), and
\[
 k:\mathrm{KerFam}(w.DZ,t.G_t,t.c,w.\mathrm{tw}(Z),G_y),\qquad
 \mathcal K:\mathrm{KerT}(t,X,G_y,k).
\]
For every \(N,h\) as in WRS1 there is a full original
\[
 \mathcal K':\mathrm{KerT}
       (t',X',G_y',k'),                              \tag{WRS.4}
\]
where the column transport is the one of GCF16 and the local kernel
transport is the one of WRS1.  If \(\mathcal K\) satisfies the
original \(\mathrm{EntC}\) predicate, then \(\mathcal K'\) does too,
with the same place embeddings and the same basis matrices.

The nine fields of \(\mathrm{KerT}\) can all be retained.  The
lifting objects \(\mathcal O_{u,\mu}\) have the same Hermitian
spaces, frames, theta functions, lattices and transformation rules.
The standardness assertion uses \(X.Q_p\), which is retained by
\(X'=\mathrm{divTr}(X)\).  All the unfolding assertions at
\(i<m\) therefore concern the same functions.  Keep the global
kernel presentation \(E\).  Its compatibility \(\mathrm{EntOk}\)
with the lifting objects keeps the same witnesses at each tuple,
place and level; the transported objects have those same fields.
The equality \(k_{u,w}.\lambda=E.\lambda_w\), the complete served
Taylor-series clause, the symmetric kernel embedding equation, and
the positive-place equation retain identical maps on identical
arguments.  This supplies, respectively, the fields
\(\mathrm{O},\mathrm{std},\mathrm{unf},\mathrm{E},\mathrm{ok},
\mathrm{lam},\mathrm{served},\mathrm{iw},\mathrm{iw\_\allowbreak pos}\).

For clarity, preserving the presentation in \(\mathrm{EntC}\) is
stronger than finding some isomorphic presentation.  Write its given
witness as
\[
 E=\mathrm{kerEntC}(P,\mathcal O,B;H),
 \qquad H:\mathrm{KerEntHyp}(P,\mathcal O,B,X,\mathrm{tw}).
 \tag{WRS.5}
\]
The polynomial weight spaces, evaluation maps and their kernels do
not involve the three changed fields in (WRS.1).  Hence their quotients
\(V_{u,\mu,i}=\mathrm{polV}_{u,i}/\ker(\mathrm{evP}_{u,\mu,i})\)
are literally the same quotients.  Retain their sections, the
coefficient embeddings \(P.i_L\), all matrices \(B_{u,\mu,w,i,c,\tau}\),
and the inverse maps used on the original polynomial image.  This
retains the source and target functionals \(\mathrm{prS}\),
\(\mathrm{prT}\), the Taylor series \(\mathrm{Th}\), and the
normalizing scalar \(\mathrm{unC}\).  Its convention replacing a
zero raw scalar by \(1\) is also unchanged.

The four clauses of \(H\) now remain the same: the target
coefficient norm is the norm of a coefficient-field element; the
source lattice is characterized by all source minor functionals
having norm at most one; the target lattice has the corresponding
target characterization; and the kernel formula is
\[
 \mathrm{prT}(v',I)=\mathrm{unC}\,
 \mathrm{prS}\!\left(v,
       \mathrm{applyAt0}(\mathrm{Th},\mathrm{Gxz}(I))\right)
 \tag{WRS.6}
\]
whenever the given global kernel maps \(\mathrm{evC}(v)\) to
\(\mathrm{evC}(v')\).  The original level-free guard also stays
unchanged, since it reads \(X\) and the twisted column data, not
the replaced division lattice.  Thus the very same \(H\) supplies
\(H'\), and every field in the definition of \(\mathrm{kerEntC}\)
is the corresponding transported field in (WRS.5).
Applying this transport to that one equality proves
\(E'=\mathrm{kerEntC}(P',\mathcal O',B;H')\).
This proves the full original
\texttt{kernel\allowbreak Taylor\allowbreak Series\_\allowbreak preserved\_\allowbreak division\allowbreak Lattice\allowbreak Change},
with its \(\mathrm{EntC}\) conclusion when supplied originally.
\hfill\(\square\)
% END WRS2

% BEGIN WRS3
\paragraph{WRS3. The dependent comparison--kernel equality transports once.}
\label{claim:WRS3}
Suppose \(O_2:\mathrm{Obj2}(w.DZ)\), a column record \(t\), and
\(k_0\) satisfy the original equality
\[
 (O_2.G_t,O_2.c,O_2.\mathrm{tw},O_2.\mathrm{ker})
    =(t.G_t,t.c,w.\mathrm{tw}(Z),k_0).                \tag{WRS.7}
\]
This is an equality in the dependent sum whose fourth entry has
the type
\[
              \mathrm{KerFam}(D,G_t,c,\mathrm{tw},O_2.G_y).
\]
For every \(N,h\), the corresponding transported equality is
\[
 (O_2'.G_t,O_2'.c,O_2'.\mathrm{tw},O_2'.\mathrm{ker})
    =(t'.G_t,t'.c,w.\mathrm{tw}(Z'),k_0').            \tag{WRS.8}
\]
Here \(Z'=\mathrm{setDivZ}(Z,N,h)\), so
\(w.DZ'=\mathrm{setDiv}(w.DZ,N,h)\).

Define a map on the entire dependent sum in (WRS.7): apply
\(\mathrm{gatTr}\), \(\mathrm{cmpTr}\) and
\(\mathrm{twTr}\) to its first three entries and apply WRS1 to
its fourth entry, using the transported \(O_2.G_y\).  WRS1 shows
that the fourth entry has exactly the type required by the first
three transported entries.  Thus this is a map on the whole sum,
not four independent identifications of records of different types.
The image of the left side is the left side of (WRS.8), by the
definition of \(\mathrm{Obj2.tr}\).  GCF16 identifies the
transported column's first entries with the specified
\(\mathrm{gatTr},\mathrm{cmpTr}\), and its twisted forms with
\(w.\mathrm{tw}(Z')\).  The right image is therefore exactly the
right side of (WRS.8).  Equality is preserved by this one map,
which proves the assertion.  If \(\mathrm{KerT}\) is supplied,
WRS2 transports it on this same \(k_0'\).
\hfill\(\square\)
% END WRS3

% BEGIN WRS4
\paragraph{WRS4. A built core supplies all moved column ties on the given object.}
\label{claim:WRS4}
Let \(b:\mathrm{Built}(\mathcal S,K,w,Z,c_0)\) be the original
built-core witness, for arbitrary original slots \(\mathcal S\).
Let \(N,h\) be arbitrary division data and let \(O\) be any full
object over \(D'=\mathrm{setDiv}(w.DZ,N,h)\) such that
\[
             \mathrm{ofObj}(O)=c_0.O.\mathrm{tr}(N,h).
 \tag{WRS.9}
\]
There is the specified transported
\(t'=\mathrm{columnLocalDataTransportDivision}(b.\mathrm{col})\)
such that
\[
 \begin{gathered}
 \mathrm{ColsExtraTiedS}(t'.\mathrm{full},O.\mathrm{ex}),\\
 \mathcal G:\mathrm{GroupCompatibilityAll}
             (D',O.G_t,O.\mathrm{ex}.c),\\
 \mathcal G=t'.g,\qquad
 \mathcal G.\mathrm{herm}=t'.\mathrm{full}.\mathrm{to0},\qquad
 \mathcal G.\mathrm{forms}(u)=t'.C.\mathrm{forms}(u).\mathrm{to0}.
 \end{gathered}                                                   \tag{WRS.10}
\]
As usual the equality of records in (WRS.10) is read after the
joint dependent identification of their indexing data.  The given
kernel of \(O.\mathrm{ex}\) is part of this conclusion.

The \(\mathrm{Built}\) witness contains its actual
\(b.\mathrm{col}:\mathrm{ColumnDataSplitPrimes}\), its kernel
family \(k_0\), its \(\mathrm{KerT}\), and precisely (WRS.7).
It also contains
\[
 c_0.O.\mathrm{lev}=b.\mathrm{col.full.base.lev}_0,
 \qquad c_0.O.\mathrm{errX}=\mathrm{toErrXS}(b.\mathrm{col.full}).
 \tag{WRS.11}
\]
GCF16 provides \(t'\) with the original coefficient field, the
full five compatibility fields, all column readings, \(\mathrm{IsRecG}\)
and \(\mathrm{OmegaCompatible}\).  WRS2 provides the transported
kernel tie, and WRS3 gives the joint identity (WRS.8).
Substitute that identity once and then use (WRS.9).

The three fields of \(\mathrm{ColsExtraTiedS}\), expanded as
\(\mathrm{SymmetricExtraCompatibility}\), are now available on
the actual \(O.\mathrm{ex}\).  Its join field consists of the two
equalities (WRS.11) and the column record's exact prime-compatibility
law.  Its kernel-embedding field is the transported
\(b.\mathrm{ker.iw}\), on the given kernel, and its positive-place
field is \(b.\mathrm{ker.iw\_\allowbreak pos}\) on that same kernel.  Explicitly,
the former preserves at every allowed \(u,w,v\) the identity
\[
 e_{Q,v}\bigl((O.\mathrm{ex.ker})_{u,w}.i_w(0,v,z)\bigr)
   =\operatorname{pr}_1\bigl(\mathrm{spl}_v(\mathrm{incl}_E z)\bigr)
       \qquad(z\in E).                              \tag{WRS.12}
\]
Use \(t'.g\) itself for \(\mathcal G\).  Its Hermitian and form
equalities are already part of the transported column record.
These prove every assertion in (WRS.10).  Neither \(\mathrm{EntC}\),
a new row choice, a full \(\mathrm{Wit}\), nor a selector is a
premise of this statement.  The original built-core witness is its
actual input; it already supplies the column and kernel data used.
\hfill\(\square\)
% END WRS4

% BEGIN WRS5
\paragraph{WRS5. The original built point has that exact object.}
\label{claim:WRS5}
Let \(\beta:\mathrm{BPt}(K)\) satisfy the original
\(\beta.\mathrm{Is}(\mathrm{IsBuilt}(\mathcal S))\).
Then (WRS.10), including its full symmetric extra compatibility,
holds on the actual \((\beta.D,\beta.O)\).
The compatibility record is on \(\beta.O.\mathrm{ex}.c\), and
its kernel equations hold on \(\beta.O.\mathrm{ex.ker}\).

By the definition of \(\mathrm{IsBuilt}\), take the original
\(w,Z,c_0,b\) with the single joint equality of the columns,
later stage and core to \(\beta.X,\beta.c\).  Substitute this
dependent equality; it identifies the base with \(D=w.DZ\)
and the core with \(c_0\).  The remaining data of \(\beta\)
contain a layout \(Y\) and its specified equality
\[
                 \mathrm{layObj}(Y)=c_0.O,
 \tag{WRS.13}
\]
together with its given \(p:\mathrm{LedPart}(K.C,Y)\).

The source's final point is not an unspecified isomorphic model.
Its division lattice and height are exactly
\[
 \begin{aligned}
 N'_u&=\mathcal O_L\operatorname{span}
       \{a y:a\in\mathrm{Iu}(D,p.L,u),\ y\in p.L.\mathrm{NuLat}(u)\},\\
 h'_u&=\mathrm{hExpl}_6(D,p.L,Y.X,K.C.\mathrm{Csz})(u),\\
 p.\mathrm{pt}.D&=\mathrm{setDiv}(D,N',h').
 \end{aligned}                                                    \tag{WRS.14}
\]
Before this operation, the full object \(Y.\mathrm{obj}(p)\)
uses the given size and group data of \(p\), and the comparison,
twisted forms, local lattices and kernel of \(Y\).
Thus its \(\mathrm{Obj2}\) is precisely \(\mathrm{layObj}(Y)\).
The definition of \(p.\mathrm{pt}\) transports that full object
by (WRS.14).  Reading each field therefore gives
\[
 \begin{aligned}
 \mathrm{Obj.ofPt}(p.\mathrm{pt})
       &=\bigl(Y.\mathrm{obj}(p)\bigr).\mathrm{tr}(N',h'),\\
 \mathrm{ofObj}(\beta.O)
       &=\mathrm{layObj}(Y).\mathrm{tr}(N',h')
        =c_0.O.\mathrm{tr}(N',h').
 \end{aligned}                                                    \tag{WRS.15}
\]
For the local lattice field, the point uses \(Y.G_y\); the
original \(p.G_y\)-equality transports the group tie to that same
field.  Hence (WRS.15) identifies the entire kernel domain as well
as its comparison, rather than merely forgetting that distinction.

Apply WRS4 with exactly \(N',h',O=\beta.O\).
No size bound on (WRS.14) is needed for this application, and none
of \(p\)'s original sizes, metrics, prime sets or constants is
replaced.  The output is the full original compatibility on the
actual point.  In particular its coefficient field is the one
already in \(\beta.D\).
\hfill\(\square\)
% END WRS5

% BEGIN WRS6
\paragraph{WRS6. Every point of the original built class receives the same-data supplier.}
\label{claim:WRS6}
Let \(Q\) be any original core predicate refining
\(\mathrm{IsBuilt}(\mathcal S)\).  Keep arbitrary original bundle
type \(\mathrm{Bd}\), linkage predicate \(\mathrm{Lk}\), and
further property \(E\).  Define the original class
\[
                \mathcal A=\mathrm{Fc.cls}(Q,\mathrm{Lk},E).
 \tag{WRS.16}
\]
For every \(\mathrm{Fm}\in\mathcal A\), every
\(x>\mathrm{Fm}.x_1\), every witness \(h_x\) of that inequality,
and every \(i\in\mathrm{Fm}.I_x\), the actual object
\[
 O_{x,h_x,i}=
 (\mathrm{Fm}.L,\mathrm{Fm}.G_t,\mathrm{Fm}.G_x,
  \mathrm{Fm}.G_y,\mathrm{Fm}.G_K,\mathrm{Fm}.X,
  \mathrm{Fm}.c_X,\mathrm{Fm}.\mathrm{ex})_{x,h_x,i}
 \tag{WRS.17}
\]
has the full conclusion (WRS.10).  No additional \(\mathrm{TSel}\)
or eventual witness-realization hypothesis is needed.

Indeed class membership supplies one original parameter record
\(\kappa\), its \(\mathrm{Link}(\kappa,\mathrm{Fm})\), and
the seven-duty class whose object property is
\(\mathrm{QOf}(Q,\mathrm{Lk},E,\kappa)\).
The link supplies the required positivity of \(\kappa\).
Apply this object property at the specified \(x,h_x,i\), using
the family's actual \(\mathrm{Csz}\).  The definition supplies
a bundle \(b\) and an actual built witness
\(K,\beta,b'\), together with
\[
  \bigl(\mathrm{Fm}.D(x,i),O_{x,h_x,i},b\bigr)
       =\bigl(\beta.D,\beta.O,b'\bigr).              \tag{WRS.18}
\]
This is the original equality in the dependent sum of the base,
full object and bundle, not a collection of numerical agreements.
Also \(\beta.\mathrm{Is}(Q)\) holds.  Its refinement gives the
built-core witness required by WRS5.  Apply that claim and substitute
(WRS.18) once.  It identifies the coefficient rings, the comparison
map, the complete kernel family, and their indexing records with
the ones in (WRS.17).  This proves the assertion for the originally
specified point and hence for all of them.

The proof does not change \(E\) to insert compatibility and does
not change \(\mathrm{Lk}\).  Compatibility follows from the
\emph{existing} built witness in their class definition.  For the
original selected predicate \(Q=\mathrm{IsGoodE}(\mathcal S)\),
the required refinement is immediate from its definition: forget
\(\mathrm{GoodO}\) and \(\mathrm{EntC}\), retaining its actual
\(\mathrm{Built}\) witness.  The same applies to any further
conjunction used in the original common class.
\hfill\(\square\)
% END WRS6

% BEGIN WRS7
\paragraph{WRS7. Full family compatibility and its original Closed order.}
\label{claim:WRS7}
For the class \(\mathcal A\) of WRS6 and every subclass
\(\mathcal C\subseteq\mathcal A\),
\[
 \begin{gathered}
 \mathcal C(\mathrm{Fm})\Longrightarrow
                 \mathrm{GroupCompatibilityFamily}(\mathrm{Fm}),\\
 \mathrm{Closed}\bigl(\mathcal C,
           (\mathrm{Fm},C_H)\mapsto
                    \mathrm{GroupCompatibilityFamily}(\mathrm{Fm})\bigr).
 \end{gathered}                                                   \tag{WRS.19}
\]
In particular these hold on \(D_3(3,\mathcal A)\), on its
target-compatible subclass, and on the original \(W_3(\mathcal A)\).

The first predicate in (WRS.19) is exactly
\[
 \begin{gathered}
 \exists x_g\ \forall x\ge x_g\ \forall i\ \forall h_x:
          \mathrm{Fm}.x_1<x,\\
 \operatorname{Nonempty}\!\left(
 \mathrm{GroupCompatibilityAll}
       (\mathrm{Fm}.D(x,i),\mathrm{Fm}.G_t(x,i),
                       \mathrm{Fm}.\mathrm{ex}(x,h_x,i).c)\right).
 \end{gathered}
 \tag{WRS.20}
\]
Take \(x_g=\mathrm{Fm}.x_1+1\).  WRS6 applies to every later
\(x,i,h_x\), on the exact \(\mathrm{Extra}\) named in (WRS.20).
This uniform threshold precedes all places, words, local tuples
and coefficient completions.  It also precedes all choices of
\(h_x\): the source class property is universal in that argument,
and each application of (WRS.18) uses its specified value.

To check the second assertion, expand the source's \(\mathrm{Closed}\):
\[
 \begin{gathered}
 \exists\epsilon_1>0\ \forall\epsilon\in(0,\epsilon_1]\
 \exists K_0:\mathrm{TransportedAlt}\\
 \forall\mathrm{Fm}:\mathrm{FamP}(\epsilon)\ \forall C_H\in\mathbb R,
 \qquad \mathcal C(\mathrm{Fm})\ \land\
       \mathrm{Fm.Large}(K_0,C_H)\Longrightarrow P(\mathrm{Fm},C_H).
 \end{gathered}
 \tag{WRS.21}
\]
Choose \(\epsilon_1=1\) and choose the bound with all four
components zero.  This choice is made before \(\mathrm{Fm}\)
and \(C_H\).  The first assertion of (WRS.19) proves the
conclusion without using the large-coefficient guard; consequently
it certainly proves (WRS.21) with that guard.  The family threshold
\(x_g\) is chosen afterwards, in its permitted position.
This argument is valid on every fixed subclass, so it also applies
after adding the original smoothness, twisting, selected-top,
\(\mathrm{TuSize}\), \(\mathrm{DGamSize}\) and target predicates.
None of those operations changes the members' objects.
\hfill\(\square\)
% END WRS7

% BEGIN WRS8
\paragraph{WRS8. The second Closed input uses the original selected tops.}
\label{claim:WRS8}
For the very same \(\mathcal A\), the original selected class
\[
 D_3(3,\mathcal A)=\mathrm{Des}(\mathrm{WithChN}(\mathcal A),3)
 \tag{WRS.22}
\]
satisfies
\[
       \mathrm{Closed}(D_3(3,\mathcal A),\mathrm{Adm2F}).
 \tag{WRS.23}
\]
Its \(\mathrm{ChosenE}\) field supplies \(\mathrm{TopsF}\)
on the actual chosen subset.  In particular this is the same
family, selected subset and degree assignment as in WRS7.

Expand membership in (WRS.22).  It includes
\(\exists C_H'\,\mathrm{ChosenE}(\mathrm{Fm},3,C_H')\)
in addition to the original base-class, smoothness and twisting
conditions and \(\mathrm{TuSize},\mathrm{DGamSize}\).
GCF18 projects its \(\mathrm{ChosenTT}\) and one positive
\(C_b\) with \(\mathrm{ChosenBad}\).  The completed estimate
GLG16 then gives \(\mathrm{Adm2F}\) with the original threshold
\[
 x_{\mathrm{top}}=
 \max\{x_a,x_b,\mathrm{Fm}.x_1+1,
                         (16C_b+2)^{3/\epsilon}\}.
 \tag{WRS.24}
\]
This is a theorem about the supplied selected tops, not a new
choice of a random subset or a later restriction of the family.
It holds for every member of the class before imposing any
coefficient bounds.  The proof of (WRS.21), again using cutoff
\(1\) and zero coefficient bounds, therefore proves (WRS.23).
The internal witness \(C_H'\) in \(\mathrm{ChosenE}\) need
not equal the external \(C_H\) universally quantified in
\(\mathrm{Closed}\); neither this projection nor its estimate
asserts such an equality.  All thresholds retain their permitted
dependence on the fixed family and its chosen-top witnesses.
\hfill\(\square\)
% END WRS8

% BEGIN WRS9
\paragraph{WRS9. The full original LinkG on the selected built class.}
\label{claim:WRS9}
For every original built class \(\mathcal A\) in WRS6,
\[
             \mathrm{Closed}(D_3(3,\mathcal A),\mathrm{LinkG}).
 \tag{WRS.25}
\]
This has no additional \(\mathrm{TSel}\), \(\mathrm{WPt}\),
kernel-existence or compatibility premise.  The actual compatibility
input is supplied by WRS7 and the actual \(\mathrm{Adm2F}\)
input by WRS8.

Apply the complete class-level theorem of LGS13 to these two
\(\mathrm{Closed}\) inputs.  Its coefficient choice can be
written explicitly here: choose exponent cutoff \(1/104\);
for a later \(\epsilon\in(0,1/104]\), take the pointwise maximum
of the zero bounds in WRS7--WRS8 and
\[
 \begin{gathered}
 K_{\mathrm{loc}}.\mathrm{pre}=0,\qquad
 K_{\mathrm{loc}}.t=0,\qquad
 K_{\mathrm{loc}}.\mathrm{mid}=4129,\\
 K_{\mathrm{loc}}.\mathrm{last}
     =\max\{5600002,C_{\mathrm{orb}}(\epsilon)\},\\
 C_{\mathrm{orb}}(\epsilon)
   =40(6r+1)\bigl((h+3)10^{6r+1}+6\bigr),\qquad
 h=2+8\epsilon,\quad r=\lceil8(h+2)/\epsilon\rceil_++1.
 \end{gathered}                                                   \tag{WRS.26}
\]
The bounds are chosen before \(\mathrm{Fm},C_H\).
The given family's \(\mathrm{Large}\) hypothesis at this
bound is exactly the sufficient coefficient region used in LGS13.
Now fix a member and all its constants.  LGS13 and GLG9 choose
one \(C_{\mathrm{count}}\ge0\) and the finitely many eventual
thresholds.  Increase their maximum to include
\(x_g\), (WRS.24), and \(\mathrm{Fm}.x_1+1\).  All later
members \(i\), places \(w\), primes, served representatives,
local degrees and coefficient embeddings use this same threshold.
Set
\[
 C_G=4C_{\mathrm{count}}+1,\qquad
 \delta_G(x)=x^{-\epsilon/2},\qquad
 b_o=(\mathrm{Fm}.L(x,i).\mathrm{sc}.b)^{\delta_G(x)}.
 \tag{WRS.27}
\]
Since \(\epsilon>0\), \(\delta_G(x)\to0\) as \(x\to+\infty\).
The function \(\delta_G\) is chosen once, before any of the
later indices or places.

For precision, the full pointwise conclusion is the literal
\[
 \mathrm{LinkG\_\allowbreak at}(\mathrm{Fm}.G_x(x,i),\mathrm{Fm}.G_K,C_G,b_o).
\]
Write \(D,L,X,K\) for that member's original records and
\(p=p_w\).  Its five conjuncts are the following.
First, for every \(u,w\) with \(p\ne T_u\),
\[
 \begin{gathered}
 \mathrm{LocLe}_{D.\mathrm{plEq}(w)}
    (\lfloor\mathrm{ur}_w(L.\tau(w))\rfloor_+,
          L.\mathrm{NuLat}(u),\mathrm{MramL}(D,u)),\\
 L.\mathrm{lenMN}(u,w)
       \le D.\mathrm{duS}(u)\{L.\tau(w)+L.\mathrm{io}(w)\},\\
 \mathrm{VV}_{p}=\varnothing\Longrightarrow
       L.\tau(w)=L.\mathrm{io}(w)=L.\mathrm{lenMN}(u,w)=0.
 \end{gathered}                                                   \tag{WRS.28}
\]
Second, for every original prime argument and every
\(t_0\in\mathrm{VV}_{p}\),
\[
 S_v(X,K,a_L,b_L,p,t_0)
      \le 2m\,N_v(X,p,t_0)\,L.\mathrm{sc}.a/p.
 \tag{WRS.29}
\]
Third, at every place with the original \(\mathrm{InertK}(D,p)\)
guard and \(p\le L.\mathrm{sc}.X''(D.R.n)\),
\[
 \begin{gathered}
 0\le L.\mathrm{gain}(j,w)
       \le4|TT|L.\mathrm{sc}.a/p
       \qquad(j\in\mathrm{lifts}(D)),\\
 \sum_{j\in\mathrm{lifts}(D)}\ \sum_{t\in L.TT_w}
   \bigl\{L.\mathrm{vonMangoldtClassSum}(j,t,p)
                  -\mathrm{loG}(L,C_G,b_o,p)\bigr\}
       -m\,\mathrm{flG}(L,C_G,p)\\
 \hspace{25mm}\le
       \sum_{j\in\mathrm{lifts}(D)}L.\mathrm{gain}(j,w)-L.\tau(w).
 \end{gathered}                                                   \tag{WRS.30}
\]
Fourth, \(\mathrm{ServedOrbitsContained}(X,K,a_L,b_L)\) holds
for the original orbit representatives.  Fifth, the full original
\(\mathrm{MeanKappaBound}(X,K)\), with both its clauses, holds.

These are exactly the conclusions of LGS13 on the supplied
compatibility and selected tops.  In particular its local length
identification uses the original valuation normalization; it
introduces no extra ramification factor into (WRS.28).  Its
auxiliary zero-chain application leaves the actual \(D,L,X,K\),
comparison, local lattices and all five displayed conclusions
unchanged.  WRS1--WRS7 identified the geometric kernel before
that auxiliary argument; they did not replace it by that zero chain.
Thus the counting, local-length, coverage and mean estimates all
concern the single original member.  Equations (WRS.27)--(WRS.30)
and the last two predicates prove the original \(\mathrm{LinkG}\),
and the preceding coefficient choice proves its full \(\mathrm{Closed}\)
form (WRS.25).
\hfill\(\square\)
% END WRS9

% BEGIN WRS10
\paragraph{WRS10. Literal universal installer and the selected closure boundary.}
\label{claim:WRS10}
The original statement
\[
 \begin{gathered}
 \forall\mathcal A:\mathrm{FamClass},\quad
 \mathrm{TSel}(\mathcal A)\ \Longrightarrow\
 \bigl(\forall\mathrm{Fm}\in\mathcal A,
                         \mathrm{Duty}(\mathrm{Fm})\bigr)
       \ \Longrightarrow\\
 \mathrm{Closed}(D_3(3,\mathcal A),\mathrm{LinkG})
 \end{gathered}                                                   \tag{WRS.31}
\]
is proved with every original quantifier.  The selected built-class
conclusion (WRS.25) is proved separately without either of those
two extra premises.  Both versions retain the same family objects,
the same \(D_3\) conditions and the exact five-field
\(\mathrm{LinkG\_\allowbreak at}\) conclusion.

To prove (WRS.31), fix the arbitrary original class and selector.
At each family member the selector gives one threshold and, after
that threshold, its original \(\mathrm{RecT}\) on every supplied
\(\mathrm{Extra}\).  Its \(T.g\) field is the required
\(\mathrm{GroupCompatibilityAll}\) on that \(\mathrm{Extra}.c\).
GCF17 therefore gives \(\mathrm{GroupCompatibilityFamily}\)
on the same member, with the maximum of the selector threshold
and \(\mathrm{Fm}.x_1+1\).  The elementary \(\mathrm{Closed}\)
argument of WRS7 supplies its first class-level input.
GCF18 supplies the original selected-top input on \(D_3\),
and LGS13 proves the conclusion.  This proof even leaves the
original duty premise unused; retaining it gives exactly
\texttt{Global\allowbreak Theory.\allowbreak Top.\allowbreak link\_\allowbreak g}.  The framed version, with
its subsequent repeated \(D_3\)-membership,
\(\mathrm{TuSize}\) and \(\mathrm{DGamSize}\) guards, follows
by ignoring those redundant guards after (WRS.31).

For the actual class selected by the construction, the route is
instead WRS5--WRS9.  Its witnesses are extracted from existing
built-class membership using (WRS.18); no \(\mathrm{TSel}\)
premise is inserted.  If the common construction intersects
\(\mathcal A\) with its target class before adding the original
\(\mathrm{WithChN},\mathrm{WithSm},\mathrm{WithTw}\) and
selected-design predicates, membership still implies membership
in \(\mathcal A\).  WRS7 applies directly to those same objects;
the selected-top projection and the coefficient choice in WRS9
are unchanged.  Equivalently one can restrict the \(\mathrm{Closed}\)
theorem to this fixed subclass, without choosing a new member.

Consequently the original local-gain supplier is complete on the
selected built class.  It does not assert
\(\mathrm{DesignExists}(\mathcal A)\): all statements above are
universal over members, and the original common-design obligation
still constructs a member simultaneously satisfying its other
geometric, row, period and arithmetic requirements.  Nor does
local compatibility manufacture the stronger \(\mathrm{RecT}\)
record needed by the denominator, period, spectral and other
second-layer arguments.  Those distinct original producers remain
separate.  This proves precisely the selected local-gain endpoint,
and makes no assertion of a completed RH or GRH proof.
\hfill\(\square\)
% END WRS10

\endgroup
\end{document}
